% ---------------------------------------------------------------------------------------------
% 30 days of simplydata: the printable guide. This file is made by docs/make_30_days.py from
% simplydata/learn_days.py and the page numbers of the compiled manual. Do not edit it by hand:
% change simplydata/learn_days.py, then run
%     python docs/make_manual.py --no-run
% ---------------------------------------------------------------------------------------------
% Page numbers from simplydata_manual.tex 72ef04b81965b111

\documentclass[10pt,a4paper]{article}
% no dates or random IDs in the PDF, so it only changes when the guide changes
\pdfinfoomitdate=1
\pdftrailerid{}
\pdfsuppressptexinfo=-1
\usepackage[T1]{fontenc}
\usepackage{XCharter}
\usepackage[xcharter,vvarbb]{newtxmath}
\usepackage[lining,scale=0.88]{FiraSans}
\usepackage[varqu,varl,scaled=0.92]{zi4}
\usepackage{microtype}
\usepackage[a4paper,top=2.0cm,bottom=1.9cm,left=1.9cm,right=1.9cm,headheight=14pt,headsep=0.5cm]{geometry}
\usepackage[dvipsnames]{xcolor}
\usepackage{array}
\usepackage{longtable}
\usepackage{booktabs}
\usepackage{enumitem}
\usepackage{needspace}
\usepackage{fancyhdr}
\usepackage{listings}
\usepackage[most]{tcolorbox}
\usepackage[hidelinks]{hyperref}
\usepackage{bookmark}

\setlength{\parindent}{0pt}
\setlength{\parskip}{0.45em plus 0.2em}
\setlength{\emergencystretch}{3em}
\tolerance=3000
\raggedbottom

% the colours of the manual and the dictionary
\definecolor{brand}{HTML}{1F5F99}
\definecolor{brandlight}{HTML}{E8F0F8}
\definecolor{codebg}{HTML}{F5F7FA}
\definecolor{coderule}{HTML}{D5DCE4}
\definecolor{codefg}{HTML}{24292F}
\definecolor{codekw}{HTML}{1F5F99}
\definecolor{codestr}{HTML}{A33B00}
\definecolor{codecom}{HTML}{6A737D}
\definecolor{goodbg}{HTML}{EEF7EE}
\definecolor{goodframe}{HTML}{2E7D32}
\definecolor{trybg}{HTML}{FFF4E5}
\definecolor{tryframe}{HTML}{C25E00}
\definecolor{inkgrey}{HTML}{5A6570}
\hypersetup{colorlinks=true,linkcolor=brand,urlcolor=brand,filecolor=brand,
  pdftitle={30 days of simplydata},pdfauthor={The simplydata project},
  pdfsubject={An hour a day for 30 days: warm-up questions, lessons and tasks}}

\newcommand{\code}[1]{\texttt{\color{codefg}#1}}
\lstset{
  basicstyle=\ttfamily\small\color{codefg},
  columns=fullflexible, keepspaces=true, upquote=true, showstringspaces=false,
  breaklines=true, breakatwhitespace=false,
  postbreak=\mbox{\textcolor{codecom}{$\hookrightarrow$}\space},
}
\lstdefinestyle{py}{
  language=Python,
  keywordstyle=\color{codekw}\bfseries, stringstyle=\color{codestr}, commentstyle=\color{codecom},
  emph={sdata}, emphstyle=\color{codekw},
  backgroundcolor=\color{codebg}, frame=single, rulecolor=\color{coderule}, framerule=0.4pt,
  framesep=4pt, xleftmargin=2.0em, xrightmargin=4pt,
  numbers=left, numberstyle=\tiny\sffamily\color{codecom}, numbersep=10pt,
  aboveskip=0.5em, belowskip=0.4em,
}
\newenvironment{titleblock}{\par\noindent\begin{minipage}{\linewidth}}{\end{minipage}\par\vspace{0.6em}
  {\color{brand}\rule{\linewidth}{1.2pt}}\par}
\makeatletter
\renewcommand{\section}{\@startsection{section}{1}{0pt}{1.1em}{0.4em}{\sffamily\large\bfseries\color{brand}}}
\renewcommand{\subsection}{\@startsection{subsection}{2}{0pt}{0.9em}{0.3em}{\sffamily\bfseries\color{brand}}}
\makeatother
\newcommand{\citeref}[1]{{\small(#1)}}
\setlist[itemize]{itemsep=0.15em,parsep=0pt,topsep=0.2em,leftmargin=1.3em}
\setlist[enumerate]{itemsep=0.15em,parsep=0pt,topsep=0.2em,leftmargin=1.6em}

% ---------------------------------------------------------------- word kinds, as in the dictionary
\newcommand{\kind}[1]{{\sffamily\scriptsize\bfseries\color{white}\colorbox{brand}{\strut #1}}}
\newcommand{\manualpage}[3]{\href{simplydata_manual.pdf\##3}{p.~#2}}

% ---------------------------------------------------------------- a day
\newcommand{\tick}{\raisebox{-0.1em}{\Large$\square$}}
\newcommand{\daystart}[4]{\clearpage\phantomsection\label{#4}\markboth{#3}{Day #1}%
  \addcontentsline{toc}{subsection}{Day #1: #2}%
  \noindent\begin{minipage}[b]{0.13\linewidth}{\sffamily\fontsize{30}{30}\selectfont\bfseries\color{brand}#1}\end{minipage}%
  \begin{minipage}[b]{0.87\linewidth}{\sffamily\small\color{inkgrey}Day #1 of 30 \textbullet{} #3}\par
  {\sffamily\Large\bfseries\color{brand}#2}\end{minipage}\par\vspace{0.2em}
  {\color{brand}\rule{\linewidth}{0.8pt}}\par\vspace{0.1em}
  {\sffamily\small\tick\ Warm-up\qquad\tick\ What to do\qquad\tick\ Task\qquad Date:\ \rule{3cm}{0.4pt}}\par}
\newcommand{\goal}[1]{\par\noindent{\sffamily\bfseries\color{brand}Goal.}\enspace#1\par}
\tcbset{boxbase/.style={enhanced,breakable,boxrule=0pt,leftrule=3pt,arc=1pt,outer arc=1pt,
  left=7pt,right=7pt,top=5pt,bottom=5pt,before skip=0.6em,after skip=0.6em,
  fonttitle=\sffamily\bfseries,coltitle=black,attach title to upper,title after break={}}}
\newtcolorbox{warmup}{boxbase,colback=brandlight,colframe=brand,
  title={\color{brand}1\enspace Warm-up, 5 minutes: from memory, without looking back\par}}
\newtcolorbox{todo}{boxbase,colback=white,colframe=brand,
  title={\color{brand}2\enspace What to do\par}}
\newtcolorbox{tryit}{boxbase,colback=trybg,colframe=tryframe,
  title={\color{tryframe}3\enspace Try it yourself, without looking back\par}}
\newtcolorbox{answerbox}{boxbase,colback=goodbg,colframe=goodframe,before skip=0.3em,after skip=0.5em}
\newcommand{\fromday}[1]{{\footnotesize\sffamily\color{inkgrey}(day~#1)}}
\newcommand{\writeline}{\par\vspace{0.9em}\noindent{\color{coderule}\rule{\linewidth}{0.4pt}}\par}
\newcommand{\field}[2]{\par\noindent\hangindent=1em{\sffamily\small\bfseries\color{inkgrey}#1}\enspace{\small#2}\par}
\newcommand{\whateachline}{\par\noindent{\sffamily\small\bfseries\color{brand}What each line does}\par}
\newlist{lineexplain}{description}{1}
\setlist[lineexplain]{leftmargin=2.0em,labelwidth=1.4em,labelsep=0.6em,align=right,
  font=\normalfont\sffamily\footnotesize\color{codecom},itemsep=0.1em,parsep=0pt,topsep=0.2em}

% ---------------------------------------------------------------- calendar and answers
\newcommand{\trackstart}[3]{\clearpage\phantomsection\label{#3}\markboth{#1}{#1}%
  \addcontentsline{toc}{section}{#1}%
  {\sffamily\fontsize{24}{28}\selectfont\bfseries\color{brand}#1\par}\vspace{0.2em}
  {\sffamily\large\color{inkgrey}#2\par}\vspace{0.2em}{\color{brand}\rule{\linewidth}{1.2pt}}\par}
\newcommand{\answerday}[3]{\par\needspace{6\baselineskip}\vspace{0.6em}\phantomsection\label{#3}%
  \noindent{\sffamily\bfseries\color{brand}Day #1: #2}\par\nopagebreak}

\pagestyle{fancy}
\fancyhf{}
\fancyhead[L]{\sffamily\small\color{inkgrey}30 days of simplydata}
\fancyhead[R]{\sffamily\small\bfseries\color{brand}\leftmark}
\fancyfoot[C]{\sffamily\small\thepage}
\renewcommand{\headrulewidth}{0.4pt}
\fancypagestyle{plain}{\fancyhf{}\fancyfoot[C]{\sffamily\small\thepage}\renewcommand{\headrulewidth}{0pt}}

\begin{document}

\thispagestyle{plain}

\begin{titleblock}
{\fontsize{30}{34}\selectfont\sffamily\bfseries\color{brand} 30 days of simplydata\par}
\vspace{0.4em}
{\large\sffamily An hour a day: a warm-up from memory, a lesson, and a task to do yourself\par}
\vspace{0.3em}
{\sffamily\small\color{inkgrey} Version 0.1.0 \textbullet{} two tracks of 30 days \textbullet{} made from the code by
\code{docs/make\_30\_days.py}\par}
\end{titleblock}

\section*{How to use this guide}
The same 30 days are in Python. Start once, then type one word each day:
\begin{lstlisting}[style=py]
import simplydata as sdata
sdata.start_30_days()                    # the Python track: you already write some Python
sdata.start_30_days(track="beginner")    # the beginner track: you have never programmed
sdata.today()                            # every day: the day to do next
sdata.answers()                          # the answers to the warm-up and the task
\end{lstlisting}
This guide is the paper version, to keep next to the manual (\href{simplydata_manual.pdf}{\code{simplydata\_\allowbreak{}manual.\allowbreak{}pdf}}): every
page number is a link that opens the manual at that page. Each day has three parts and takes about
an hour:
\begin{enumerate}[leftmargin=1.6em]
  \item \textbf{Warm-up, 5 minutes.} A few questions about words from earlier days. Answer them
    \emph{from memory}, on the lines, before you look anything up. Then check the answers at the
    back of the guide (each day has a link to them) or with \code{sdata.answers(n)}.
  \item \textbf{What to do, about 45 minutes.} The lesson to work through (\code{sdata.lesson(n)}),
    the pages of the manual to read and the words to look up with \code{sdata.help("word")}. Type every
    example yourself: typing it, and getting it wrong, teaches more than copying.
  \item \textbf{Try it yourself, 10 minutes.} A task to do without looking back. If yours runs and
    answers the question, it is right; the worked answers are at the back.
\end{enumerate}
Tick the boxes at the top of the page as you go, and the day on the calendar of your track.
\textbf{Missed a day, or a week?} That is fine. Carry on where you stopped: the days build on each
other, and nothing is lost by taking a break. In Python, \code{sdata.today()} does the same.

\section*{Why a warm-up from memory?}
Trying to remember something fixes it in memory better than reading it again: in experiments,
people who were tested on a text remembered more of it days later than people who read it again
\citeref{Roediger and Karpicke 2006}. Practice spread over several days also works better than the
same practice in one go \citeref{Cepeda et al.\ 2006}. So each warm-up asks one question about
yesterday, one about three days ago and one about a week ago (and, later on, two weeks ago). Review
days (7, 14, 21 and 28) ask about the whole course, and the questions change each time a day comes back.

\section*{Which track?}
\begin{center}\small
\begin{tabular}{@{}>{\raggedright\arraybackslash}p{0.17\linewidth}>{\raggedright\arraybackslash}p{0.75\linewidth}@{}}
\hyperref[track:python]{\textbf{Python track}} & For people who already write some Python. About two lessons a day for 20 days, with review days on days 7, 14, 21 and 28, then an eight-day project of your own (days 21 to 28), a final test from memory and a look at the same work in pandas and matplotlib. Start: \code{sdata.\allowbreak{}start\_\allowbreak{}30\_\allowbreak{}days()}.\\[0.4em]
\hyperref[track:beginner]{\textbf{Beginner track}} & For people who have never programmed. Three days of just enough Python first, then about one lesson a day, with review days on days 7, 14 and 21, a small project of your own (days 27 and 28), a final test from memory and a look back. Start: \code{sdata.\allowbreak{}start\_\allowbreak{}30\_\allowbreak{}days(track=\textquotedbl{}beginner\textquotedbl{})}.\\[0.4em]
\end{tabular}
\end{center}

\section*{The words in Python}
\begin{center}\small
\begin{tabular}{@{}>{\raggedright\arraybackslash}p{0.36\linewidth}>{\raggedright\arraybackslash}p{0.58\linewidth}@{}}
\kind{sdata}\,\code{sdata.\allowbreak{}start\_\allowbreak{}30\_\allowbreak{}days()}~\manualpage{word:sdata.start-30-days}{239}{section*.336} & Start 30 days of simplydata: an hour a day of warm-up questions, lessons and tasks.\\[0.3em]
\kind{sdata}\,\code{sdata.\allowbreak{}today()}~\manualpage{word:sdata.today}{240}{section*.337} & Show today's day of the 30 days of simplydata: warm-up questions, what to do and a task.\\[0.3em]
\kind{sdata}\,\code{sdata.\allowbreak{}day()}~\manualpage{word:sdata.day}{240}{section*.338} & Show any day of the 30 days, without changing where you are.\\[0.3em]
\kind{sdata}\,\code{sdata.\allowbreak{}answers()}~\manualpage{word:sdata.answers}{240}{section*.339} & The answers to a day's warm-up questions, and a worked answer to its task.\\[0.3em]
\kind{sdata}\,\code{sdata.\allowbreak{}next\_\allowbreak{}day()}~\manualpage{word:sdata.next-day}{241}{section*.340} & Go on to the next day now, without waiting for tomorrow.\\[0.3em]
\end{tabular}\end{center}
Your progress is kept only on your computer, in a small file (\code{\textasciitilde{}/.\allowbreak{}simplydata/progress.\allowbreak{}json}). Nothing is sent anywhere.

\section*{In this guide}
\begin{center}\small
\begin{tabular}{@{}l r@{}}
\hyperref[track:python]{Python track: the calendar and the 30 days} & page~\pageref*{track:python}\\
\hyperref[track:beginner]{Beginner track: the calendar and the 30 days} & page~\pageref*{track:beginner}\\
\hyperref[answers:python]{Answers: Python track} & page~\pageref*{answers:python}\\
\hyperref[answers:beginner]{Answers: Beginner track} & page~\pageref*{answers:beginner}\\
\end{tabular}\end{center}

\section*{References}
{\small
Roediger, H.~L., \& Karpicke, J.~D. (2006). Test-enhanced learning: taking memory tests improves long-term
retention. \emph{Psychological Science}, 17(3), 249--255. \url{https://doi.org/10.1111/j.1467-9280.2006.01693.x}

Cepeda, N.~J., Pashler, H., Vul, E., Wixted, J.~T., \& Rohrer, D. (2006). Distributed practice in verbal recall
tasks: a review and quantitative synthesis. \emph{Psychological Bulletin}, 132(3), 354--380.
\url{https://doi.org/10.1037/0033-2909.132.3.354}\par}

\trackstart{Python track}{For people who already write some Python}{track:python}
About two lessons a day for 20 days, with review days on days 7, 14, 21 and 28, then an eight-day project of your own (days 21 to 28), a final test from memory and a look at the same work in pandas and matplotlib. Tick each day when you have done it; the day's number is a link to its page.
\begin{center}\small
\begin{longtable}{@{}c r >{\raggedright\arraybackslash}p{0.56\linewidth} >{\raggedright\arraybackslash}p{0.14\linewidth} r@{}}
\toprule & \textbf{Day} & \textbf{What you do} & \textbf{Lesson} & \textbf{Page}\\\midrule\endhead
\tick & \hyperref[day:python:1]{1} & Your first chart and your first table & 1, 2 & \pageref*{day:python:1}\\[0.25em]
\tick & \hyperref[day:python:2]{2} & Colours, word families and saving & 3 & \pageref*{day:python:2}\\[0.25em]
\tick & \hyperref[day:python:3]{3} & Every kind of chart & 4 & \pageref*{day:python:3}\\[0.25em]
\tick & \hyperref[day:python:4]{4} & Statistics on charts, and axes & 5, 6 & \pageref*{day:python:4}\\[0.25em]
\tick & \hyperref[day:python:5]{5} & Several charts, and a figure for a journal & 7, 8 & \pageref*{day:python:5}\\[0.25em]
\tick & \hyperref[day:python:6]{6} & Open a table, choose rows and columns & 9, 10 & \pageref*{day:python:6}\\[0.25em]
\tick & \hyperref[day:python:7]{7} & Review of the first week & review & \pageref*{day:python:7}\\[0.25em]
\tick & \hyperref[day:python:8]{8} & Clean the data and add columns & 11 & \pageref*{day:python:8}\\[0.25em]
\tick & \hyperref[day:python:9]{9} & Summaries, numbers and grids & 12, 13 & \pageref*{day:python:9}\\[0.25em]
\tick & \hyperref[day:python:10]{10} & Is a difference real? Tests and effect sizes & 14 & \pageref*{day:python:10}\\[0.25em]
\tick & \hyperref[day:python:11]{11} & Counts, outliers and many tests & 14 & \pageref*{day:python:11}\\[0.25em]
\tick & \hyperref[day:python:12]{12} & Choosing a test, and writing it up & - & \pageref*{day:python:12}\\[0.25em]
\tick & \hyperref[day:python:13]{13} & A whole analysis & 15 & \pageref*{day:python:13}\\[0.25em]
\tick & \hyperref[day:python:14]{14} & Review of the second week & review & \pageref*{day:python:14}\\[0.25em]
\tick & \hyperref[day:python:15]{15} & Models: explain one column from several & 16 & \pageref*{day:python:15}\\[0.25em]
\tick & \hyperref[day:python:16]{16} & Patterns across many columns & 17 & \pageref*{day:python:16}\\[0.25em]
\tick & \hyperref[day:python:17]{17} & Survival, and planning a study & 18 & \pageref*{day:python:17}\\[0.25em]
\tick & \hyperref[day:python:18]{18} & Fitting lines, curves and formulas & 19 & \pageref*{day:python:18}\\[0.25em]
\tick & \hyperref[day:python:19]{19} & Finding, solving and distributions & 19, 20 & \pageref*{day:python:19}\\[0.25em]
\tick & \hyperref[day:python:20]{20} & Values between points & 20 & \pageref*{day:python:20}\\[0.25em]
\tick & \hyperref[day:python:21]{21} & Review, and the question of your project & review & \pageref*{day:python:21}\\[0.25em]
\tick & \hyperref[day:python:22]{22} & Project: clean & project & \pageref*{day:python:22}\\[0.25em]
\tick & \hyperref[day:python:23]{23} & Project: explore & project & \pageref*{day:python:23}\\[0.25em]
\tick & \hyperref[day:python:24]{24} & Project: figures & project & \pageref*{day:python:24}\\[0.25em]
\tick & \hyperref[day:python:25]{25} & Project: tests & project & \pageref*{day:python:25}\\[0.25em]
\tick & \hyperref[day:python:26]{26} & Project: models and fits & project & \pageref*{day:python:26}\\[0.25em]
\tick & \hyperref[day:python:27]{27} & Project: results text & project & \pageref*{day:python:27}\\[0.25em]
\tick & \hyperref[day:python:28]{28} & Review, and finish your project & review & \pageref*{day:python:28}\\[0.25em]
\tick & \hyperref[day:python:29]{29} & Final test from memory & final test & \pageref*{day:python:29}\\[0.25em]
\tick & \hyperref[day:python:30]{30} & The same work in pandas and matplotlib & 15 & \pageref*{day:python:30}\\[0.25em]
\bottomrule\end{longtable}\end{center}

\daystart{1}{Your first chart and your first table}{Python track}{day:python:1}
\goal{Make a bar chart with titles and labels, and a first table you keep rows of and sum up. About an hour.}
\begin{warmup}
Nothing to remember yet on day 1: just start.
\end{warmup}
\begin{todo}
\begin{itemize}
  \item Work through \code{sdata.\allowbreak{}lesson(1)}: type each example yourself (do not copy and paste), run it and read what each line does. Then do the lesson's tasks and check them with \code{sdata.\allowbreak{}lesson(1,\allowbreak{} solutions=True)}.
  \item Work through \code{sdata.\allowbreak{}lesson(2)}: type each example yourself (do not copy and paste), run it and read what each line does. Then do the lesson's tasks and check them with \code{sdata.\allowbreak{}lesson(2,\allowbreak{} solutions=True)}.
  \item Run the example below: your first table. A table and a chart are separate things: one clean table can feed many charts. Tables get their own lessons on day 6.
\end{itemize}
\field{In the manual}{lesson~1, \manualpage{lesson:1}{10}{chapter.II.1}; lesson~2, \manualpage{lesson:2}{13}{chapter.II.2}; ``Welcome'', \manualpage{start:welcome}{2}{section*.3}.}
\field{Words to look up}{\kind{sdata}\,\code{sdata.\allowbreak{}example()}~\manualpage{word:sdata.example}{111}{section*.16}; \kind{sdata}\,\code{sdata.\allowbreak{}chart()}~\manualpage{word:sdata.chart}{111}{section*.15}; \kind{chart}\,\code{chart.\allowbreak{}bars()}~\manualpage{word:bars}{113}{section*.20}; \kind{chart}\,\code{chart.\allowbreak{}save()}~\manualpage{word:save}{112}{section*.18}; \kind{chart}\,\code{chart.\allowbreak{}title\_\allowbreak{}top()}~\manualpage{word:title-top}{131}{section*.59}; \kind{chart}\,\code{chart.\allowbreak{}label\_\allowbreak{}x()}~\manualpage{word:label-x}{132}{section*.61}; \kind{chart}\,\code{chart.\allowbreak{}label\_\allowbreak{}y()}~\manualpage{word:label-y}{132}{section*.62}; \kind{chart}\,\code{chart.\allowbreak{}rotate\_\allowbreak{}labels\_\allowbreak{}x()}~\manualpage{word:rotate-labels-x}{138}{section*.75}.}
\par\noindent{\sffamily\small\bfseries\color{brand}Example: type it yourself}\par
\begin{lstlisting}[style=py]
import simplydata as sdata
table = sdata.table(sdata.example("marks"))
table.keep_rows("marks", above=50)
table.average_of("marks", for_each="class")
chart = table.chart()
chart.bars(x="class", y="marks")
chart.save("good_marks.png")
\end{lstlisting}
\whateachline
\begin{lineexplain}
  \item[1] Load simplydata.
  \item[2] Start a table from the marks practice table.
  \item[3] Keep only the rows with marks above 50. It says how many rows are left.
  \item[4] The average mark of each class, in a sentence.
  \item[5] Start a chart from the table as it is now.
  \item[6] One bar for each class.
  \item[7] Save it.
\end{lineexplain}
\end{todo}
\begin{tryit}
\begin{enumerate}
  \item Make a bar chart of the \code{\textquotedbl{}weather\textquotedbl{}} table with \code{x=\textquotedbl{}city\textquotedbl{}} and \code{y=\textquotedbl{}rain\textquotedbl{}}, and save it as \code{\textquotedbl{}rain.\allowbreak{}png\textquotedbl{}}.
  \item Draw dots of \code{\textquotedbl{}hours\textquotedbl{}} (x) against \code{\textquotedbl{}marks\textquotedbl{}} (y) from the marks table, with a title, a name for each axis and a caption at the bottom.
  \item Start a table of the weather, keep the rows with \code{\textquotedbl{}temperature\textquotedbl{}} above 15, and give the average \code{\textquotedbl{}rain\textquotedbl{}} for each city.
\end{enumerate}
{\small Worked answers: \hyperref[answers:python:1]{page~\pageref*{answers:python:1}}, or \code{sdata.\allowbreak{}answers(1)}.}
\end{tryit}

\daystart{2}{Colours, word families and saving}{Python track}{day:python:2}
\goal{Give groups their own colours, learn how simplydata words are built, and save in every file type. About an hour.}
\begin{warmup}
\begin{enumerate}
  \item Which word lists the practice tables? \fromday{1}\writeline
  \item How do you start a chart of a table called marks? \fromday{1}\writeline
  \item Which word draws one bar for each name, as tall as the marks? \fromday{1}\writeline
\end{enumerate}
{\small Answers: \hyperref[answers:python:2]{page~\pageref*{answers:python:2}}, or \code{sdata.\allowbreak{}answers(2)}.}
\end{warmup}
\begin{todo}
\begin{itemize}
  \item Work through \code{sdata.\allowbreak{}lesson(3)}: type each example yourself (do not copy and paste), run it and read what each line does. Then do the lesson's tasks and check them with \code{sdata.\allowbreak{}lesson(3,\allowbreak{} solutions=True)}.
  \item Read \code{\textquotedbl{}How the words are built\textquotedbl{}} in the manual (two pages). It is the key to guessing words you have never seen.
  \item Type \code{sdata.\allowbreak{}words()} and look for the families: \code{keep\_\allowbreak{}} and \code{remove\_\allowbreak{}}, ...\_x and ...\_y, \code{add\_\allowbreak{}}...\_column, ...\_of.
  \item Then do the task: guess first, check with \code{sdata.\allowbreak{}help} after.
\end{itemize}
\field{In the manual}{lesson~3, \manualpage{lesson:3}{16}{chapter.II.3}; ``How the words are built'', \manualpage{start:grammar}{5}{section*.7}.}
\field{Words to look up}{\kind{chart}\,\code{chart.\allowbreak{}colour()}~\manualpage{word:colour}{140}{section*.83}; \kind{chart}\,\code{chart.\allowbreak{}colours()}~\manualpage{word:colours}{141}{section*.84}; \kind{chart}\,\code{chart.\allowbreak{}legend\_\allowbreak{}bottom()}~\manualpage{word:legend-bottom}{144}{section*.91}; \kind{chart}\,\code{chart.\allowbreak{}order\_\allowbreak{}x()}~\manualpage{word:order-x}{146}{section*.99}; \kind{chart}\,\code{chart.\allowbreak{}save()}~\manualpage{word:save}{112}{section*.18}; \kind{chart}\,\code{chart.\allowbreak{}hide\_\allowbreak{}legend()}~\manualpage{word:hide-legend}{145}{section*.94}; \kind{chart}\,\code{chart.\allowbreak{}legend\_\allowbreak{}left()}~\manualpage{word:legend-left}{144}{section*.92}; \kind{sdata}\,\code{sdata.\allowbreak{}words()}~\manualpage{word:sdata.words}{238}{section*.332}.}
\end{todo}
\begin{tryit}
\begin{enumerate}
  \item Weather table: draw lines of \code{\textquotedbl{}temperature\textquotedbl{}} over \code{\textquotedbl{}month\textquotedbl{}}, one colour for each city. Put the legend on top and use the style \code{\textquotedbl{}paper\textquotedbl{}}.
  \item Save one chart as \code{.\allowbreak{}png}, \code{.\allowbreak{}pdf} and \code{.\allowbreak{}svg}. Then guess, before you look anything up: which word hides the legend, and which puts it on the left? Try your guesses on a chart with \code{colour=}.
\end{enumerate}
{\small Worked answers: \hyperref[answers:python:2]{page~\pageref*{answers:python:2}}, or \code{sdata.\allowbreak{}answers(2)}.}
\end{tryit}

\daystart{3}{Every kind of chart}{Python track}{day:python:3}
\goal{Draw lines, dots, histograms, boxes, violins, pies and areas, and know when to use each. About an hour.}
\begin{warmup}
\begin{enumerate}
  \item How do you give each class its own colour in a bar chart? \fromday{2}\writeline
  \item Which word moves the legend below the chart? \fromday{2}\writeline
  \item How do you save a chart as a PDF? \fromday{1}\writeline
\end{enumerate}
{\small Answers: \hyperref[answers:python:3]{page~\pageref*{answers:python:3}}, or \code{sdata.\allowbreak{}answers(3)}.}
\end{warmup}
\begin{todo}
\begin{itemize}
  \item Work through \code{sdata.\allowbreak{}lesson(4)}: type each example yourself (do not copy and paste), run it and read what each line does. Then do the lesson's tasks and check them with \code{sdata.\allowbreak{}lesson(4,\allowbreak{} solutions=True)}.
\end{itemize}
\field{In the manual}{lesson~4, \manualpage{lesson:4}{21}{chapter.II.4}.}
\field{Words to look up}{\kind{chart}\,\code{chart.\allowbreak{}lines()}~\manualpage{word:lines}{113}{section*.21}; \kind{chart}\,\code{chart.\allowbreak{}dots()}~\manualpage{word:dots}{114}{section*.22}; \kind{chart}\,\code{chart.\allowbreak{}histogram()}~\manualpage{word:histogram}{115}{section*.24}; \kind{chart}\,\code{chart.\allowbreak{}boxes()}~\manualpage{word:boxes}{116}{section*.25}.}
\end{todo}
\begin{tryit}
\begin{enumerate}
  \item Draw violins of \code{\textquotedbl{}temperature\textquotedbl{}} for each city in the weather table, with all the dots on top.
\end{enumerate}
{\small Worked answers: \hyperref[answers:python:3]{page~\pageref*{answers:python:3}}, or \code{sdata.\allowbreak{}answers(3)}.}
\end{tryit}

\daystart{4}{Statistics on charts, and axes}{Python track}{day:python:4}
\goal{Show averages with error bars and fit lines, and control the axes with ranges and reference lines. About an hour.}
\begin{warmup}
\begin{enumerate}
  \item Which chart shows how the values of one column are spread, in bins? \fromday{3}\writeline
  \item Which word puts a title above the chart? \fromday{1}\writeline
  \item Which two words show the spread of each group, with every value as a dot on top? \fromday{3}\writeline
\end{enumerate}
{\small Answers: \hyperref[answers:python:4]{page~\pageref*{answers:python:4}}, or \code{sdata.\allowbreak{}answers(4)}.}
\end{warmup}
\begin{todo}
\begin{itemize}
  \item Work through \code{sdata.\allowbreak{}lesson(5)}: type each example yourself (do not copy and paste), run it and read what each line does. Then do the lesson's tasks and check them with \code{sdata.\allowbreak{}lesson(5,\allowbreak{} solutions=True)}.
  \item Work through \code{sdata.\allowbreak{}lesson(6)}: type each example yourself (do not copy and paste), run it and read what each line does. Then do the lesson's tasks and check them with \code{sdata.\allowbreak{}lesson(6,\allowbreak{} solutions=True)}.
\end{itemize}
\field{In the manual}{lesson~5, \manualpage{lesson:5}{26}{chapter.II.5}; lesson~6, \manualpage{lesson:6}{31}{chapter.II.6}; ``Standard deviation or standard error?'', \manualpage{ideas:sd-se}{244}{section*.348}.}
\field{Words to look up}{\kind{option}\,\code{show=}~\manualpage{word:show-eq}{121}{section*.36}; \kind{option}\,\code{error=}~\manualpage{word:error-eq}{121}{section*.37}; \kind{chart}\,\code{chart.\allowbreak{}numbers\_\allowbreak{}on\_\allowbreak{}bars()}~\manualpage{word:numbers-on-bars}{133}{section*.64}; \kind{chart}\,\code{chart.\allowbreak{}fit\_\allowbreak{}line()}~\manualpage{word:fit-line}{122}{section*.39}; \kind{chart}\,\code{chart.\allowbreak{}range\_\allowbreak{}y()}~\manualpage{word:range-y}{136}{section*.70}; \kind{chart}\,\code{chart.\allowbreak{}line\_\allowbreak{}at\_\allowbreak{}y()}~\manualpage{word:line-at-y}{130}{section*.54}; \kind{chart}\,\code{chart.\allowbreak{}shade\_\allowbreak{}between\_\allowbreak{}y()}~\manualpage{word:shade-between-y}{131}{section*.57}; \kind{chart}\,\code{chart.\allowbreak{}log\_\allowbreak{}scale\_\allowbreak{}y()}~\manualpage{word:log-scale-y}{136}{section*.72}.}
\end{todo}
\begin{tryit}
\begin{enumerate}
  \item Weather table: draw bars of the average \code{\textquotedbl{}rain\textquotedbl{}} for each city with 95 \% confidence intervals, and write the numbers on the bars.
  \item Growth table: draw lines of \code{\textquotedbl{}population\textquotedbl{}} over \code{\textquotedbl{}year\textquotedbl{}}, one for each country, with a line at x = 2010 and the y-axis starting at 0.
\end{enumerate}
{\small Worked answers: \hyperref[answers:python:4]{page~\pageref*{answers:python:4}}, or \code{sdata.\allowbreak{}answers(4)}.}
\end{tryit}

\daystart{5}{Several charts, and a figure for a journal}{Python track}{day:python:5}
\goal{Put charts side by side or in a grid, and save a figure at the size and resolution a journal asks for. About an hour.}
\begin{warmup}
\begin{enumerate}
  \item Which option makes a bar show the median instead of the average? \fromday{4}\writeline
  \item Which word lets you choose the colours of the groups? \fromday{2}\writeline
  \item Which option adds error bars that say how sure we are of the average? \fromday{4}\writeline
\end{enumerate}
{\small Answers: \hyperref[answers:python:5]{page~\pageref*{answers:python:5}}, or \code{sdata.\allowbreak{}answers(5)}.}
\end{warmup}
\begin{todo}
\begin{itemize}
  \item Work through \code{sdata.\allowbreak{}lesson(7)}: type each example yourself (do not copy and paste), run it and read what each line does. Then do the lesson's tasks and check them with \code{sdata.\allowbreak{}lesson(7,\allowbreak{} solutions=True)}.
  \item Work through \code{sdata.\allowbreak{}lesson(8)}: type each example yourself (do not copy and paste), run it and read what each line does. Then do the lesson's tasks and check them with \code{sdata.\allowbreak{}lesson(8,\allowbreak{} solutions=True)}.
\end{itemize}
\field{In the manual}{lesson~7, \manualpage{lesson:7}{36}{chapter.II.7}; lesson~8, \manualpage{lesson:8}{41}{chapter.II.8}.}
\field{Words to look up}{\kind{chart}\,\code{chart.\allowbreak{}one\_\allowbreak{}chart\_\allowbreak{}per()}~\manualpage{word:one-chart-per}{148}{section*.103}; \kind{sdata}\,\code{sdata.\allowbreak{}side\_\allowbreak{}by\_\allowbreak{}side()}~\manualpage{word:sdata.side-by-side}{148}{section*.104}; \kind{sdata}\,\code{sdata.\allowbreak{}one\_\allowbreak{}above\_\allowbreak{}other()}~\manualpage{word:sdata.one-above-other}{149}{section*.105}; \kind{sdata}\,\code{sdata.\allowbreak{}grid()}~\manualpage{word:sdata.grid}{149}{section*.106}; \kind{chart}\,\code{chart.\allowbreak{}second\_\allowbreak{}y()}~\manualpage{word:second-y}{139}{section*.79}; \kind{chart}\,\code{chart.\allowbreak{}size()}~\manualpage{word:size}{142}{section*.87}; \kind{chart}\,\code{chart.\allowbreak{}style()}~\manualpage{word:style}{142}{section*.86}; \kind{chart}\,\code{chart.\allowbreak{}note()}~\manualpage{word:note}{133}{section*.63}.}
\end{todo}
\begin{tryit}
\begin{enumerate}
  \item Make a histogram and a box chart of the marks, and put them one above the other on one page.
  \item A figure for a journal: dots of \code{\textquotedbl{}hours\textquotedbl{}} against \code{\textquotedbl{}marks\textquotedbl{}}, style \code{\textquotedbl{}paper\textquotedbl{}}, 8.5 cm wide and 7 cm high, font size 8, saved as a \code{.\allowbreak{}tif} with 600 dpi.
\end{enumerate}
{\small Worked answers: \hyperref[answers:python:5]{page~\pageref*{answers:python:5}}, or \code{sdata.\allowbreak{}answers(5)}.}
\end{tryit}

\daystart{6}{Open a table, choose rows and columns}{Python track}{day:python:6}
\goal{Look at a table, keep and remove rows and columns (by value and by position), sort, and undo. About an hour.}
\begin{warmup}
\begin{enumerate}
  \item Which word draws one small chart for each class? \fromday{5}\writeline
  \item Which word draws one dot for each row, x against y? \fromday{3}\writeline
  \item Which word puts two charts next to each other on one page? \fromday{5}\writeline
\end{enumerate}
{\small Answers: \hyperref[answers:python:6]{page~\pageref*{answers:python:6}}, or \code{sdata.\allowbreak{}answers(6)}.}
\end{warmup}
\begin{todo}
\begin{itemize}
  \item Work through \code{sdata.\allowbreak{}lesson(9)}: type each example yourself (do not copy and paste), run it and read what each line does. Then do the lesson's tasks and check them with \code{sdata.\allowbreak{}lesson(9,\allowbreak{} solutions=True)}.
  \item Work through \code{sdata.\allowbreak{}lesson(10)}: type each example yourself (do not copy and paste), run it and read what each line does. Then do the lesson's tasks and check them with \code{sdata.\allowbreak{}lesson(10,\allowbreak{} solutions=True)}.
\end{itemize}
\field{In the manual}{lesson~9, \manualpage{lesson:9}{48}{chapter.II.9}; lesson~10, \manualpage{lesson:10}{52}{chapter.II.10}.}
\field{Words to look up}{\kind{sdata}\,\code{sdata.\allowbreak{}table()}~\manualpage{word:sdata.table}{151}{section*.112}; \kind{table}\,\code{table.\allowbreak{}show\_\allowbreak{}first\_\allowbreak{}rows()}~\manualpage{word:table.show-first-rows}{152}{section*.113}; \kind{table}\,\code{table.\allowbreak{}describe\_\allowbreak{}columns()}~\manualpage{word:table.describe-columns}{153}{section*.117}; \kind{table}\,\code{table.\allowbreak{}count\_\allowbreak{}rows()}~\manualpage{word:table.count-rows}{176}{section*.171}; \kind{table}\,\code{table.\allowbreak{}keep\_\allowbreak{}rows()}~\manualpage{word:table.keep-rows}{157}{section*.128}; \kind{table}\,\code{table.\allowbreak{}keep\_\allowbreak{}columns()}~\manualpage{word:table.keep-columns}{155}{section*.122}; \kind{table}\,\code{table.\allowbreak{}sort\_\allowbreak{}by()}~\manualpage{word:table.sort-by}{169}{section*.157}; \kind{table}\,\code{table.\allowbreak{}undo()}~\manualpage{word:table.undo}{211}{section*.250}.}
\end{todo}
\begin{tryit}
\begin{enumerate}
  \item Open the \code{\textquotedbl{}shop\textquotedbl{}} practice table. How many rows does it have, and how many rows does each shop have?
  \item Weather table: keep the city \code{\textquotedbl{}Sunnyvale\textquotedbl{}}, keep only the columns month, temperature and rain, sort by rain with the biggest first, and save it as \code{\textquotedbl{}sunny.\allowbreak{}csv\textquotedbl{}}.
\end{enumerate}
{\small Worked answers: \hyperref[answers:python:6]{page~\pageref*{answers:python:6}}, or \code{sdata.\allowbreak{}answers(6)}.}
\end{tryit}

\daystart{7}{Review of the first week}{Python track}{day:python:7}
\goal{Remember the words of days 1 to 6 without looking, and use them together in one program. About an hour.}
\begin{warmup}
\begin{enumerate}
  \item How do you look at the first 5 rows of a table? \fromday{6}\writeline
  \item How do you name the y-axis? \fromday{1}\writeline
  \item Which word draws a straight line through the dots? \fromday{4}\writeline
  \item How do you put the bars in your own order? \fromday{2}\writeline
  \item How do you put four charts in two rows, with the letters A to D? \fromday{5}\writeline
  \item Which word draws a line with a dot at each point? \fromday{3}\writeline
\end{enumerate}
{\small Answers: \hyperref[answers:python:7]{page~\pageref*{answers:python:7}}, or \code{sdata.\allowbreak{}answers(7)}.}
\end{warmup}
\begin{todo}
\begin{itemize}
  \item Answer the warm-up on paper, from memory, before you look anything up.
  \item For each question you could not answer, open its day with \code{sdata.\allowbreak{}day(n)} and run its example again.
  \item Then do the mixed task without looking back. It uses words from the whole week.
\end{itemize}
\end{todo}
\begin{tryit}
\begin{enumerate}
  \item Marks table: keep only class \code{\textquotedbl{}A\textquotedbl{}}, draw bars of the average \code{\textquotedbl{}marks\textquotedbl{}} for each subject with standard error bars, give the chart a title and save it as a PDF.
\end{enumerate}
{\small Worked answers: \hyperref[answers:python:7]{page~\pageref*{answers:python:7}}, or \code{sdata.\allowbreak{}answers(7)}.}
\end{tryit}

\daystart{8}{Clean the data and add columns}{Python track}{day:python:8}
\goal{Find and fill empty cells, fix text and numbers, and add new columns, groups and dates. About an hour.}
\begin{warmup}
\begin{enumerate}
  \item Which word describes each column: its kind, empty cells and numbers? \fromday{6}\writeline
  \item Which word adds a second y-axis on the right? \fromday{5}\writeline
  \item The names under the bars bump into each other. Which word helps? \fromday{1}\writeline
\end{enumerate}
{\small Answers: \hyperref[answers:python:8]{page~\pageref*{answers:python:8}}, or \code{sdata.\allowbreak{}answers(8)}.}
\end{warmup}
\begin{todo}
\begin{itemize}
  \item Work through \code{sdata.\allowbreak{}lesson(11)}: type each example yourself (do not copy and paste), run it and read what each line does. Then do the lesson's tasks and check them with \code{sdata.\allowbreak{}lesson(11,\allowbreak{} solutions=True)}.
\end{itemize}
\field{In the manual}{lesson~11, \manualpage{lesson:11}{57}{chapter.II.11}; ``Missing values and imputation'', \manualpage{ideas:missing}{248}{section*.365}.}
\field{Words to look up}{\kind{table}\,\code{table.\allowbreak{}count\_\allowbreak{}empty\_\allowbreak{}cells()}~\manualpage{word:table.count-empty-cells}{163}{section*.142}; \kind{table}\,\code{table.\allowbreak{}fill\_\allowbreak{}empty\_\allowbreak{}cells()}~\manualpage{word:table.fill-empty-cells}{164}{section*.144}; \kind{table}\,\code{table.\allowbreak{}make\_\allowbreak{}numbers()}~\manualpage{word:table.make-numbers}{165}{section*.147}; \kind{table}\,\code{table.\allowbreak{}add\_\allowbreak{}column()}~\manualpage{word:table.add-column}{170}{section*.158}.}
\end{todo}
\begin{tryit}
\begin{enumerate}
  \item Shop table: count the empty cells, fill the empty \code{\textquotedbl{}sales\textquotedbl{}} with the line between neighbours for each shop (sorted by \code{\textquotedbl{}date\textquotedbl{}}), and add a column \code{\textquotedbl{}weekday\textquotedbl{}} with the name of the weekday.
\end{enumerate}
{\small Worked answers: \hyperref[answers:python:8]{page~\pageref*{answers:python:8}}, or \code{sdata.\allowbreak{}answers(8)}.}
\end{tryit}

\daystart{9}{Summaries, numbers and grids}{Python track}{day:python:9}
\goal{Sum up columns for each group, join and reshape tables, and work with lists and grids of numbers. About an hour.}
\begin{warmup}
\begin{enumerate}
  \item Which word counts the empty cells in each column? \fromday{8}\writeline
  \item Which word counts the rows of each fruit? \fromday{6}\writeline
  \item If \code{keep\_\allowbreak{}rows} keeps rows, what would you guess takes columns away? \fromday{2}\writeline
\end{enumerate}
{\small Answers: \hyperref[answers:python:9]{page~\pageref*{answers:python:9}}, or \code{sdata.\allowbreak{}answers(9)}.}
\end{warmup}
\begin{todo}
\begin{itemize}
  \item Work through \code{sdata.\allowbreak{}lesson(12)}: type each example yourself (do not copy and paste), run it and read what each line does. Then do the lesson's tasks and check them with \code{sdata.\allowbreak{}lesson(12,\allowbreak{} solutions=True)}.
  \item Work through \code{sdata.\allowbreak{}lesson(13)}: type each example yourself (do not copy and paste), run it and read what each line does. Then do the lesson's tasks and check them with \code{sdata.\allowbreak{}lesson(13,\allowbreak{} solutions=True)}.
\end{itemize}
\field{In the manual}{lesson~12, \manualpage{lesson:12}{64}{chapter.II.12}; lesson~13, \manualpage{lesson:13}{72}{chapter.II.13}.}
\field{Words to look up}{\kind{table}\,\code{table.\allowbreak{}average\_\allowbreak{}of()}~\manualpage{word:table.average-of}{177}{section*.173}; \kind{table}\,\code{table.\allowbreak{}summary\_\allowbreak{}of()}~\manualpage{word:table.summary-of}{179}{section*.179}; \kind{table}\,\code{table.\allowbreak{}correlation\_\allowbreak{}between()}~\manualpage{word:table.correlation-between}{182}{section*.184}; \kind{table}\,\code{table.\allowbreak{}join\_\allowbreak{}with()}~\manualpage{word:table.join-with}{209}{section*.245}; \kind{sdata}\,\code{sdata.\allowbreak{}numbers()}~\manualpage{word:sdata.numbers}{214}{section*.260}; \kind{sdata}\,\code{sdata.\allowbreak{}random\_\allowbreak{}numbers()}~\manualpage{word:sdata.random-numbers}{214}{section*.261}; \kind{nums}\,\code{nums.\allowbreak{}median()}~\manualpage{word:nums.median}{226}{section*.296}; \kind{nums}\,\code{nums.\allowbreak{}percentile()}~\manualpage{word:nums.percentile}{228}{section*.301}.}
\end{todo}
\begin{tryit}
\begin{enumerate}
  \item Fruit table: the total \code{\textquotedbl{}count\textquotedbl{}} of each fruit, drawn as bars with the numbers on the bars.
  \item Make 200 random heights with average 150 and standard deviation 8 (\code{seed=1}). What is their 90th percentile? Draw them as a histogram.
\end{enumerate}
{\small Worked answers: \hyperref[answers:python:9]{page~\pageref*{answers:python:9}}, or \code{sdata.\allowbreak{}answers(9)}.}
\end{tryit}

\daystart{10}{Is a difference real? Tests and effect sizes}{Python track}{day:python:10}
\goal{Test whether averages differ, find which groups differ, and say how big the difference is. About an hour.}
\begin{warmup}
\begin{enumerate}
  \item Which word gives the average of a column for each group? \fromday{9}\writeline
  \item Which word counts the rows? \fromday{6}\writeline
  \item Which chart shows how the values of one column are spread, in bins? \fromday{3}\writeline
\end{enumerate}
{\small Answers: \hyperref[answers:python:10]{page~\pageref*{answers:python:10}}, or \code{sdata.\allowbreak{}answers(10)}.}
\end{warmup}
\begin{todo}
\begin{itemize}
  \item Work through \code{sdata.\allowbreak{}lesson(14)}, examples 1 and 3: type each example yourself (do not copy and paste), run it and read what each line does. Then do the lesson's tasks and check them with \code{sdata.\allowbreak{}lesson(14,\allowbreak{} solutions=True)}.
  \item Read \code{\textquotedbl{}The null hypothesis and the p-value\textquotedbl{}} and \code{\textquotedbl{}Effect sizes\textquotedbl{}} in the manual.
\end{itemize}
\field{In the manual}{lesson~14, \manualpage{lesson:14}{76}{chapter.II.14}; ``The null hypothesis and the p-value'', \manualpage{ideas:p-value}{245}{section*.353}; ``Effect sizes: how big is the difference?'', \manualpage{ideas:effect-size}{246}{section*.356}.}
\field{Words to look up}{\kind{table}\,\code{table.\allowbreak{}compare\_\allowbreak{}averages()}~\manualpage{word:table.compare-averages}{183}{section*.186}; \kind{table}\,\code{table.\allowbreak{}compare\_\allowbreak{}each\_\allowbreak{}pair()}~\manualpage{word:table.compare-each-pair}{183}{section*.187}; \kind{table}\,\code{table.\allowbreak{}fit\_\allowbreak{}line()}~\manualpage{word:table.fit-line}{188}{section*.197}.}
\end{todo}
\begin{tryit}
\begin{enumerate}
  \item School table: do the average marks differ between pupils who had breakfast and pupils who did not? Write down the difference, the effect size and the p-value.
\end{enumerate}
{\small Worked answers: \hyperref[answers:python:10]{page~\pageref*{answers:python:10}}, or \code{sdata.\allowbreak{}answers(10)}.}
\end{tryit}

\daystart{11}{Counts, outliers and many tests}{Python track}{day:python:11}
\goal{Test counts of categories, check shapes and outliers, and correct the p-values of many tests. About an hour.}
\begin{warmup}
\begin{enumerate}
  \item Which word tests whether the average marks differ between the classes? \fromday{10}\writeline
  \item How do you fill the empty marks with the median? \fromday{8}\writeline
  \item Which word checks by eye whether values are bell-shaped? \fromday{4}\writeline
\end{enumerate}
{\small Answers: \hyperref[answers:python:11]{page~\pageref*{answers:python:11}}, or \code{sdata.\allowbreak{}answers(11)}.}
\end{warmup}
\begin{todo}
\begin{itemize}
  \item Work through \code{sdata.\allowbreak{}lesson(14)}, examples 2, 4 and 5: type each example yourself (do not copy and paste), run it and read what each line does. Then do the lesson's tasks and check them with \code{sdata.\allowbreak{}lesson(14,\allowbreak{} solutions=True)}.
  \item Read \code{\textquotedbl{}Many tests at once\textquotedbl{}} in the manual.
\end{itemize}
\field{In the manual}{lesson~14, \manualpage{lesson:14}{76}{chapter.II.14}; ``Many tests at once'', \manualpage{ideas:multiple}{246}{section*.355}.}
\field{Words to look up}{\kind{table}\,\code{table.\allowbreak{}compare\_\allowbreak{}counts()}~\manualpage{word:table.compare-counts}{185}{section*.190}; \kind{table}\,\code{table.\allowbreak{}show\_\allowbreak{}rows\_\allowbreak{}with\_\allowbreak{}outliers()}~\manualpage{word:table.show-rows-with-outliers}{162}{section*.139}; \kind{table}\,\code{table.\allowbreak{}compare\_\allowbreak{}shape\_\allowbreak{}to()}~\manualpage{word:table.compare-shape-to}{185}{section*.191}; \kind{table}\,\code{table.\allowbreak{}add\_\allowbreak{}adjusted\_\allowbreak{}p\_\allowbreak{}value\_\allowbreak{}column()}~\manualpage{word:table.add-adjusted-p-value-column}{186}{section*.192}.}
\end{todo}
\begin{tryit}
\begin{enumerate}
  \item School table: is the share of pupils who \code{\textquotedbl{}passed\textquotedbl{}} different with and without breakfast?
\end{enumerate}
{\small Worked answers: \hyperref[answers:python:11]{page~\pageref*{answers:python:11}}, or \code{sdata.\allowbreak{}answers(11)}.}
\end{tryit}

\daystart{12}{Choosing a test, and writing it up}{Python track}{day:python:12}
\goal{Choose the right test with four questions, and write a results sentence a reader can check. About an hour.}
\begin{warmup}
\begin{enumerate}
  \item Which word tests whether the counts of categories differ between groups? \fromday{11}\writeline
  \item Which word tells whether two columns go up and down together? \fromday{9}\writeline
  \item Which word sets the size of a figure in cm? \fromday{5}\writeline
\end{enumerate}
{\small Answers: \hyperref[answers:python:12]{page~\pageref*{answers:python:12}}, or \code{sdata.\allowbreak{}answers(12)}.}
\end{warmup}
\begin{todo}
\begin{itemize}
  \item Read \code{\textquotedbl{}Which test when?\textquotedbl{}} and \code{\textquotedbl{}Writing it up\textquotedbl{}} in the manual. Answer its four questions for each task below before you run anything.
  \item Run \code{sdata.\allowbreak{}help(\textquotedbl{}compare\_\allowbreak{}before\_\allowbreak{}after\textquotedbl{})} and \code{sdata.\allowbreak{}help(\textquotedbl{}compare\_\allowbreak{}average\_\allowbreak{}to\textquotedbl{})} and type their examples.
\end{itemize}
\field{In the manual}{``Which test when?'', \manualpage{ideas:which-test}{249}{section*.367}; ``Writing it up'', \manualpage{ideas:report}{250}{section*.368}.}
\field{Words to look up}{\kind{table}\,\code{table.\allowbreak{}compare\_\allowbreak{}average\_\allowbreak{}to()}~\manualpage{word:table.compare-average-to}{184}{section*.189}; \kind{table}\,\code{table.\allowbreak{}compare\_\allowbreak{}before\_\allowbreak{}after()}~\manualpage{word:table.compare-before-after}{184}{section*.188}; \kind{table}\,\code{table.\allowbreak{}compare\_\allowbreak{}percent\_\allowbreak{}to()}~\manualpage{word:table.compare-percent-to}{187}{section*.195}.}
\end{todo}
\begin{tryit}
\begin{enumerate}
  \item Marks table: is each child's Maths mark different from the same child's Science mark? Choose the right word first.
  \item School table: is the average \code{\textquotedbl{}sleep\textquotedbl{}} different from 8 hours?
\end{enumerate}
{\small Worked answers: \hyperref[answers:python:12]{page~\pageref*{answers:python:12}}, or \code{sdata.\allowbreak{}answers(12)}.}
\end{tryit}

\daystart{13}{A whole analysis}{Python track}{day:python:13}
\goal{Go from a question and a file to a tested answer and a figure with panels A and B. About an hour.}
\begin{warmup}
\begin{enumerate}
  \item The same children were measured twice. Which word compares the two? \fromday{12}\writeline
  \item Which word says which pairs of groups differ? \fromday{10}\writeline
  \item Which word keeps only the rows where subject is Maths? \fromday{6}\writeline
\end{enumerate}
{\small Answers: \hyperref[answers:python:13]{page~\pageref*{answers:python:13}}, or \code{sdata.\allowbreak{}answers(13)}.}
\end{warmup}
\begin{todo}
\begin{itemize}
  \item Work through \code{sdata.\allowbreak{}lesson(15)}: type each example yourself (do not copy and paste), run it and read what each line does. Then do the lesson's tasks and check them with \code{sdata.\allowbreak{}lesson(15,\allowbreak{} solutions=True)}.
\end{itemize}
\field{In the manual}{lesson~15, \manualpage{lesson:15}{81}{chapter.II.15}.}
\field{Words to look up}{\kind{table}\,\code{table.\allowbreak{}steps()}~\manualpage{word:table.steps}{211}{section*.251}; \kind{sdata}\,\code{sdata.\allowbreak{}grid()}~\manualpage{word:sdata.grid}{149}{section*.106}; \kind{table}\,\code{table.\allowbreak{}save()}~\manualpage{word:table.save}{212}{section*.253}; \kind{table}\,\code{table.\allowbreak{}as\_\allowbreak{}pandas()}~\manualpage{word:table.as-pandas}{212}{section*.255}.}
\end{todo}
\begin{tryit}
\begin{enumerate}
  \item From memory: weather table, compare the average \code{\textquotedbl{}temperature\textquotedbl{}} of the cities, draw boxes with all the dots, and save the figure as a PDF and the table as \code{\textquotedbl{}weather.\allowbreak{}xlsx\textquotedbl{}}.
\end{enumerate}
{\small Worked answers: \hyperref[answers:python:13]{page~\pageref*{answers:python:13}}, or \code{sdata.\allowbreak{}answers(13)}.}
\end{tryit}

\daystart{14}{Review of the second week}{Python track}{day:python:14}
\goal{Remember the table and test words of days 6 to 13 without looking, and use them together. About an hour.}
\begin{warmup}
\begin{enumerate}
  \item Which word lists every change you made to a table, for your methods? \fromday{13}\writeline
  \item Which word makes all the text on a chart bigger or smaller? \fromday{1}\writeline
  \item A word for one axis ends in ...? \fromday{2}\writeline
  \item Which word turns a column of text such as \code{\textquotedbl{}144\textquotedbl{}} into numbers? \fromday{8}\writeline
  \item Which word sets the y-axis from 0 to 100? \fromday{4}\writeline
  \item Which word shows the rows with outliers? \fromday{11}\writeline
  \item Which word joins two tables on a column they share? \fromday{9}\writeline
  \item How do you save a picture sharp enough for a journal? \fromday{5}\writeline
\end{enumerate}
{\small Answers: \hyperref[answers:python:14]{page~\pageref*{answers:python:14}}, or \code{sdata.\allowbreak{}answers(14)}.}
\end{warmup}
\begin{todo}
\begin{itemize}
  \item Answer the warm-up on paper, from memory, before you look anything up.
  \item For each question you could not answer, open its day with \code{sdata.\allowbreak{}day(n)} and run its example again.
  \item Then do the mixed task without looking back.
\end{itemize}
\end{todo}
\begin{tryit}
\begin{enumerate}
  \item School table: count the empty cells, remove the rows with an empty \code{\textquotedbl{}sleep\textquotedbl{}}, give the average sleep for each value of \code{\textquotedbl{}passed\textquotedbl{}}, test the difference, and draw boxes with all the dots.
\end{enumerate}
{\small Worked answers: \hyperref[answers:python:14]{page~\pageref*{answers:python:14}}, or \code{sdata.\allowbreak{}answers(14)}.}
\end{tryit}

\daystart{15}{Models: explain one column from several}{Python track}{day:python:15}
\goal{Explain one column from several others: multiple, logistic, Poisson and mixed regression. About an hour.}
\begin{warmup}
\begin{enumerate}
  \item How do you give the panels of a figure the letters A and B? \fromday{13}\writeline
  \item Which word compares one average with a number you choose? \fromday{12}\writeline
  \item Which word adds a new column made from other columns? \fromday{8}\writeline
  \item How do you start a table from the marks practice table? \fromday{1}\writeline
\end{enumerate}
{\small Answers: \hyperref[answers:python:15]{page~\pageref*{answers:python:15}}, or \code{sdata.\allowbreak{}answers(15)}.}
\end{warmup}
\begin{todo}
\begin{itemize}
  \item Work through \code{sdata.\allowbreak{}lesson(16)}: type each example yourself (do not copy and paste), run it and read what each line does. Then do the lesson's tasks and check them with \code{sdata.\allowbreak{}lesson(16,\allowbreak{} solutions=True)}.
\end{itemize}
\field{In the manual}{lesson~16, \manualpage{lesson:16}{86}{chapter.II.16}.}
\field{Words to look up}{\kind{table}\,\code{table.\allowbreak{}explain()}~\manualpage{word:table.explain}{200}{section*.227}; \kind{table}\,\code{table.\allowbreak{}add\_\allowbreak{}predicted\_\allowbreak{}column()}~\manualpage{word:table.add-predicted-column}{200}{section*.228}.}
\end{todo}
\begin{tryit}
\begin{enumerate}
  \item School table: explain \code{\textquotedbl{}days\_\allowbreak{}absent\textquotedbl{}} using \code{\textquotedbl{}sleep\textquotedbl{}} and \code{\textquotedbl{}breakfast\textquotedbl{}} with \code{outcome=\textquotedbl{}counts\textquotedbl{}}. Does breakfast matter?
\end{enumerate}
{\small Worked answers: \hyperref[answers:python:15]{page~\pageref*{answers:python:15}}, or \code{sdata.\allowbreak{}answers(15)}.}
\end{tryit}

\daystart{16}{Patterns across many columns}{Python track}{day:python:16}
\goal{Find patterns across many columns with principal components, clusters, NMDS and PERMANOVA. About an hour.}
\begin{warmup}
\begin{enumerate}
  \item Which word explains marks from hours, sleep and breakfast together? \fromday{15}\writeline
  \item Which word hands a table on to pandas? \fromday{13}\writeline
  \item Which word gives several summaries at once? \fromday{9}\writeline
  \item Which file endings give a picture that stays sharp at any size? \fromday{2}\writeline
\end{enumerate}
{\small Answers: \hyperref[answers:python:16]{page~\pageref*{answers:python:16}}, or \code{sdata.\allowbreak{}answers(16)}.}
\end{warmup}
\begin{todo}
\begin{itemize}
  \item Work through \code{sdata.\allowbreak{}lesson(17)}: type each example yourself (do not copy and paste), run it and read what each line does. Then do the lesson's tasks and check them with \code{sdata.\allowbreak{}lesson(17,\allowbreak{} solutions=True)}.
\end{itemize}
\field{In the manual}{lesson~17, \manualpage{lesson:17}{90}{chapter.II.17}.}
\field{Words to look up}{\kind{table}\,\code{table.\allowbreak{}add\_\allowbreak{}principal\_\allowbreak{}components()}~\manualpage{word:table.add-principal-components}{201}{section*.230}; \kind{table}\,\code{table.\allowbreak{}add\_\allowbreak{}clusters()}~\manualpage{word:table.add-clusters}{203}{section*.232}; \kind{table}\,\code{table.\allowbreak{}add\_\allowbreak{}nmds()}~\manualpage{word:table.add-nmds}{204}{section*.234}; \kind{table}\,\code{table.\allowbreak{}compare\_\allowbreak{}profiles()}~\manualpage{word:table.compare-profiles}{205}{section*.235}.}
\end{todo}
\begin{tryit}
\begin{enumerate}
  \item Meadow table: add 3 clusters and count the rows of each cluster across \code{\textquotedbl{}field\textquotedbl{}}.
\end{enumerate}
{\small Worked answers: \hyperref[answers:python:16]{page~\pageref*{answers:python:16}}, or \code{sdata.\allowbreak{}answers(16)}.}
\end{tryit}

\daystart{17}{Survival, and planning a study}{Python track}{day:python:17}
\goal{Draw and test survival curves, and work out how many values a study needs. About an hour.}
\begin{warmup}
\begin{enumerate}
  \item Which word adds the principal components PC1 and PC2? \fromday{16}\writeline
  \item Which word lists every change you made to a table, for your methods? \fromday{13}\writeline
  \item A p-value answers \code{\textquotedbl{}could it be chance?\textquotedbl{}}. What answers \code{\textquotedbl{}how big is it?\textquotedbl{}} \fromday{10}\writeline
  \item Which two words show the spread of each group, with every value as a dot on top? \fromday{3}\writeline
\end{enumerate}
{\small Answers: \hyperref[answers:python:17]{page~\pageref*{answers:python:17}}, or \code{sdata.\allowbreak{}answers(17)}.}
\end{warmup}
\begin{todo}
\begin{itemize}
  \item Work through \code{sdata.\allowbreak{}lesson(18)}: type each example yourself (do not copy and paste), run it and read what each line does. Then do the lesson's tasks and check them with \code{sdata.\allowbreak{}lesson(18,\allowbreak{} solutions=True)}.
\end{itemize}
\field{In the manual}{lesson~18, \manualpage{lesson:18}{95}{chapter.II.18}; ``Two kinds of mistake, and power'', \manualpage{ideas:errors}{246}{section*.354}.}
\field{Words to look up}{\kind{table}\,\code{table.\allowbreak{}survival\_\allowbreak{}over\_\allowbreak{}time()}~\manualpage{word:table.survival-over-time}{205}{section*.237}; \kind{chart}\,\code{chart.\allowbreak{}survival\_\allowbreak{}curves()}~\manualpage{word:survival-curves}{207}{section*.240}; \kind{table}\,\code{table.\allowbreak{}compare\_\allowbreak{}survival()}~\manualpage{word:table.compare-survival}{206}{section*.238}; \kind{sdata}\,\code{sdata.\allowbreak{}how\_\allowbreak{}many\_\allowbreak{}needed()}~\manualpage{word:sdata.how-many-needed}{207}{section*.242}.}
\end{todo}
\begin{tryit}
\begin{enumerate}
  \item How many values does each group need to find a large difference (\code{effect\_\allowbreak{}size=0.\allowbreak{}8}) with power 0.8? And with a small one (0.2)?
\end{enumerate}
{\small Worked answers: \hyperref[answers:python:17]{page~\pageref*{answers:python:17}}, or \code{sdata.\allowbreak{}answers(17)}.}
\end{tryit}

\daystart{18}{Fitting lines, curves and formulas}{Python track}{day:python:18}
\goal{Fit straight lines, curves and formulas of your own to a table, and draw them on a chart. About an hour.}
\begin{warmup}
\begin{enumerate}
  \item Which word draws survival curves? \fromday{17}\writeline
  \item How do you explain a yes-or-no column such as \code{\textquotedbl{}passed\textquotedbl{}}? \fromday{15}\writeline
  \item Which word checks whether values are bell-shaped (normal)? \fromday{11}\writeline
  \item Which word draws a line across at y = 50, with a label? \fromday{4}\writeline
\end{enumerate}
{\small Answers: \hyperref[answers:python:18]{page~\pageref*{answers:python:18}}, or \code{sdata.\allowbreak{}answers(18)}.}
\end{warmup}
\begin{todo}
\begin{itemize}
  \item Work through \code{sdata.\allowbreak{}lesson(19)}, examples 1 and 2: type each example yourself (do not copy and paste), run it and read what each line does. Then do the lesson's tasks and check them with \code{sdata.\allowbreak{}lesson(19,\allowbreak{} solutions=True)}.
\end{itemize}
\field{In the manual}{lesson~19, \manualpage{lesson:19}{99}{chapter.II.19}; ``Fitting a line or a formula'', \manualpage{ideas:regression}{247}{section*.361}.}
\field{Words to look up}{\kind{table}\,\code{table.\allowbreak{}fit\_\allowbreak{}line()}~\manualpage{word:table.fit-line}{188}{section*.197}; \kind{table}\,\code{table.\allowbreak{}fit\_\allowbreak{}curve()}~\manualpage{word:table.fit-curve}{189}{section*.198}; \kind{table}\,\code{table.\allowbreak{}fit\_\allowbreak{}formula()}~\manualpage{word:table.fit-formula}{189}{section*.199}; \kind{chart}\,\code{chart.\allowbreak{}fit\_\allowbreak{}formula()}~\manualpage{word:fit-formula}{123}{section*.41}.}
\end{todo}
\begin{tryit}
\begin{enumerate}
  \item Growth table: fit the formula \code{\textquotedbl{}forest = a + b * (year - 2000)\textquotedbl{}} for each country. Which country loses forest fastest?
\end{enumerate}
{\small Worked answers: \hyperref[answers:python:18]{page~\pageref*{answers:python:18}}, or \code{sdata.\allowbreak{}answers(18)}.}
\end{tryit}

\daystart{19}{Finding, solving and distributions}{Python track}{day:python:19}
\goal{Find where a formula is smallest or biggest, solve equations, and work out chances from a distribution. About an hour.}
\begin{warmup}
\begin{enumerate}
  \item Which word fits a straight line on a table? \fromday{18}\writeline
  \item Which word puts rows that are alike into groups? \fromday{16}\writeline
  \item What should a results sentence give? \fromday{12}\writeline
  \item Which word writes a note with an arrow on the chart? \fromday{5}\writeline
\end{enumerate}
{\small Answers: \hyperref[answers:python:19]{page~\pageref*{answers:python:19}}, or \code{sdata.\allowbreak{}answers(19)}.}
\end{warmup}
\begin{todo}
\begin{itemize}
  \item Work through \code{sdata.\allowbreak{}lesson(19)}, example 3: type each example yourself (do not copy and paste), run it and read what each line does. Then do the lesson's tasks and check them with \code{sdata.\allowbreak{}lesson(19,\allowbreak{} solutions=True)}.
  \item Work through \code{sdata.\allowbreak{}lesson(20)}, examples 1 and 2: type each example yourself (do not copy and paste), run it and read what each line does. Then do the lesson's tasks and check them with \code{sdata.\allowbreak{}lesson(20,\allowbreak{} solutions=True)}.
\end{itemize}
\field{In the manual}{lesson~19, \manualpage{lesson:19}{99}{chapter.II.19}; lesson~20, \manualpage{lesson:20}{104}{chapter.II.20}; ``The shape of the values: distributions'', \manualpage{ideas:distribution}{244}{section*.345}.}
\field{Words to look up}{\kind{sdata}\,\code{sdata.\allowbreak{}find\_\allowbreak{}smallest()}~\manualpage{word:sdata.find-smallest}{190}{section*.201}; \kind{sdata}\,\code{sdata.\allowbreak{}find\_\allowbreak{}biggest()}~\manualpage{word:sdata.find-biggest}{191}{section*.202}; \kind{sdata}\,\code{sdata.\allowbreak{}solve()}~\manualpage{word:sdata.solve}{191}{section*.203}; \kind{sdata}\,\code{sdata.\allowbreak{}distribution()}~\manualpage{word:sdata.distribution}{193}{section*.208}; \kind{dist}\,\code{dist.\allowbreak{}chance()}~\manualpage{word:dist.chance}{195}{section*.212}; \kind{dist}\,\code{dist.\allowbreak{}percentile()}~\manualpage{word:dist.percentile}{196}{section*.214}; \kind{table}\,\code{table.\allowbreak{}fit\_\allowbreak{}distribution()}~\manualpage{word:table.fit-distribution}{194}{section*.209}.}
\end{todo}
\begin{tryit}
\begin{enumerate}
  \item A box has a square bottom of side x and a height of 10 - x, so its volume is x * x * (10 - x). For which x between 0 and 10 is the volume biggest?
  \item Heights are normal with average 150 and standard deviation 8. What share is taller than 165? Which height do 95 \% of them stay below?
\end{enumerate}
{\small Worked answers: \hyperref[answers:python:19]{page~\pageref*{answers:python:19}}, or \code{sdata.\allowbreak{}answers(19)}.}
\end{tryit}

\daystart{20}{Values between points}{Python track}{day:python:20}
\goal{Estimate values between the points you measured, and fill an even map for a contour chart. About an hour.}
\begin{warmup}
\begin{enumerate}
  \item Which word finds where a formula is smallest? \fromday{19}\writeline
  \item Which word tests whether survival differs between groups? \fromday{17}\writeline
  \item How do you give the panels of a figure the letters A and B? \fromday{13}\writeline
  \item Which word sorts with the biggest marks first? \fromday{6}\writeline
\end{enumerate}
{\small Answers: \hyperref[answers:python:20]{page~\pageref*{answers:python:20}}, or \code{sdata.\allowbreak{}answers(20)}.}
\end{warmup}
\begin{todo}
\begin{itemize}
  \item Work through \code{sdata.\allowbreak{}lesson(20)}, example 3: type each example yourself (do not copy and paste), run it and read what each line does. Then do the lesson's tasks and check them with \code{sdata.\allowbreak{}lesson(20,\allowbreak{} solutions=True)}.
\end{itemize}
\field{In the manual}{lesson~20, \manualpage{lesson:20}{104}{chapter.II.20}.}
\field{Words to look up}{\kind{table}\,\code{table.\allowbreak{}values\_\allowbreak{}at()}~\manualpage{word:table.values-at}{192}{section*.205}; \kind{table}\,\code{table.\allowbreak{}map\_\allowbreak{}of()}~\manualpage{word:table.map-of}{193}{section*.206}.}
\end{todo}
\begin{tryit}
\begin{enumerate}
  \item Weather table: estimate the \code{\textquotedbl{}rain\textquotedbl{}} of each city at \code{month\_\allowbreak{}number} 2.5 and 7.5, joined by straight lines.
\end{enumerate}
{\small Worked answers: \hyperref[answers:python:20]{page~\pageref*{answers:python:20}}, or \code{sdata.\allowbreak{}answers(20)}.}
\end{tryit}

\daystart{21}{Review, and the question of your project}{Python track}{day:python:21}
\goal{Remember the words of the third week, then choose the question and the raw file of your own project. About an hour.}
\begin{warmup}
\begin{enumerate}
  \item Which word estimates a column between the points you measured? \fromday{20}\writeline
  \item Which word cuts a column of numbers into named groups? \fromday{8}\writeline
  \item Which word keeps only the rows with marks above 50? \fromday{1}\writeline
  \item How do you make simplydata numbers from a list? \fromday{9}\writeline
  \item Which option gives the rank-based test, for small groups that are not bell-shaped? \fromday{10}\writeline
  \item Which word adds the values a model predicts as a new column? \fromday{15}\writeline
  \item You did many tests at once. Which word corrects the p-values? \fromday{11}\writeline
  \item A few values are much bigger than the rest. Which word helps? \fromday{4}\writeline
  \item How do you fit a curve with x and x squared? \fromday{18}\writeline
\end{enumerate}
{\small Answers: \hyperref[answers:python:21]{page~\pageref*{answers:python:21}}, or \code{sdata.\allowbreak{}answers(21)}.}
\end{warmup}
\begin{todo}
\begin{itemize}
  \item Answer the warm-up on paper, from memory, before you look anything up.
  \item For each question you could not answer, open its day with \code{sdata.\allowbreak{}day(n)} and run its example again.
  \item Choose data of your own: a file from your work, your studies or a public data set. No data? Use the \code{\textquotedbl{}school\textquotedbl{}} practice table, as the answers do.
  \item Write your question in one sentence, for example: do pupils who sleep more get better marks?
  \item Keep the raw file as it is: never change it. All the cleaning will be steps on a table, which you can repeat and undo.
  \item Open the file with \code{sdata.\allowbreak{}table} and look at it.
\end{itemize}
\field{In the manual}{``Tidy data: long and wide tables'', \manualpage{ideas:tidy}{248}{section*.364}.}
\field{Words to look up}{\kind{sdata}\,\code{sdata.\allowbreak{}table()}~\manualpage{word:sdata.table}{151}{section*.112}; \kind{table}\,\code{table.\allowbreak{}describe\_\allowbreak{}columns()}~\manualpage{word:table.describe-columns}{153}{section*.117}; \kind{table}\,\code{table.\allowbreak{}save()}~\manualpage{word:table.save}{212}{section*.253}.}
\end{todo}
\begin{tryit}
\begin{enumerate}
  \item Trial table: test whether survival differs between the treatments, then explain the risk of dying from \code{\textquotedbl{}treatment\textquotedbl{}} and \code{\textquotedbl{}age\textquotedbl{}} together.
  \item Write your question in one sentence. Open your raw file and write down how many rows and columns it has, and which columns have empty cells.
\end{enumerate}
{\small Worked answers: \hyperref[answers:python:21]{page~\pageref*{answers:python:21}}, or \code{sdata.\allowbreak{}answers(21)}.}
\end{tryit}

\daystart{22}{Project: clean}{Python track}{day:python:22}
\goal{Clean your raw file into a tidy table, with every step written down. About an hour.}
\begin{warmup}
\begin{enumerate}
  \item Which word tests whether groups differ in their make-up across many columns (PERMANOVA)? \fromday{16}\writeline
  \item Which word finds where the two sides of an equation are equal? \fromday{19}\writeline
  \item Which word says how many values you need in each group? \fromday{17}\writeline
  \item You made a mistake in the last step. Which word takes it back? \fromday{6}\writeline
\end{enumerate}
{\small Answers: \hyperref[answers:python:22]{page~\pageref*{answers:python:22}}, or \code{sdata.\allowbreak{}answers(22)}.}
\end{warmup}
\begin{todo}
\begin{itemize}
  \item Project day 2: clean. Go column by column through what \code{describe\_\allowbreak{}columns} showed you.
  \item Fix text with spaces or mixed capitals, numbers stored as text, odd codes such as -999, and decide what to do with empty cells (remove, fill, or keep and let the tests skip them).
  \item Save the clean table under a new name, and keep the \code{steps()}.
\end{itemize}
\field{In the manual}{lesson~11, \manualpage{lesson:11}{57}{chapter.II.11}.}
\field{Words to look up}{\kind{table}\,\code{table.\allowbreak{}count\_\allowbreak{}empty\_\allowbreak{}cells()}~\manualpage{word:table.count-empty-cells}{163}{section*.142}; \kind{table}\,\code{table.\allowbreak{}remove\_\allowbreak{}extra\_\allowbreak{}spaces()}~\manualpage{word:table.remove-extra-spaces}{167}{section*.152}; \kind{table}\,\code{table.\allowbreak{}replace\_\allowbreak{}values()}~\manualpage{word:table.replace-values}{167}{section*.153}; \kind{table}\,\code{table.\allowbreak{}fill\_\allowbreak{}empty\_\allowbreak{}cells()}~\manualpage{word:table.fill-empty-cells}{164}{section*.144}.}
\end{todo}
\begin{tryit}
\begin{enumerate}
  \item Clean your table and save it as a new file. Write one line for each decision you made about empty cells.
\end{enumerate}
{\small Worked answers: \hyperref[answers:python:22]{page~\pageref*{answers:python:22}}, or \code{sdata.\allowbreak{}answers(22)}.}
\end{tryit}

\daystart{23}{Project: explore}{Python track}{day:python:23}
\goal{Explore your clean table with summaries and charts before any test. About an hour.}
\begin{warmup}
\begin{enumerate}
  \item Which word fills an even map from scattered points, ready for a contour map? \fromday{20}\writeline
  \item Which word fits a formula of your own, such as growth? \fromday{18}\writeline
  \item Which word gives the middle value of numbers? \fromday{9}\writeline
  \item Which word shows the numbers on the y-axis as percents? \fromday{4}\writeline
\end{enumerate}
{\small Answers: \hyperref[answers:python:23]{page~\pageref*{answers:python:23}}, or \code{sdata.\allowbreak{}answers(23)}.}
\end{warmup}
\begin{todo}
\begin{itemize}
  \item Project day 3: explore. Draw before you test: charts show odd values, skew and outliers that a single number hides.
  \item Count each category, sum up the numbers for each group, look at histograms and at \code{sdata.\allowbreak{}pairs}.
\end{itemize}
\field{In the manual}{lesson~9, \manualpage{lesson:9}{48}{chapter.II.9}; lesson~12, \manualpage{lesson:12}{64}{chapter.II.12}.}
\field{Words to look up}{\kind{table}\,\code{table.\allowbreak{}count\_\allowbreak{}each()}~\manualpage{word:table.count-each}{176}{section*.172}; \kind{table}\,\code{table.\allowbreak{}summary\_\allowbreak{}of()}~\manualpage{word:table.summary-of}{179}{section*.179}; \kind{table}\,\code{table.\allowbreak{}correlation\_\allowbreak{}between()}~\manualpage{word:table.correlation-between}{182}{section*.184}; \kind{chart}\,\code{chart.\allowbreak{}histogram()}~\manualpage{word:histogram}{115}{section*.24}.}
\end{todo}
\begin{tryit}
\begin{enumerate}
  \item Make three summaries and two charts of your table that help to answer your question.
\end{enumerate}
{\small Worked answers: \hyperref[answers:python:23]{page~\pageref*{answers:python:23}}, or \code{sdata.\allowbreak{}answers(23)}.}
\end{tryit}

\daystart{24}{Project: figures}{Python track}{day:python:24}
\goal{Make the main figure of your project, ready for a journal. About an hour.}
\begin{warmup}
\begin{enumerate}
  \item Which word gives the average mark of each class? \fromday{1}\writeline
  \item How do you make a normal distribution with average 50 and standard deviation 10? \fromday{19}\writeline
  \item Which word keeps rows 10 to 20, by their place in the table? \fromday{6}\writeline
  \item How do you make 100 random numbers with average 50 and standard deviation 10, the same every time? \fromday{9}\writeline
\end{enumerate}
{\small Answers: \hyperref[answers:python:24]{page~\pageref*{answers:python:24}}, or \code{sdata.\allowbreak{}answers(24)}.}
\end{warmup}
\begin{todo}
\begin{itemize}
  \item Project day 4: figures. Make the one figure that answers your question best, with two panels (A and B), axis names with units, and a caption that says what the error bars or bands show.
  \item Use style \code{\textquotedbl{}paper\textquotedbl{}}, a size of 17 cm wide, and save it as \code{.\allowbreak{}pdf} and as a 600 dpi \code{.\allowbreak{}tif}.
\end{itemize}
\field{In the manual}{lesson~7, \manualpage{lesson:7}{36}{chapter.II.7}; lesson~8, \manualpage{lesson:8}{41}{chapter.II.8}; lesson~15, \manualpage{lesson:15}{81}{chapter.II.15}.}
\field{Words to look up}{\kind{sdata}\,\code{sdata.\allowbreak{}grid()}~\manualpage{word:sdata.grid}{149}{section*.106}; \kind{chart}\,\code{chart.\allowbreak{}style()}~\manualpage{word:style}{142}{section*.86}; \kind{chart}\,\code{chart.\allowbreak{}size()}~\manualpage{word:size}{142}{section*.87}; \kind{chart}\,\code{chart.\allowbreak{}save()}~\manualpage{word:save}{112}{section*.18}.}
\end{todo}
\begin{tryit}
\begin{enumerate}
  \item Make your main figure with two panels and save it ready for a journal.
\end{enumerate}
{\small Worked answers: \hyperref[answers:python:24]{page~\pageref*{answers:python:24}}, or \code{sdata.\allowbreak{}answers(24)}.}
\end{tryit}

\daystart{25}{Project: tests}{Python track}{day:python:25}
\goal{Choose and run the test that answers your question, with its effect size. About an hour.}
\begin{warmup}
\begin{enumerate}
  \item Which word gives the chance of a value above 80? \fromday{19}\writeline
  \item How do you get the number in row 1, column 2 of a grid? \fromday{9}\writeline
  \item Which word lists every step you did to a table? \fromday{6}\writeline
  \item How do you give each class its own colour in a bar chart? \fromday{2}\writeline
\end{enumerate}
{\small Answers: \hyperref[answers:python:25]{page~\pageref*{answers:python:25}}, or \code{sdata.\allowbreak{}answers(25)}.}
\end{warmup}
\begin{todo}
\begin{itemize}
  \item Project day 5: tests. Choose the test with the four questions of \code{\textquotedbl{}Which test when?\textquotedbl{}} before you run it, and run one test per question.
  \item Check the shape and the outliers first; use \code{based\_\allowbreak{}on=\textquotedbl{}ranks\textquotedbl{}} if the groups are small and far from bell-shaped.
\end{itemize}
\field{In the manual}{lesson~14, \manualpage{lesson:14}{76}{chapter.II.14}; ``Which test when?'', \manualpage{ideas:which-test}{249}{section*.367}.}
\field{Words to look up}{\kind{table}\,\code{table.\allowbreak{}compare\_\allowbreak{}averages()}~\manualpage{word:table.compare-averages}{183}{section*.186}; \kind{table}\,\code{table.\allowbreak{}compare\_\allowbreak{}counts()}~\manualpage{word:table.compare-counts}{185}{section*.190}; \kind{table}\,\code{table.\allowbreak{}correlation\_\allowbreak{}between()}~\manualpage{word:table.correlation-between}{182}{section*.184}; \kind{table}\,\code{table.\allowbreak{}compare\_\allowbreak{}shape\_\allowbreak{}to()}~\manualpage{word:table.compare-shape-to}{185}{section*.191}.}
\end{todo}
\begin{tryit}
\begin{enumerate}
  \item Run the test that answers your question. Write down the test's name, the effect size and the p-value.
\end{enumerate}
{\small Worked answers: \hyperref[answers:python:25]{page~\pageref*{answers:python:25}}, or \code{sdata.\allowbreak{}answers(25)}.}
\end{tryit}

\daystart{26}{Project: models and fits}{Python track}{day:python:26}
\goal{Check your answer with a model of several columns, or a fitted line or formula. About an hour.}
\begin{warmup}
\begin{enumerate}
  \item Which word fits a distribution to a column? \fromday{19}\writeline
  \item Which word draws one dot for each row, x against y? \fromday{3}\writeline
  \item The same children were measured twice. Which word compares the two? \fromday{12}\writeline
  \item Which word draws one small chart for each class? \fromday{5}\writeline
\end{enumerate}
{\small Answers: \hyperref[answers:python:26]{page~\pageref*{answers:python:26}}, or \code{sdata.\allowbreak{}answers(26)}.}
\end{warmup}
\begin{todo}
\begin{itemize}
  \item Project day 6: go further. If several columns could explain the answer, put them in one model with \code{explain()}; if one column follows another, fit a line or a formula.
  \item Compare the model with the simple test of yesterday: does the answer hold when the other columns are taken into account?
\end{itemize}
\field{In the manual}{lesson~16, \manualpage{lesson:16}{86}{chapter.II.16}; lesson~19, \manualpage{lesson:19}{99}{chapter.II.19}.}
\field{Words to look up}{\kind{table}\,\code{table.\allowbreak{}explain()}~\manualpage{word:table.explain}{200}{section*.227}; \kind{table}\,\code{table.\allowbreak{}fit\_\allowbreak{}line()}~\manualpage{word:table.fit-line}{188}{section*.197}; \kind{table}\,\code{table.\allowbreak{}add\_\allowbreak{}predicted\_\allowbreak{}column()}~\manualpage{word:table.add-predicted-column}{200}{section*.228}.}
\end{todo}
\begin{tryit}
\begin{enumerate}
  \item Explain your main column from two or three others, and say in one sentence what changed compared with the single test.
\end{enumerate}
{\small Worked answers: \hyperref[answers:python:26]{page~\pageref*{answers:python:26}}, or \code{sdata.\allowbreak{}answers(26)}.}
\end{tryit}

\daystart{27}{Project: results text}{Python track}{day:python:27}
\goal{Write the results, the methods and the figure caption of your project from the printed answers. About an hour.}
\begin{warmup}
\begin{enumerate}
  \item Which word hands a table on to pandas? \fromday{13}\writeline
  \item Which word explains marks from hours, sleep and breakfast together? \fromday{15}\writeline
  \item Which word tests whether the counts of categories differ between groups? \fromday{11}\writeline
  \item Which word tests whether the average marks differ between the classes? \fromday{10}\writeline
\end{enumerate}
{\small Answers: \hyperref[answers:python:27]{page~\pageref*{answers:python:27}}, or \code{sdata.\allowbreak{}answers(27)}.}
\end{warmup}
\begin{todo}
\begin{itemize}
  \item Project day 7: write. Turn the printed answers into a results paragraph: for each finding, the summary of each group with n, the effect size, the name of the test and the exact p-value.
  \item Write the methods from \code{table.\allowbreak{}steps()}: what you removed or filled, and which tests you used.
  \item Write the caption of your figure: what each panel shows, what the bars or bands are, and n.
\end{itemize}
\field{In the manual}{``Writing it up'', \manualpage{ideas:report}{250}{section*.368}.}
\field{Words to look up}{\kind{table}\,\code{table.\allowbreak{}steps()}~\manualpage{word:table.steps}{211}{section*.251}.}
\end{todo}
\begin{tryit}
\begin{enumerate}
  \item Write a results paragraph, a methods paragraph and a figure caption for your project.
\end{enumerate}
{\small Worked answers: \hyperref[answers:python:27]{page~\pageref*{answers:python:27}}, or \code{sdata.\allowbreak{}answers(27)}.}
\end{tryit}

\daystart{28}{Review, and finish your project}{Python track}{day:python:28}
\goal{Remember words from the whole course, and run your project as one script from the raw file to the figure. About an hour.}
\begin{warmup}
\begin{enumerate}
  \item Which word counts the empty cells in each column? \fromday{8}\writeline
  \item Which word draws survival curves? \fromday{17}\writeline
  \item Which word adds the principal components PC1 and PC2? \fromday{16}\writeline
  \item Which word estimates a column between the points you measured? \fromday{20}\writeline
  \item Which word fits a straight line on a table? \fromday{18}\writeline
  \item Which option makes a bar show the median instead of the average? \fromday{4}\writeline
  \item Which word lists the practice tables? \fromday{1}\writeline
  \item Which word gives the average of a column for each group? \fromday{9}\writeline
  \item How do you look at the first 5 rows of a table? \fromday{6}\writeline
  \item Which word moves the legend below the chart? \fromday{2}\writeline
\end{enumerate}
{\small Answers: \hyperref[answers:python:28]{page~\pageref*{answers:python:28}}, or \code{sdata.\allowbreak{}answers(28)}.}
\end{warmup}
\begin{todo}
\begin{itemize}
  \item Answer the warm-up on paper, from memory, before you look anything up.
  \item For each question you could not answer, open its day with \code{sdata.\allowbreak{}day(n)} and run its example again.
  \item Project day 8: finish. Put every step of your project in one script, from the raw file to the figure, and run it from the top in a fresh session. It must run without errors.
  \item Read \code{steps()} out loud: would someone else understand what you did?
\end{itemize}
\field{Words to look up}{\kind{table}\,\code{table.\allowbreak{}steps()}~\manualpage{word:table.steps}{211}{section*.251}; \kind{table}\,\code{table.\allowbreak{}save()}~\manualpage{word:table.save}{212}{section*.253}.}
\end{todo}
\begin{tryit}
\begin{enumerate}
  \item Without looking back: which word fits a formula, which finds where it is biggest, and which gives the chance of a value above a number? Use one of them on any table or formula.
  \item Run your whole project script from the top. Check that it makes the clean table, the figure and the numbers in your text.
\end{enumerate}
{\small Worked answers: \hyperref[answers:python:28]{page~\pageref*{answers:python:28}}, or \code{sdata.\allowbreak{}answers(28)}.}
\end{tryit}

\daystart{29}{Final test from memory}{Python track}{day:python:29}
\goal{Answer questions on the whole course and write a whole analysis from memory. About an hour.}
\begin{warmup}
\begin{enumerate}
  \item Which word finds where a formula is smallest? \fromday{19}\writeline
  \item Which word compares one average with a number you choose? \fromday{12}\writeline
  \item Which word puts two charts next to each other on one page? \fromday{5}\writeline
  \item Which word draws a line with a dot at each point? \fromday{3}\writeline
  \item How do you explain a yes-or-no column such as \code{\textquotedbl{}passed\textquotedbl{}}? \fromday{15}\writeline
  \item Which word lists every change you made to a table, for your methods? \fromday{13}\writeline
  \item Which word shows the rows with outliers? \fromday{11}\writeline
  \item Which word says which pairs of groups differ? \fromday{10}\writeline
  \item Which word fills an even map from scattered points, ready for a contour map? \fromday{20}\writeline
  \item How do you fit a curve with x and x squared? \fromday{18}\writeline
  \item Which word tests whether survival differs between groups? \fromday{17}\writeline
  \item Which word puts rows that are alike into groups? \fromday{16}\writeline
  \item Which word tells whether two columns go up and down together? \fromday{9}\writeline
  \item How do you fill the empty marks with the median? \fromday{8}\writeline
\end{enumerate}
{\small Answers: \hyperref[answers:python:29]{page~\pageref*{answers:python:29}}, or \code{sdata.\allowbreak{}answers(29)}.}
\end{warmup}
\begin{todo}
\begin{itemize}
  \item Today's warm-up is the final test: answer every question on paper, from memory. Then check your answers with \code{sdata.\allowbreak{}answers()}.
  \item Then write the whole program of the task in an empty file, without looking anything up. Only when it is finished, check it with \code{sdata.\allowbreak{}help}.
\end{itemize}
\field{Words to look up}{\kind{sdata}\,\code{sdata.\allowbreak{}words()}~\manualpage{word:sdata.words}{238}{section*.332}; \kind{sdata}\,\code{sdata.\allowbreak{}help()}~\manualpage{word:sdata.help}{238}{section*.333}.}
\end{todo}
\begin{tryit}
\begin{enumerate}
  \item From memory, in an empty file: open the \code{\textquotedbl{}marks\textquotedbl{}} table, remove the rows with empty cells, give the average marks for each class, test the difference, and save a two-panel figure (boxes, and dots of hours against marks with a fitted line) as a PDF.
\end{enumerate}
{\small Worked answers: \hyperref[answers:python:29]{page~\pageref*{answers:python:29}}, or \code{sdata.\allowbreak{}answers(29)}.}
\end{tryit}

\daystart{30}{The same work in pandas and matplotlib}{Python track}{day:python:30}
\goal{See what simplydata does underneath by writing the same analysis in pandas, SciPy and matplotlib. About an hour.}
\begin{warmup}
\begin{enumerate}
  \item Which word estimates a column between the points you measured? \fromday{20}\writeline
  \item Which word finds where the two sides of an equation are equal? \fromday{19}\writeline
  \item Which word fits a formula of your own, such as growth? \fromday{18}\writeline
  \item Which word tests whether groups differ in their make-up across many columns (PERMANOVA)? \fromday{16}\writeline
\end{enumerate}
{\small Answers: \hyperref[answers:python:30]{page~\pageref*{answers:python:30}}, or \code{sdata.\allowbreak{}answers(30)}.}
\end{warmup}
\begin{todo}
\begin{itemize}
  \item Today: the same work in pandas and matplotlib, the libraries simplydata uses underneath. Each word of the manual's dictionary says which one does the work (\code{\textquotedbl{}Underneath\textquotedbl{}}).
  \item Run the simplydata version and the pandas version of the task. Compare them line by line.
  \item Look back at the calendar with \code{sdata.\allowbreak{}day()} and choose three days to do again next month.
\end{itemize}
\field{In the manual}{lesson~15, \manualpage{lesson:15}{81}{chapter.II.15}.}
\field{Words to look up}{\kind{table}\,\code{table.\allowbreak{}as\_\allowbreak{}pandas()}~\manualpage{word:table.as-pandas}{212}{section*.255}; \kind{table}\,\code{table.\allowbreak{}as\_\allowbreak{}polars()}~\manualpage{word:table.as-polars}{213}{section*.256}; \kind{table}\,\code{table.\allowbreak{}steps()}~\manualpage{word:table.steps}{211}{section*.251}.}
\end{todo}
\begin{tryit}
\begin{enumerate}
  \item In simplydata: keep classes A and B of the marks table, give the average marks of each class, test the difference and draw it as bars.
  \item Write the same analysis with pandas, SciPy and matplotlib, then compare the two line by line.
\end{enumerate}
{\small Worked answers: \hyperref[answers:python:30]{page~\pageref*{answers:python:30}}, or \code{sdata.\allowbreak{}answers(30)}.}
\end{tryit}

\trackstart{Beginner track}{For people who have never programmed}{track:beginner}
Three days of just enough Python first, then about one lesson a day, with review days on days 7, 14 and 21, a small project of your own (days 27 and 28), a final test from memory and a look back. Tick each day when you have done it; the day's number is a link to its page.
\begin{center}\small
\begin{longtable}{@{}c r >{\raggedright\arraybackslash}p{0.56\linewidth} >{\raggedright\arraybackslash}p{0.14\linewidth} r@{}}
\toprule & \textbf{Day} & \textbf{What you do} & \textbf{Lesson} & \textbf{Page}\\\midrule\endhead
\tick & \hyperref[day:beginner:1]{1} & Python first steps: values and names & - & \pageref*{day:beginner:1}\\[0.25em]
\tick & \hyperref[day:beginner:2]{2} & Python first steps: lists and words & - & \pageref*{day:beginner:2}\\[0.25em]
\tick & \hyperref[day:beginner:3]{3} & Reading an error, and your first chart & 1 & \pageref*{day:beginner:3}\\[0.25em]
\tick & \hyperref[day:beginner:4]{4} & Titles and labels & 2 & \pageref*{day:beginner:4}\\[0.25em]
\tick & \hyperref[day:beginner:5]{5} & Colours, groups and legends & 3 & \pageref*{day:beginner:5}\\[0.25em]
\tick & \hyperref[day:beginner:6]{6} & Different chart types & 4 & \pageref*{day:beginner:6}\\[0.25em]
\tick & \hyperref[day:beginner:7]{7} & Review of the first week & review & \pageref*{day:beginner:7}\\[0.25em]
\tick & \hyperref[day:beginner:8]{8} & Averages, error bars and fit lines & 5 & \pageref*{day:beginner:8}\\[0.25em]
\tick & \hyperref[day:beginner:9]{9} & Axes and reference lines & 6 & \pageref*{day:beginner:9}\\[0.25em]
\tick & \hyperref[day:beginner:10]{10} & Several charts together & 7 & \pageref*{day:beginner:10}\\[0.25em]
\tick & \hyperref[day:beginner:11]{11} & A figure for a journal, and word families & 8 & \pageref*{day:beginner:11}\\[0.25em]
\tick & \hyperref[day:beginner:12]{12} & Open a table and look at it & 9 & \pageref*{day:beginner:12}\\[0.25em]
\tick & \hyperref[day:beginner:13]{13} & Choose rows and columns, and undo & 10 & \pageref*{day:beginner:13}\\[0.25em]
\tick & \hyperref[day:beginner:14]{14} & Review of the second week & review & \pageref*{day:beginner:14}\\[0.25em]
\tick & \hyperref[day:beginner:15]{15} & Clean the data and add columns & 11 & \pageref*{day:beginner:15}\\[0.25em]
\tick & \hyperref[day:beginner:16]{16} & Summaries, and from a table to a chart & 12 & \pageref*{day:beginner:16}\\[0.25em]
\tick & \hyperref[day:beginner:17]{17} & Lists and grids of numbers & 13 & \pageref*{day:beginner:17}\\[0.25em]
\tick & \hyperref[day:beginner:18]{18} & Is a difference real? Tests and effect sizes & 14 & \pageref*{day:beginner:18}\\[0.25em]
\tick & \hyperref[day:beginner:19]{19} & Counts, outliers and many tests & 14 & \pageref*{day:beginner:19}\\[0.25em]
\tick & \hyperref[day:beginner:20]{20} & Choosing a test, and a whole analysis & 15 & \pageref*{day:beginner:20}\\[0.25em]
\tick & \hyperref[day:beginner:21]{21} & Review of the third week & review & \pageref*{day:beginner:21}\\[0.25em]
\tick & \hyperref[day:beginner:22]{22} & Models: explain one column from several & 16 & \pageref*{day:beginner:22}\\[0.25em]
\tick & \hyperref[day:beginner:23]{23} & Patterns across many columns & 17 & \pageref*{day:beginner:23}\\[0.25em]
\tick & \hyperref[day:beginner:24]{24} & Survival, and planning a study & 18 & \pageref*{day:beginner:24}\\[0.25em]
\tick & \hyperref[day:beginner:25]{25} & Fitting, finding and solving & 19 & \pageref*{day:beginner:25}\\[0.25em]
\tick & \hyperref[day:beginner:26]{26} & Distributions and values between points & 20 & \pageref*{day:beginner:26}\\[0.25em]
\tick & \hyperref[day:beginner:27]{27} & Small project: look, clean and sum up & project & \pageref*{day:beginner:27}\\[0.25em]
\tick & \hyperref[day:beginner:28]{28} & Small project: test, figure and results & project & \pageref*{day:beginner:28}\\[0.25em]
\tick & \hyperref[day:beginner:29]{29} & Final test from memory & final test & \pageref*{day:beginner:29}\\[0.25em]
\tick & \hyperref[day:beginner:30]{30} & Look back, and what comes next & - & \pageref*{day:beginner:30}\\[0.25em]
\bottomrule\end{longtable}\end{center}

\daystart{1}{Python first steps: values and names}{Beginner track}{day:beginner:1}
\goal{Run your first lines of Python: text in quotes, numbers, names for values, and print. About an hour.}
\begin{warmup}
Nothing to remember yet on day 1: just start.
\end{warmup}
\begin{todo}
\begin{itemize}
  \item Open a place to type Python: a Jupyter notebook, VS Code, Spyder, or IDLE (it comes with Python). The manual's \code{\textquotedbl{}Install and first steps\textquotedbl{}} says how.
  \item Type each line of the example below, run it, and read what it prints.
  \item Then save the same lines in a file called day1.py and run the whole file: in a terminal, type python day1.py. A file of lines like this is called a script.
\end{itemize}
\field{In the manual}{``Install and first steps'', \manualpage{start:install}{4}{section*.5}.}
\par\noindent{\sffamily\small\bfseries\color{brand}Example: type it yourself}\par
\begin{lstlisting}[style=py]
name = "Asha"
age = 12
print(name)
print(age + 1)
print("Hello", name)
height_m = 1.44
print(height_m * 100)
\end{lstlisting}
\whateachline
\begin{lineexplain}
  \item[1] Give the text \code{\textquotedbl{}Asha\textquotedbl{}} the name name. Text always goes in quotes.
  \item[2] Give the number 12 the name age. Numbers have no quotes.
  \item[3] print shows a value on the screen: Asha.
  \item[4] Maths works as on a calculator: + - * / and ** for a power. This prints 13.
  \item[5] print can show several values, separated by commas.
  \item[6] A name has no spaces (use \_ instead), and a decimal number has a point.
  \item[7] Prints 144.0: the height in centimetres.
\end{lineexplain}
\end{todo}
\begin{tryit}
\begin{enumerate}
  \item Give the text \code{\textquotedbl{}Bergen\textquotedbl{}} the name city and the number 2250 the name rain (the rain in a year, in mm). Print the words \code{\textquotedbl{}Rain in\textquotedbl{}}, the city, and the rain of an average month (rain / 12).
\end{enumerate}
{\small Worked answers: \hyperref[answers:beginner:1]{page~\pageref*{answers:beginner:1}}, or \code{sdata.\allowbreak{}answers(1)}.}
\end{tryit}

\daystart{2}{Python first steps: lists and words}{Beginner track}{day:beginner:2}
\goal{Make lists, call a word with brackets and options, and use simplydata for the first time. About an hour.}
\begin{warmup}
\begin{enumerate}
  \item How do you write the text Asha in Python? \fromday{1}\writeline
  \item How do you give the number 12 the name age? \fromday{1}\writeline
  \item What does \code{print(age)} do? \fromday{1}\writeline
\end{enumerate}
{\small Answers: \hyperref[answers:beginner:2]{page~\pageref*{answers:beginner:2}}, or \code{sdata.\allowbreak{}answers(2)}.}
\end{warmup}
\begin{todo}
\begin{itemize}
  \item Type and run the example below, one line at a time.
  \item Look closely at the brackets: square brackets [ ] make a list, round brackets ( ) call a word, and curly brackets \{ \} make a table by hand.
  \item Then type \code{sdata.\allowbreak{}help(\textquotedbl{}average\textquotedbl{})} and read what it says.
\end{itemize}
\field{In the manual}{``Install and first steps'', \manualpage{start:install}{4}{section*.5}.}
\field{Words to look up}{\kind{sdata}\,\code{sdata.\allowbreak{}numbers()}~\manualpage{word:sdata.numbers}{214}{section*.260}; \kind{nums}\,\code{nums.\allowbreak{}average()}~\manualpage{word:nums.average}{226}{section*.295}; \kind{nums}\,\code{nums.\allowbreak{}keep\_\allowbreak{}numbers()}~\manualpage{word:nums.keep-numbers}{221}{section*.281}.}
\par\noindent{\sffamily\small\bfseries\color{brand}Example: type it yourself}\par
\begin{lstlisting}[style=py]
marks = [48, 71, 55, 88, 62]
print(len(marks))
print(round(3.14159, 2))
import simplydata as sdata
nums = sdata.numbers(marks)
nums.average()
nums.keep_numbers(above=50)
nums.show()
kids = {"name": ["Asha", "Ben"], "height": [144, 152]}
table = sdata.table(kids)
table.show()
\end{lstlisting}
\whateachline
\begin{lineexplain}
  \item[1] A list: values in square brackets, with commas between.
  \item[2] len is a word of Python that counts: 5. A word is called with round brackets around what it works on.
  \item[3] Two things in the brackets, separated by a comma: round 3.14159 to 2 decimals.
  \item[4] Load simplydata. From now on we call it by its short name, sdata.
  \item[5] numbers is a word of sdata: the point joins a thing and its word. We call the result nums.
  \item[6] A step on the numbers: it prints the answer in a sentence. Empty brackets still have to be there.
  \item[7] \code{above=50} is an option with a name: the name, =, the value.
  \item[8] Show the numbers that are left.
  \item[9] A table typed by hand: each column name with its list of values.
  \item[10] Turn it into a simplydata table.
  \item[11] Show the table.
\end{lineexplain}
\end{todo}
\begin{tryit}
\begin{enumerate}
  \item Make a list of five temperatures, turn it into simplydata numbers, and find the biggest one and the median.
\end{enumerate}
{\small Worked answers: \hyperref[answers:beginner:2]{page~\pageref*{answers:beginner:2}}, or \code{sdata.\allowbreak{}answers(2)}.}
\end{tryit}

\daystart{3}{Reading an error, and your first chart}{Beginner track}{day:beginner:3}
\goal{Read an error message calmly from the bottom up, then make and save your first chart. About an hour.}
\begin{warmup}
\begin{enumerate}
  \item How do you write a list of the numbers 3, 5 and 2? \fromday{2}\writeline
  \item In \code{nums.\allowbreak{}keep\_\allowbreak{}numbers(above=50)}, what is \code{above=50}? \fromday{2}\writeline
  \item How do you write the text Asha in Python? \fromday{1}\writeline
\end{enumerate}
{\small Answers: \hyperref[answers:beginner:3]{page~\pageref*{answers:beginner:3}}, or \code{sdata.\allowbreak{}answers(3)}.}
\end{warmup}
\begin{todo}
\begin{itemize}
  \item Errors are normal: everybody gets them, every day. Today you make some on purpose.
  \item Type each of the errors to make on purpose, below, and read the error from the bottom up: the last line says what went wrong.
  \item Then start lesson 1, below.
  \item Work through \code{sdata.\allowbreak{}lesson(1)}: type each example yourself (do not copy and paste), run it and read what each line does. Then do the lesson's tasks and check them with \code{sdata.\allowbreak{}lesson(1,\allowbreak{} solutions=True)}.
\end{itemize}
\field{In the manual}{lesson~1, \manualpage{lesson:1}{10}{chapter.II.1}; ``Install and first steps'', \manualpage{start:install}{4}{section*.5}.}
\field{Words to look up}{\kind{sdata}\,\code{sdata.\allowbreak{}help()}~\manualpage{word:sdata.help}{238}{section*.333}; \kind{sdata}\,\code{sdata.\allowbreak{}example()}~\manualpage{word:sdata.example}{111}{section*.16}; \kind{sdata}\,\code{sdata.\allowbreak{}chart()}~\manualpage{word:sdata.chart}{111}{section*.15}; \kind{chart}\,\code{chart.\allowbreak{}bars()}~\manualpage{word:bars}{113}{section*.20}; \kind{chart}\,\code{chart.\allowbreak{}save()}~\manualpage{word:save}{112}{section*.18}.}
\par\noindent{\sffamily\small\bfseries\color{brand}Errors to make on purpose} {\small(read the last line of each error first)}\par
\begin{lstlisting}[style=py]
print(nmae)
\end{lstlisting}
\vspace{-0.3em}{\small NameError: a name is misspelt (nmae, not name), or was never given a value with = first.}\par
\begin{lstlisting}[style=py]
print("Asha)
\end{lstlisting}
\vspace{-0.3em}{\small SyntaxError: a quote mark or a bracket is not closed. Every opening quote mark needs a closing one, and so does every bracket.}\par
\begin{lstlisting}[style=py]
sdata.example("wether")
\end{lstlisting}
\vspace{-0.3em}{\small simplydata's own errors say what went wrong and suggest the word you probably meant.}\par
\begin{lstlisting}[style=py]
nums = sdata.numbers([3, 5, 2])
nums.avrage()
\end{lstlisting}
\vspace{-0.3em}{\small A misspelt step: simplydata lists the steps you can use and the nearest one.}\par
\end{todo}
\begin{tryit}
\begin{enumerate}
  \item Find and fix the two mistakes in these lines, then run them: nums = \code{sdata.\allowbreak{}numbers([4,\allowbreak{} 8,\allowbreak{} 6)} and then \code{nums.\allowbreak{}averag()}
  \item Make a bar chart of the \code{\textquotedbl{}weather\textquotedbl{}} table with \code{x=\textquotedbl{}city\textquotedbl{}} and \code{y=\textquotedbl{}rain\textquotedbl{}}, and save it as \code{\textquotedbl{}rain.\allowbreak{}png\textquotedbl{}}.
\end{enumerate}
{\small Worked answers: \hyperref[answers:beginner:3]{page~\pageref*{answers:beginner:3}}, or \code{sdata.\allowbreak{}answers(3)}.}
\end{tryit}

\daystart{4}{Titles and labels}{Beginner track}{day:beginner:4}
\goal{Put a title, axis names and a caption on a chart, and make its text easy to read. About an hour.}
\begin{warmup}
\begin{enumerate}
  \item Where in an error message do you look first? \fromday{3}\writeline
  \item How do you give the number 12 the name age? \fromday{1}\writeline
  \item What does NameError: name 'nmae' is not defined usually mean? \fromday{3}\writeline
\end{enumerate}
{\small Answers: \hyperref[answers:beginner:4]{page~\pageref*{answers:beginner:4}}, or \code{sdata.\allowbreak{}answers(4)}.}
\end{warmup}
\begin{todo}
\begin{itemize}
  \item Work through \code{sdata.\allowbreak{}lesson(2)}: type each example yourself (do not copy and paste), run it and read what each line does. Then do the lesson's tasks and check them with \code{sdata.\allowbreak{}lesson(2,\allowbreak{} solutions=True)}.
\end{itemize}
\field{In the manual}{lesson~2, \manualpage{lesson:2}{13}{chapter.II.2}.}
\field{Words to look up}{\kind{chart}\,\code{chart.\allowbreak{}title\_\allowbreak{}top()}~\manualpage{word:title-top}{131}{section*.59}; \kind{chart}\,\code{chart.\allowbreak{}label\_\allowbreak{}x()}~\manualpage{word:label-x}{132}{section*.61}; \kind{chart}\,\code{chart.\allowbreak{}label\_\allowbreak{}y()}~\manualpage{word:label-y}{132}{section*.62}; \kind{chart}\,\code{chart.\allowbreak{}rotate\_\allowbreak{}labels\_\allowbreak{}x()}~\manualpage{word:rotate-labels-x}{138}{section*.75}.}
\end{todo}
\begin{tryit}
\begin{enumerate}
  \item Draw dots of \code{\textquotedbl{}hours\textquotedbl{}} (x) against \code{\textquotedbl{}marks\textquotedbl{}} (y) from the marks table, with a title, a name for each axis and a caption at the bottom.
\end{enumerate}
{\small Worked answers: \hyperref[answers:beginner:4]{page~\pageref*{answers:beginner:4}}, or \code{sdata.\allowbreak{}answers(4)}.}
\end{tryit}

\daystart{5}{Colours, groups and legends}{Beginner track}{day:beginner:5}
\goal{Give each group its own colour and place the legend where it helps. About an hour.}
\begin{warmup}
\begin{enumerate}
  \item Which word puts a title above the chart? \fromday{4}\writeline
  \item What does every simplydata program start with? \fromday{2}\writeline
  \item How do you name the y-axis? \fromday{4}\writeline
\end{enumerate}
{\small Answers: \hyperref[answers:beginner:5]{page~\pageref*{answers:beginner:5}}, or \code{sdata.\allowbreak{}answers(5)}.}
\end{warmup}
\begin{todo}
\begin{itemize}
  \item Work through \code{sdata.\allowbreak{}lesson(3)}: type each example yourself (do not copy and paste), run it and read what each line does. Then do the lesson's tasks and check them with \code{sdata.\allowbreak{}lesson(3,\allowbreak{} solutions=True)}.
\end{itemize}
\field{In the manual}{lesson~3, \manualpage{lesson:3}{16}{chapter.II.3}.}
\field{Words to look up}{\kind{chart}\,\code{chart.\allowbreak{}colour()}~\manualpage{word:colour}{140}{section*.83}; \kind{chart}\,\code{chart.\allowbreak{}colours()}~\manualpage{word:colours}{141}{section*.84}; \kind{chart}\,\code{chart.\allowbreak{}legend\_\allowbreak{}bottom()}~\manualpage{word:legend-bottom}{144}{section*.91}; \kind{chart}\,\code{chart.\allowbreak{}order\_\allowbreak{}x()}~\manualpage{word:order-x}{146}{section*.99}.}
\end{todo}
\begin{tryit}
\begin{enumerate}
  \item Weather table: draw lines of \code{\textquotedbl{}temperature\textquotedbl{}} over \code{\textquotedbl{}month\textquotedbl{}}, one colour for each city. Put the legend on top and use the style \code{\textquotedbl{}paper\textquotedbl{}}.
\end{enumerate}
{\small Worked answers: \hyperref[answers:beginner:5]{page~\pageref*{answers:beginner:5}}, or \code{sdata.\allowbreak{}answers(5)}.}
\end{tryit}

\daystart{6}{Different chart types}{Beginner track}{day:beginner:6}
\goal{Draw lines, dots, histograms, boxes, violins, pies and areas. About an hour.}
\begin{warmup}
\begin{enumerate}
  \item How do you give each class its own colour in a bar chart? \fromday{5}\writeline
  \item Python says SyntaxError. What do you check first? \fromday{3}\writeline
  \item Which word moves the legend below the chart? \fromday{5}\writeline
\end{enumerate}
{\small Answers: \hyperref[answers:beginner:6]{page~\pageref*{answers:beginner:6}}, or \code{sdata.\allowbreak{}answers(6)}.}
\end{warmup}
\begin{todo}
\begin{itemize}
  \item Work through \code{sdata.\allowbreak{}lesson(4)}: type each example yourself (do not copy and paste), run it and read what each line does. Then do the lesson's tasks and check them with \code{sdata.\allowbreak{}lesson(4,\allowbreak{} solutions=True)}.
\end{itemize}
\field{In the manual}{lesson~4, \manualpage{lesson:4}{21}{chapter.II.4}.}
\field{Words to look up}{\kind{chart}\,\code{chart.\allowbreak{}lines()}~\manualpage{word:lines}{113}{section*.21}; \kind{chart}\,\code{chart.\allowbreak{}dots()}~\manualpage{word:dots}{114}{section*.22}; \kind{chart}\,\code{chart.\allowbreak{}histogram()}~\manualpage{word:histogram}{115}{section*.24}; \kind{chart}\,\code{chart.\allowbreak{}boxes()}~\manualpage{word:boxes}{116}{section*.25}.}
\end{todo}
\begin{tryit}
\begin{enumerate}
  \item Draw violins of \code{\textquotedbl{}temperature\textquotedbl{}} for each city in the weather table, with all the dots on top.
\end{enumerate}
{\small Worked answers: \hyperref[answers:beginner:6]{page~\pageref*{answers:beginner:6}}, or \code{sdata.\allowbreak{}answers(6)}.}
\end{tryit}

\daystart{7}{Review of the first week}{Beginner track}{day:beginner:7}
\goal{Remember the Python and chart words of days 1 to 6 without looking, and use them together. About an hour.}
\begin{warmup}
\begin{enumerate}
  \item Which chart shows how the values of one column are spread, in bins? \fromday{6}\writeline
  \item The names under the bars bump into each other. Which word helps? \fromday{4}\writeline
  \item Which word gives the average of numbers called nums? \fromday{2}\writeline
  \item Which word lets you choose the colours of the groups? \fromday{5}\writeline
  \item Which word lists the practice tables? \fromday{3}\writeline
  \item What does \code{print(age)} do? \fromday{1}\writeline
\end{enumerate}
{\small Answers: \hyperref[answers:beginner:7]{page~\pageref*{answers:beginner:7}}, or \code{sdata.\allowbreak{}answers(7)}.}
\end{warmup}
\begin{todo}
\begin{itemize}
  \item Answer the warm-up on paper, from memory, before you look anything up.
  \item For each question you could not answer, open its day with \code{sdata.\allowbreak{}day(n)} and run its example again.
  \item Then do the mixed task without looking back.
\end{itemize}
\end{todo}
\begin{tryit}
\begin{enumerate}
  \item Make numbers of five marks and give their average. Then draw bars of the marks table with \code{x=\textquotedbl{}subject\textquotedbl{}}, \code{y=\textquotedbl{}marks\textquotedbl{}}, one colour for each class, and a title. Save it.
\end{enumerate}
{\small Worked answers: \hyperref[answers:beginner:7]{page~\pageref*{answers:beginner:7}}, or \code{sdata.\allowbreak{}answers(7)}.}
\end{tryit}

\daystart{8}{Averages, error bars and fit lines}{Beginner track}{day:beginner:8}
\goal{Show averages and medians with error bars, and draw a line through dots. About an hour.}
\begin{warmup}
\begin{enumerate}
  \item Which two words show the spread of each group, with every value as a dot on top? \fromday{6}\writeline
  \item How do you put the bars in your own order? \fromday{5}\writeline
  \item How do you write the text Asha in Python? \fromday{1}\writeline
\end{enumerate}
{\small Answers: \hyperref[answers:beginner:8]{page~\pageref*{answers:beginner:8}}, or \code{sdata.\allowbreak{}answers(8)}.}
\end{warmup}
\begin{todo}
\begin{itemize}
  \item Work through \code{sdata.\allowbreak{}lesson(5)}: type each example yourself (do not copy and paste), run it and read what each line does. Then do the lesson's tasks and check them with \code{sdata.\allowbreak{}lesson(5,\allowbreak{} solutions=True)}.
\end{itemize}
\field{In the manual}{lesson~5, \manualpage{lesson:5}{26}{chapter.II.5}; ``Standard deviation or standard error?'', \manualpage{ideas:sd-se}{244}{section*.348}.}
\field{Words to look up}{\kind{option}\,\code{show=}~\manualpage{word:show-eq}{121}{section*.36}; \kind{option}\,\code{error=}~\manualpage{word:error-eq}{121}{section*.37}; \kind{chart}\,\code{chart.\allowbreak{}numbers\_\allowbreak{}on\_\allowbreak{}bars()}~\manualpage{word:numbers-on-bars}{133}{section*.64}; \kind{chart}\,\code{chart.\allowbreak{}fit\_\allowbreak{}line()}~\manualpage{word:fit-line}{122}{section*.39}.}
\end{todo}
\begin{tryit}
\begin{enumerate}
  \item Weather table: draw bars of the average \code{\textquotedbl{}rain\textquotedbl{}} for each city with 95 \% confidence intervals, and write the numbers on the bars.
\end{enumerate}
{\small Worked answers: \hyperref[answers:beginner:8]{page~\pageref*{answers:beginner:8}}, or \code{sdata.\allowbreak{}answers(8)}.}
\end{tryit}

\daystart{9}{Axes and reference lines}{Beginner track}{day:beginner:9}
\goal{Choose the range of an axis, and add lines and shading at the values that matter. About an hour.}
\begin{warmup}
\begin{enumerate}
  \item Which option makes a bar show the median instead of the average? \fromday{8}\writeline
  \item Which word draws one dot for each row, x against y? \fromday{6}\writeline
  \item How do you write a list of the numbers 3, 5 and 2? \fromday{2}\writeline
\end{enumerate}
{\small Answers: \hyperref[answers:beginner:9]{page~\pageref*{answers:beginner:9}}, or \code{sdata.\allowbreak{}answers(9)}.}
\end{warmup}
\begin{todo}
\begin{itemize}
  \item Work through \code{sdata.\allowbreak{}lesson(6)}: type each example yourself (do not copy and paste), run it and read what each line does. Then do the lesson's tasks and check them with \code{sdata.\allowbreak{}lesson(6,\allowbreak{} solutions=True)}.
\end{itemize}
\field{In the manual}{lesson~6, \manualpage{lesson:6}{31}{chapter.II.6}.}
\field{Words to look up}{\kind{chart}\,\code{chart.\allowbreak{}range\_\allowbreak{}y()}~\manualpage{word:range-y}{136}{section*.70}; \kind{chart}\,\code{chart.\allowbreak{}line\_\allowbreak{}at\_\allowbreak{}y()}~\manualpage{word:line-at-y}{130}{section*.54}; \kind{chart}\,\code{chart.\allowbreak{}shade\_\allowbreak{}between\_\allowbreak{}y()}~\manualpage{word:shade-between-y}{131}{section*.57}; \kind{chart}\,\code{chart.\allowbreak{}log\_\allowbreak{}scale\_\allowbreak{}y()}~\manualpage{word:log-scale-y}{136}{section*.72}.}
\end{todo}
\begin{tryit}
\begin{enumerate}
  \item Growth table: draw lines of \code{\textquotedbl{}population\textquotedbl{}} over \code{\textquotedbl{}year\textquotedbl{}}, one for each country, with a line at x = 2010 and the y-axis starting at 0.
\end{enumerate}
{\small Worked answers: \hyperref[answers:beginner:9]{page~\pageref*{answers:beginner:9}}, or \code{sdata.\allowbreak{}answers(9)}.}
\end{tryit}

\daystart{10}{Several charts together}{Beginner track}{day:beginner:10}
\goal{Put several charts on one page: one chart per group, side by side, or in a grid. About an hour.}
\begin{warmup}
\begin{enumerate}
  \item Which word sets the y-axis from 0 to 100? \fromday{9}\writeline
  \item Which word draws a line with a dot at each point? \fromday{6}\writeline
  \item How do you start a chart of a table called marks? \fromday{3}\writeline
\end{enumerate}
{\small Answers: \hyperref[answers:beginner:10]{page~\pageref*{answers:beginner:10}}, or \code{sdata.\allowbreak{}answers(10)}.}
\end{warmup}
\begin{todo}
\begin{itemize}
  \item Work through \code{sdata.\allowbreak{}lesson(7)}: type each example yourself (do not copy and paste), run it and read what each line does. Then do the lesson's tasks and check them with \code{sdata.\allowbreak{}lesson(7,\allowbreak{} solutions=True)}.
\end{itemize}
\field{In the manual}{lesson~7, \manualpage{lesson:7}{36}{chapter.II.7}.}
\field{Words to look up}{\kind{chart}\,\code{chart.\allowbreak{}one\_\allowbreak{}chart\_\allowbreak{}per()}~\manualpage{word:one-chart-per}{148}{section*.103}; \kind{sdata}\,\code{sdata.\allowbreak{}side\_\allowbreak{}by\_\allowbreak{}side()}~\manualpage{word:sdata.side-by-side}{148}{section*.104}; \kind{sdata}\,\code{sdata.\allowbreak{}one\_\allowbreak{}above\_\allowbreak{}other()}~\manualpage{word:sdata.one-above-other}{149}{section*.105}; \kind{sdata}\,\code{sdata.\allowbreak{}grid()}~\manualpage{word:sdata.grid}{149}{section*.106}.}
\end{todo}
\begin{tryit}
\begin{enumerate}
  \item Make a histogram and a box chart of the marks, and put them one above the other on one page.
\end{enumerate}
{\small Worked answers: \hyperref[answers:beginner:10]{page~\pageref*{answers:beginner:10}}, or \code{sdata.\allowbreak{}answers(10)}.}
\end{tryit}

\daystart{11}{A figure for a journal, and word families}{Beginner track}{day:beginner:11}
\goal{Size and save a figure for a journal, and learn the patterns that make words easy to guess. About an hour.}
\begin{warmup}
\begin{enumerate}
  \item Which word draws one small chart for each class? \fromday{10}\writeline
  \item Which option adds error bars that say how sure we are of the average? \fromday{8}\writeline
  \item Which word makes all the text on a chart bigger or smaller? \fromday{4}\writeline
\end{enumerate}
{\small Answers: \hyperref[answers:beginner:11]{page~\pageref*{answers:beginner:11}}, or \code{sdata.\allowbreak{}answers(11)}.}
\end{warmup}
\begin{todo}
\begin{itemize}
  \item Work through \code{sdata.\allowbreak{}lesson(8)}: type each example yourself (do not copy and paste), run it and read what each line does. Then do the lesson's tasks and check them with \code{sdata.\allowbreak{}lesson(8,\allowbreak{} solutions=True)}.
  \item Read \code{\textquotedbl{}How the words are built\textquotedbl{}} in the manual (two pages). It is the key to guessing words you have never seen.
  \item Type \code{sdata.\allowbreak{}words()} and look for the families: \code{keep\_\allowbreak{}} and \code{remove\_\allowbreak{}}, ...\_x and ...\_y, \code{add\_\allowbreak{}}...\_column, ...\_of.
  \item Then do the task: guess first, check with \code{sdata.\allowbreak{}help} after.
\end{itemize}
\field{In the manual}{lesson~8, \manualpage{lesson:8}{41}{chapter.II.8}; ``How the words are built'', \manualpage{start:grammar}{5}{section*.7}.}
\field{Words to look up}{\kind{chart}\,\code{chart.\allowbreak{}second\_\allowbreak{}y()}~\manualpage{word:second-y}{139}{section*.79}; \kind{chart}\,\code{chart.\allowbreak{}size()}~\manualpage{word:size}{142}{section*.87}; \kind{chart}\,\code{chart.\allowbreak{}style()}~\manualpage{word:style}{142}{section*.86}; \kind{chart}\,\code{chart.\allowbreak{}note()}~\manualpage{word:note}{133}{section*.63}; \kind{chart}\,\code{chart.\allowbreak{}save()}~\manualpage{word:save}{112}{section*.18}; \kind{chart}\,\code{chart.\allowbreak{}hide\_\allowbreak{}legend()}~\manualpage{word:hide-legend}{145}{section*.94}; \kind{chart}\,\code{chart.\allowbreak{}legend\_\allowbreak{}left()}~\manualpage{word:legend-left}{144}{section*.92}; \kind{sdata}\,\code{sdata.\allowbreak{}words()}~\manualpage{word:sdata.words}{238}{section*.332}.}
\end{todo}
\begin{tryit}
\begin{enumerate}
  \item A figure for a journal: dots of \code{\textquotedbl{}hours\textquotedbl{}} against \code{\textquotedbl{}marks\textquotedbl{}}, style \code{\textquotedbl{}paper\textquotedbl{}}, 8.5 cm wide and 7 cm high, font size 8, saved as a \code{.\allowbreak{}tif} with 600 dpi.
  \item Save one chart as \code{.\allowbreak{}png}, \code{.\allowbreak{}pdf} and \code{.\allowbreak{}svg}. Then guess, before you look anything up: which word hides the legend, and which puts it on the left? Try your guesses on a chart with \code{colour=}.
\end{enumerate}
{\small Worked answers: \hyperref[answers:beginner:11]{page~\pageref*{answers:beginner:11}}, or \code{sdata.\allowbreak{}answers(11)}.}
\end{tryit}

\daystart{12}{Open a table and look at it}{Beginner track}{day:beginner:12}
\goal{Open a table, look at its first rows and columns, and count its rows. About an hour.}
\begin{warmup}
\begin{enumerate}
  \item Which word adds a second y-axis on the right? \fromday{11}\writeline
  \item Which word draws a line across at y = 50, with a label? \fromday{9}\writeline
  \item How do you give each class its own colour in a bar chart? \fromday{5}\writeline
\end{enumerate}
{\small Answers: \hyperref[answers:beginner:12]{page~\pageref*{answers:beginner:12}}, or \code{sdata.\allowbreak{}answers(12)}.}
\end{warmup}
\begin{todo}
\begin{itemize}
  \item Work through \code{sdata.\allowbreak{}lesson(9)}: type each example yourself (do not copy and paste), run it and read what each line does. Then do the lesson's tasks and check them with \code{sdata.\allowbreak{}lesson(9,\allowbreak{} solutions=True)}.
\end{itemize}
\field{In the manual}{lesson~9, \manualpage{lesson:9}{48}{chapter.II.9}.}
\field{Words to look up}{\kind{sdata}\,\code{sdata.\allowbreak{}table()}~\manualpage{word:sdata.table}{151}{section*.112}; \kind{table}\,\code{table.\allowbreak{}show\_\allowbreak{}first\_\allowbreak{}rows()}~\manualpage{word:table.show-first-rows}{152}{section*.113}; \kind{table}\,\code{table.\allowbreak{}describe\_\allowbreak{}columns()}~\manualpage{word:table.describe-columns}{153}{section*.117}; \kind{table}\,\code{table.\allowbreak{}count\_\allowbreak{}rows()}~\manualpage{word:table.count-rows}{176}{section*.171}.}
\end{todo}
\begin{tryit}
\begin{enumerate}
  \item Open the \code{\textquotedbl{}shop\textquotedbl{}} practice table. How many rows does it have, and how many rows does each shop have?
\end{enumerate}
{\small Worked answers: \hyperref[answers:beginner:12]{page~\pageref*{answers:beginner:12}}, or \code{sdata.\allowbreak{}answers(12)}.}
\end{tryit}

\daystart{13}{Choose rows and columns, and undo}{Beginner track}{day:beginner:13}
\goal{Keep and remove rows and columns, sort, and take back any step. About an hour.}
\begin{warmup}
\begin{enumerate}
  \item How do you look at the first 5 rows of a table? \fromday{12}\writeline
  \item Which word puts two charts next to each other on one page? \fromday{10}\writeline
  \item Which chart shows how the values of one column are spread, in bins? \fromday{6}\writeline
\end{enumerate}
{\small Answers: \hyperref[answers:beginner:13]{page~\pageref*{answers:beginner:13}}, or \code{sdata.\allowbreak{}answers(13)}.}
\end{warmup}
\begin{todo}
\begin{itemize}
  \item Work through \code{sdata.\allowbreak{}lesson(10)}: type each example yourself (do not copy and paste), run it and read what each line does. Then do the lesson's tasks and check them with \code{sdata.\allowbreak{}lesson(10,\allowbreak{} solutions=True)}.
\end{itemize}
\field{In the manual}{lesson~10, \manualpage{lesson:10}{52}{chapter.II.10}.}
\field{Words to look up}{\kind{table}\,\code{table.\allowbreak{}keep\_\allowbreak{}rows()}~\manualpage{word:table.keep-rows}{157}{section*.128}; \kind{table}\,\code{table.\allowbreak{}keep\_\allowbreak{}columns()}~\manualpage{word:table.keep-columns}{155}{section*.122}; \kind{table}\,\code{table.\allowbreak{}sort\_\allowbreak{}by()}~\manualpage{word:table.sort-by}{169}{section*.157}; \kind{table}\,\code{table.\allowbreak{}undo()}~\manualpage{word:table.undo}{211}{section*.250}.}
\end{todo}
\begin{tryit}
\begin{enumerate}
  \item Weather table: keep the city \code{\textquotedbl{}Sunnyvale\textquotedbl{}}, keep only the columns month, temperature and rain, sort by rain with the biggest first, and save it as \code{\textquotedbl{}sunny.\allowbreak{}csv\textquotedbl{}}.
\end{enumerate}
{\small Worked answers: \hyperref[answers:beginner:13]{page~\pageref*{answers:beginner:13}}, or \code{sdata.\allowbreak{}answers(13)}.}
\end{tryit}

\daystart{14}{Review of the second week}{Beginner track}{day:beginner:14}
\goal{Remember the words of days 8 to 13 without looking, and use them together. About an hour.}
\begin{warmup}
\begin{enumerate}
  \item Which word keeps only the rows where subject is Maths? \fromday{13}\writeline
  \item Which word draws one bar for each name, as tall as the marks? \fromday{3}\writeline
  \item Which word draws a straight line through the dots? \fromday{8}\writeline
  \item Which word sets the size of a figure in cm? \fromday{11}\writeline
  \item A few values are much bigger than the rest. Which word helps? \fromday{9}\writeline
  \item Which word describes each column: its kind, empty cells and numbers? \fromday{12}\writeline
  \item How do you put four charts in two rows, with the letters A to D? \fromday{10}\writeline
  \item How do you give the number 12 the name age? \fromday{1}\writeline
\end{enumerate}
{\small Answers: \hyperref[answers:beginner:14]{page~\pageref*{answers:beginner:14}}, or \code{sdata.\allowbreak{}answers(14)}.}
\end{warmup}
\begin{todo}
\begin{itemize}
  \item Answer the warm-up on paper, from memory, before you look anything up.
  \item For each question you could not answer, open its day with \code{sdata.\allowbreak{}day(n)} and run its example again.
  \item Then do the mixed task without looking back.
\end{itemize}
\end{todo}
\begin{tryit}
\begin{enumerate}
  \item Weather table: keep the city \code{\textquotedbl{}Frostby\textquotedbl{}}, sort by \code{\textquotedbl{}temperature\textquotedbl{}}, and draw a line of the temperature over the month, with a line at y = 0 and the label \code{\textquotedbl{}freezing\textquotedbl{}}.
\end{enumerate}
{\small Worked answers: \hyperref[answers:beginner:14]{page~\pageref*{answers:beginner:14}}, or \code{sdata.\allowbreak{}answers(14)}.}
\end{tryit}

\daystart{15}{Clean the data and add columns}{Beginner track}{day:beginner:15}
\goal{Find and fill empty cells, fix text and numbers, and add new columns. About an hour.}
\begin{warmup}
\begin{enumerate}
  \item Which word sorts with the biggest marks first? \fromday{13}\writeline
  \item Which word counts the rows of each fruit? \fromday{12}\writeline
  \item Which word checks by eye whether values are bell-shaped? \fromday{8}\writeline
  \item What does \code{print(age)} do? \fromday{1}\writeline
\end{enumerate}
{\small Answers: \hyperref[answers:beginner:15]{page~\pageref*{answers:beginner:15}}, or \code{sdata.\allowbreak{}answers(15)}.}
\end{warmup}
\begin{todo}
\begin{itemize}
  \item Work through \code{sdata.\allowbreak{}lesson(11)}: type each example yourself (do not copy and paste), run it and read what each line does. Then do the lesson's tasks and check them with \code{sdata.\allowbreak{}lesson(11,\allowbreak{} solutions=True)}.
\end{itemize}
\field{In the manual}{lesson~11, \manualpage{lesson:11}{57}{chapter.II.11}; ``Missing values and imputation'', \manualpage{ideas:missing}{248}{section*.365}.}
\field{Words to look up}{\kind{table}\,\code{table.\allowbreak{}count\_\allowbreak{}empty\_\allowbreak{}cells()}~\manualpage{word:table.count-empty-cells}{163}{section*.142}; \kind{table}\,\code{table.\allowbreak{}fill\_\allowbreak{}empty\_\allowbreak{}cells()}~\manualpage{word:table.fill-empty-cells}{164}{section*.144}; \kind{table}\,\code{table.\allowbreak{}make\_\allowbreak{}numbers()}~\manualpage{word:table.make-numbers}{165}{section*.147}; \kind{table}\,\code{table.\allowbreak{}add\_\allowbreak{}column()}~\manualpage{word:table.add-column}{170}{section*.158}.}
\end{todo}
\begin{tryit}
\begin{enumerate}
  \item Shop table: count the empty cells, fill the empty \code{\textquotedbl{}sales\textquotedbl{}} with the line between neighbours for each shop (sorted by \code{\textquotedbl{}date\textquotedbl{}}), and add a column \code{\textquotedbl{}weekday\textquotedbl{}} with the name of the weekday.
\end{enumerate}
{\small Worked answers: \hyperref[answers:beginner:15]{page~\pageref*{answers:beginner:15}}, or \code{sdata.\allowbreak{}answers(15)}.}
\end{tryit}

\daystart{16}{Summaries, and from a table to a chart}{Beginner track}{day:beginner:16}
\goal{Sum up columns for each group, and draw a summary as a chart. About an hour.}
\begin{warmup}
\begin{enumerate}
  \item Which word counts the empty cells in each column? \fromday{15}\writeline
  \item You made a mistake in the last step. Which word takes it back? \fromday{13}\writeline
  \item Which word shows the numbers on the y-axis as percents? \fromday{9}\writeline
  \item In \code{nums.\allowbreak{}keep\_\allowbreak{}numbers(above=50)}, what is \code{above=50}? \fromday{2}\writeline
\end{enumerate}
{\small Answers: \hyperref[answers:beginner:16]{page~\pageref*{answers:beginner:16}}, or \code{sdata.\allowbreak{}answers(16)}.}
\end{warmup}
\begin{todo}
\begin{itemize}
  \item Work through \code{sdata.\allowbreak{}lesson(12)}: type each example yourself (do not copy and paste), run it and read what each line does. Then do the lesson's tasks and check them with \code{sdata.\allowbreak{}lesson(12,\allowbreak{} solutions=True)}.
\end{itemize}
\field{In the manual}{lesson~12, \manualpage{lesson:12}{64}{chapter.II.12}.}
\field{Words to look up}{\kind{table}\,\code{table.\allowbreak{}average\_\allowbreak{}of()}~\manualpage{word:table.average-of}{177}{section*.173}; \kind{table}\,\code{table.\allowbreak{}summary\_\allowbreak{}of()}~\manualpage{word:table.summary-of}{179}{section*.179}; \kind{table}\,\code{table.\allowbreak{}correlation\_\allowbreak{}between()}~\manualpage{word:table.correlation-between}{182}{section*.184}; \kind{table}\,\code{table.\allowbreak{}join\_\allowbreak{}with()}~\manualpage{word:table.join-with}{209}{section*.245}.}
\end{todo}
\begin{tryit}
\begin{enumerate}
  \item Fruit table: the total \code{\textquotedbl{}count\textquotedbl{}} of each fruit, drawn as bars with the numbers on the bars.
\end{enumerate}
{\small Worked answers: \hyperref[answers:beginner:16]{page~\pageref*{answers:beginner:16}}, or \code{sdata.\allowbreak{}answers(16)}.}
\end{tryit}

\daystart{17}{Lists and grids of numbers}{Beginner track}{day:beginner:17}
\goal{Do maths on lists and grids of numbers, and make random numbers. About an hour.}
\begin{warmup}
\begin{enumerate}
  \item Which word gives the average of a column for each group? \fromday{16}\writeline
  \item Which word keeps rows 10 to 20, by their place in the table? \fromday{13}\writeline
  \item Which word draws one small chart for each class? \fromday{10}\writeline
  \item How do you save a chart as a PDF? \fromday{3}\writeline
\end{enumerate}
{\small Answers: \hyperref[answers:beginner:17]{page~\pageref*{answers:beginner:17}}, or \code{sdata.\allowbreak{}answers(17)}.}
\end{warmup}
\begin{todo}
\begin{itemize}
  \item Work through \code{sdata.\allowbreak{}lesson(13)}: type each example yourself (do not copy and paste), run it and read what each line does. Then do the lesson's tasks and check them with \code{sdata.\allowbreak{}lesson(13,\allowbreak{} solutions=True)}.
\end{itemize}
\field{In the manual}{lesson~13, \manualpage{lesson:13}{72}{chapter.II.13}.}
\field{Words to look up}{\kind{sdata}\,\code{sdata.\allowbreak{}numbers()}~\manualpage{word:sdata.numbers}{214}{section*.260}; \kind{sdata}\,\code{sdata.\allowbreak{}random\_\allowbreak{}numbers()}~\manualpage{word:sdata.random-numbers}{214}{section*.261}; \kind{nums}\,\code{nums.\allowbreak{}median()}~\manualpage{word:nums.median}{226}{section*.296}; \kind{nums}\,\code{nums.\allowbreak{}percentile()}~\manualpage{word:nums.percentile}{228}{section*.301}.}
\end{todo}
\begin{tryit}
\begin{enumerate}
  \item Make 200 random heights with average 150 and standard deviation 8 (\code{seed=1}). What is their 90th percentile? Draw them as a histogram.
\end{enumerate}
{\small Worked answers: \hyperref[answers:beginner:17]{page~\pageref*{answers:beginner:17}}, or \code{sdata.\allowbreak{}answers(17)}.}
\end{tryit}

\daystart{18}{Is a difference real? Tests and effect sizes}{Beginner track}{day:beginner:18}
\goal{Test whether averages differ, and say how big the difference is. About an hour.}
\begin{warmup}
\begin{enumerate}
  \item How do you make simplydata numbers from a list? \fromday{17}\writeline
  \item How do you fill the empty marks with the median? \fromday{15}\writeline
  \item How do you save a picture sharp enough for a journal? \fromday{11}\writeline
  \item Which word puts a title above the chart? \fromday{4}\writeline
\end{enumerate}
{\small Answers: \hyperref[answers:beginner:18]{page~\pageref*{answers:beginner:18}}, or \code{sdata.\allowbreak{}answers(18)}.}
\end{warmup}
\begin{todo}
\begin{itemize}
  \item Work through \code{sdata.\allowbreak{}lesson(14)}, examples 1 and 3: type each example yourself (do not copy and paste), run it and read what each line does. Then do the lesson's tasks and check them with \code{sdata.\allowbreak{}lesson(14,\allowbreak{} solutions=True)}.
  \item Read \code{\textquotedbl{}The null hypothesis and the p-value\textquotedbl{}} and \code{\textquotedbl{}Effect sizes\textquotedbl{}} in the manual.
\end{itemize}
\field{In the manual}{lesson~14, \manualpage{lesson:14}{76}{chapter.II.14}; ``The null hypothesis and the p-value'', \manualpage{ideas:p-value}{245}{section*.353}; ``Effect sizes: how big is the difference?'', \manualpage{ideas:effect-size}{246}{section*.356}.}
\field{Words to look up}{\kind{table}\,\code{table.\allowbreak{}compare\_\allowbreak{}averages()}~\manualpage{word:table.compare-averages}{183}{section*.186}; \kind{table}\,\code{table.\allowbreak{}compare\_\allowbreak{}each\_\allowbreak{}pair()}~\manualpage{word:table.compare-each-pair}{183}{section*.187}; \kind{table}\,\code{table.\allowbreak{}fit\_\allowbreak{}line()}~\manualpage{word:table.fit-line}{188}{section*.197}.}
\end{todo}
\begin{tryit}
\begin{enumerate}
  \item School table: do the average marks differ between pupils who had breakfast and pupils who did not? Write down the difference, the effect size and the p-value.
\end{enumerate}
{\small Worked answers: \hyperref[answers:beginner:18]{page~\pageref*{answers:beginner:18}}, or \code{sdata.\allowbreak{}answers(18)}.}
\end{tryit}

\daystart{19}{Counts, outliers and many tests}{Beginner track}{day:beginner:19}
\goal{Test counts of categories, and check shapes and outliers. About an hour.}
\begin{warmup}
\begin{enumerate}
  \item Which word tests whether the average marks differ between the classes? \fromday{18}\writeline
  \item Which word tells whether two columns go up and down together? \fromday{16}\writeline
  \item Which word counts the rows? \fromday{12}\writeline
  \item Which word moves the legend below the chart? \fromday{5}\writeline
\end{enumerate}
{\small Answers: \hyperref[answers:beginner:19]{page~\pageref*{answers:beginner:19}}, or \code{sdata.\allowbreak{}answers(19)}.}
\end{warmup}
\begin{todo}
\begin{itemize}
  \item Work through \code{sdata.\allowbreak{}lesson(14)}, examples 2, 4 and 5: type each example yourself (do not copy and paste), run it and read what each line does. Then do the lesson's tasks and check them with \code{sdata.\allowbreak{}lesson(14,\allowbreak{} solutions=True)}.
  \item Read \code{\textquotedbl{}Many tests at once\textquotedbl{}} in the manual.
\end{itemize}
\field{In the manual}{lesson~14, \manualpage{lesson:14}{76}{chapter.II.14}; ``Many tests at once'', \manualpage{ideas:multiple}{246}{section*.355}.}
\field{Words to look up}{\kind{table}\,\code{table.\allowbreak{}compare\_\allowbreak{}counts()}~\manualpage{word:table.compare-counts}{185}{section*.190}; \kind{table}\,\code{table.\allowbreak{}show\_\allowbreak{}rows\_\allowbreak{}with\_\allowbreak{}outliers()}~\manualpage{word:table.show-rows-with-outliers}{162}{section*.139}; \kind{table}\,\code{table.\allowbreak{}compare\_\allowbreak{}shape\_\allowbreak{}to()}~\manualpage{word:table.compare-shape-to}{185}{section*.191}; \kind{table}\,\code{table.\allowbreak{}add\_\allowbreak{}adjusted\_\allowbreak{}p\_\allowbreak{}value\_\allowbreak{}column()}~\manualpage{word:table.add-adjusted-p-value-column}{186}{section*.192}.}
\end{todo}
\begin{tryit}
\begin{enumerate}
  \item School table: is the share of pupils who \code{\textquotedbl{}passed\textquotedbl{}} different with and without breakfast?
\end{enumerate}
{\small Worked answers: \hyperref[answers:beginner:19]{page~\pageref*{answers:beginner:19}}, or \code{sdata.\allowbreak{}answers(19)}.}
\end{tryit}

\daystart{20}{Choosing a test, and a whole analysis}{Beginner track}{day:beginner:20}
\goal{Choose the right test, and go from a question to a tested answer and a figure. About an hour.}
\begin{warmup}
\begin{enumerate}
  \item Which word tests whether the counts of categories differ between groups? \fromday{19}\writeline
  \item Which word gives the middle value of numbers? \fromday{17}\writeline
  \item Which word lists every step you did to a table? \fromday{13}\writeline
  \item Which two words show the spread of each group, with every value as a dot on top? \fromday{6}\writeline
\end{enumerate}
{\small Answers: \hyperref[answers:beginner:20]{page~\pageref*{answers:beginner:20}}, or \code{sdata.\allowbreak{}answers(20)}.}
\end{warmup}
\begin{todo}
\begin{itemize}
  \item Read \code{\textquotedbl{}Which test when?\textquotedbl{}} and \code{\textquotedbl{}Writing it up\textquotedbl{}} in the manual. Answer its four questions for each task below before you run anything.
  \item Run \code{sdata.\allowbreak{}help(\textquotedbl{}compare\_\allowbreak{}before\_\allowbreak{}after\textquotedbl{})} and \code{sdata.\allowbreak{}help(\textquotedbl{}compare\_\allowbreak{}average\_\allowbreak{}to\textquotedbl{})} and type their examples.
  \item Work through \code{sdata.\allowbreak{}lesson(15)}: type each example yourself (do not copy and paste), run it and read what each line does. Then do the lesson's tasks and check them with \code{sdata.\allowbreak{}lesson(15,\allowbreak{} solutions=True)}.
\end{itemize}
\field{In the manual}{lesson~15, \manualpage{lesson:15}{81}{chapter.II.15}; ``Which test when?'', \manualpage{ideas:which-test}{249}{section*.367}; ``Writing it up'', \manualpage{ideas:report}{250}{section*.368}.}
\field{Words to look up}{\kind{table}\,\code{table.\allowbreak{}compare\_\allowbreak{}average\_\allowbreak{}to()}~\manualpage{word:table.compare-average-to}{184}{section*.189}; \kind{table}\,\code{table.\allowbreak{}compare\_\allowbreak{}before\_\allowbreak{}after()}~\manualpage{word:table.compare-before-after}{184}{section*.188}; \kind{table}\,\code{table.\allowbreak{}compare\_\allowbreak{}percent\_\allowbreak{}to()}~\manualpage{word:table.compare-percent-to}{187}{section*.195}; \kind{table}\,\code{table.\allowbreak{}steps()}~\manualpage{word:table.steps}{211}{section*.251}; \kind{sdata}\,\code{sdata.\allowbreak{}grid()}~\manualpage{word:sdata.grid}{149}{section*.106}; \kind{table}\,\code{table.\allowbreak{}save()}~\manualpage{word:table.save}{212}{section*.253}; \kind{table}\,\code{table.\allowbreak{}as\_\allowbreak{}pandas()}~\manualpage{word:table.as-pandas}{212}{section*.255}.}
\end{todo}
\begin{tryit}
\begin{enumerate}
  \item Marks table: is each child's Maths mark different from the same child's Science mark? Choose the right word first.
  \item School table: is the average \code{\textquotedbl{}sleep\textquotedbl{}} different from 8 hours?
  \item From memory: weather table, compare the average \code{\textquotedbl{}temperature\textquotedbl{}} of the cities, draw boxes with all the dots, and save the figure as a PDF and the table as \code{\textquotedbl{}weather.\allowbreak{}xlsx\textquotedbl{}}.
\end{enumerate}
{\small Worked answers: \hyperref[answers:beginner:20]{page~\pageref*{answers:beginner:20}}, or \code{sdata.\allowbreak{}answers(20)}.}
\end{tryit}

\daystart{21}{Review of the third week}{Beginner track}{day:beginner:21}
\goal{Remember the words of days 15 to 20 without looking, and use them together. About an hour.}
\begin{warmup}
\begin{enumerate}
  \item The same children were measured twice. Which word compares the two? \fromday{20}\writeline
  \item Which word turns a column of text such as \code{\textquotedbl{}144\textquotedbl{}} into numbers? \fromday{15}\writeline
  \item Which word writes a note with an arrow on the chart? \fromday{11}\writeline
  \item Which word says which pairs of groups differ? \fromday{18}\writeline
  \item Which word joins two tables on a column they share? \fromday{16}\writeline
  \item Which word shows the rows with outliers? \fromday{19}\writeline
  \item How do you make 100 random numbers with average 50 and standard deviation 10, the same every time? \fromday{17}\writeline
  \item Which option makes a bar show the median instead of the average? \fromday{8}\writeline
  \item How do you write the text Asha in Python? \fromday{1}\writeline
\end{enumerate}
{\small Answers: \hyperref[answers:beginner:21]{page~\pageref*{answers:beginner:21}}, or \code{sdata.\allowbreak{}answers(21)}.}
\end{warmup}
\begin{todo}
\begin{itemize}
  \item Answer the warm-up on paper, from memory, before you look anything up.
  \item For each question you could not answer, open its day with \code{sdata.\allowbreak{}day(n)} and run its example again.
  \item Then do the mixed task without looking back.
\end{itemize}
\end{todo}
\begin{tryit}
\begin{enumerate}
  \item School table: the average \code{\textquotedbl{}marks\textquotedbl{}} for each value of \code{\textquotedbl{}breakfast\textquotedbl{}}, a test of the difference, and boxes of the marks for each group.
\end{enumerate}
{\small Worked answers: \hyperref[answers:beginner:21]{page~\pageref*{answers:beginner:21}}, or \code{sdata.\allowbreak{}answers(21)}.}
\end{tryit}

\daystart{22}{Models: explain one column from several}{Beginner track}{day:beginner:22}
\goal{Explain one column from several others with a model. About an hour.}
\begin{warmup}
\begin{enumerate}
  \item Which word compares one average with a number you choose? \fromday{20}\writeline
  \item Which word checks whether values are bell-shaped (normal)? \fromday{19}\writeline
  \item Which word adds a new column made from other columns? \fromday{15}\writeline
  \item Which option adds error bars that say how sure we are of the average? \fromday{8}\writeline
\end{enumerate}
{\small Answers: \hyperref[answers:beginner:22]{page~\pageref*{answers:beginner:22}}, or \code{sdata.\allowbreak{}answers(22)}.}
\end{warmup}
\begin{todo}
\begin{itemize}
  \item Work through \code{sdata.\allowbreak{}lesson(16)}: type each example yourself (do not copy and paste), run it and read what each line does. Then do the lesson's tasks and check them with \code{sdata.\allowbreak{}lesson(16,\allowbreak{} solutions=True)}.
\end{itemize}
\field{In the manual}{lesson~16, \manualpage{lesson:16}{86}{chapter.II.16}.}
\field{Words to look up}{\kind{table}\,\code{table.\allowbreak{}explain()}~\manualpage{word:table.explain}{200}{section*.227}; \kind{table}\,\code{table.\allowbreak{}add\_\allowbreak{}predicted\_\allowbreak{}column()}~\manualpage{word:table.add-predicted-column}{200}{section*.228}.}
\end{todo}
\begin{tryit}
\begin{enumerate}
  \item School table: explain \code{\textquotedbl{}days\_\allowbreak{}absent\textquotedbl{}} using \code{\textquotedbl{}sleep\textquotedbl{}} and \code{\textquotedbl{}breakfast\textquotedbl{}} with \code{outcome=\textquotedbl{}counts\textquotedbl{}}. Does breakfast matter?
\end{enumerate}
{\small Worked answers: \hyperref[answers:beginner:22]{page~\pageref*{answers:beginner:22}}, or \code{sdata.\allowbreak{}answers(22)}.}
\end{tryit}

\daystart{23}{Patterns across many columns}{Beginner track}{day:beginner:23}
\goal{Find patterns across many columns with principal components and clusters. About an hour.}
\begin{warmup}
\begin{enumerate}
  \item Which word explains marks from hours, sleep and breakfast together? \fromday{22}\writeline
  \item What should a results sentence give? \fromday{20}\writeline
  \item Which word gives several summaries at once? \fromday{16}\writeline
  \item Which word sets the y-axis from 0 to 100? \fromday{9}\writeline
\end{enumerate}
{\small Answers: \hyperref[answers:beginner:23]{page~\pageref*{answers:beginner:23}}, or \code{sdata.\allowbreak{}answers(23)}.}
\end{warmup}
\begin{todo}
\begin{itemize}
  \item Work through \code{sdata.\allowbreak{}lesson(17)}: type each example yourself (do not copy and paste), run it and read what each line does. Then do the lesson's tasks and check them with \code{sdata.\allowbreak{}lesson(17,\allowbreak{} solutions=True)}.
\end{itemize}
\field{In the manual}{lesson~17, \manualpage{lesson:17}{90}{chapter.II.17}.}
\field{Words to look up}{\kind{table}\,\code{table.\allowbreak{}add\_\allowbreak{}principal\_\allowbreak{}components()}~\manualpage{word:table.add-principal-components}{201}{section*.230}; \kind{table}\,\code{table.\allowbreak{}add\_\allowbreak{}clusters()}~\manualpage{word:table.add-clusters}{203}{section*.232}; \kind{table}\,\code{table.\allowbreak{}add\_\allowbreak{}nmds()}~\manualpage{word:table.add-nmds}{204}{section*.234}; \kind{table}\,\code{table.\allowbreak{}compare\_\allowbreak{}profiles()}~\manualpage{word:table.compare-profiles}{205}{section*.235}.}
\end{todo}
\begin{tryit}
\begin{enumerate}
  \item Meadow table: add 3 clusters and count the rows of each cluster across \code{\textquotedbl{}field\textquotedbl{}}.
\end{enumerate}
{\small Worked answers: \hyperref[answers:beginner:23]{page~\pageref*{answers:beginner:23}}, or \code{sdata.\allowbreak{}answers(23)}.}
\end{tryit}

\daystart{24}{Survival, and planning a study}{Beginner track}{day:beginner:24}
\goal{Draw survival curves, and work out how many values a study needs. About an hour.}
\begin{warmup}
\begin{enumerate}
  \item Which word adds the principal components PC1 and PC2? \fromday{23}\writeline
  \item Which word lists every change you made to a table, for your methods? \fromday{20}\writeline
  \item How do you get the number in row 1, column 2 of a grid? \fromday{17}\writeline
  \item Which word puts two charts next to each other on one page? \fromday{10}\writeline
\end{enumerate}
{\small Answers: \hyperref[answers:beginner:24]{page~\pageref*{answers:beginner:24}}, or \code{sdata.\allowbreak{}answers(24)}.}
\end{warmup}
\begin{todo}
\begin{itemize}
  \item Work through \code{sdata.\allowbreak{}lesson(18)}: type each example yourself (do not copy and paste), run it and read what each line does. Then do the lesson's tasks and check them with \code{sdata.\allowbreak{}lesson(18,\allowbreak{} solutions=True)}.
\end{itemize}
\field{In the manual}{lesson~18, \manualpage{lesson:18}{95}{chapter.II.18}; ``Two kinds of mistake, and power'', \manualpage{ideas:errors}{246}{section*.354}.}
\field{Words to look up}{\kind{table}\,\code{table.\allowbreak{}survival\_\allowbreak{}over\_\allowbreak{}time()}~\manualpage{word:table.survival-over-time}{205}{section*.237}; \kind{chart}\,\code{chart.\allowbreak{}survival\_\allowbreak{}curves()}~\manualpage{word:survival-curves}{207}{section*.240}; \kind{table}\,\code{table.\allowbreak{}compare\_\allowbreak{}survival()}~\manualpage{word:table.compare-survival}{206}{section*.238}; \kind{sdata}\,\code{sdata.\allowbreak{}how\_\allowbreak{}many\_\allowbreak{}needed()}~\manualpage{word:sdata.how-many-needed}{207}{section*.242}.}
\end{todo}
\begin{tryit}
\begin{enumerate}
  \item How many values does each group need to find a large difference (\code{effect\_\allowbreak{}size=0.\allowbreak{}8}) with power 0.8? And with a small one (0.2)?
\end{enumerate}
{\small Worked answers: \hyperref[answers:beginner:24]{page~\pageref*{answers:beginner:24}}, or \code{sdata.\allowbreak{}answers(24)}.}
\end{tryit}

\daystart{25}{Fitting, finding and solving}{Beginner track}{day:beginner:25}
\goal{Fit lines, curves and formulas, find where a formula is biggest, and solve equations. About an hour.}
\begin{warmup}
\begin{enumerate}
  \item Which word draws survival curves? \fromday{24}\writeline
  \item How do you explain a yes-or-no column such as \code{\textquotedbl{}passed\textquotedbl{}}? \fromday{22}\writeline
  \item A p-value answers \code{\textquotedbl{}could it be chance?\textquotedbl{}}. What answers \code{\textquotedbl{}how big is it?\textquotedbl{}} \fromday{18}\writeline
  \item If \code{keep\_\allowbreak{}rows} keeps rows, what would you guess takes columns away? \fromday{11}\writeline
\end{enumerate}
{\small Answers: \hyperref[answers:beginner:25]{page~\pageref*{answers:beginner:25}}, or \code{sdata.\allowbreak{}answers(25)}.}
\end{warmup}
\begin{todo}
\begin{itemize}
  \item Work through \code{sdata.\allowbreak{}lesson(19)}, examples 1 and 2: type each example yourself (do not copy and paste), run it and read what each line does. Then do the lesson's tasks and check them with \code{sdata.\allowbreak{}lesson(19,\allowbreak{} solutions=True)}.
  \item Work through \code{sdata.\allowbreak{}lesson(19)}, example 3: type each example yourself (do not copy and paste), run it and read what each line does. Then do the lesson's tasks and check them with \code{sdata.\allowbreak{}lesson(19,\allowbreak{} solutions=True)}.
\end{itemize}
\field{In the manual}{lesson~19, \manualpage{lesson:19}{99}{chapter.II.19}; ``Fitting a line or a formula'', \manualpage{ideas:regression}{247}{section*.361}.}
\field{Words to look up}{\kind{table}\,\code{table.\allowbreak{}fit\_\allowbreak{}line()}~\manualpage{word:table.fit-line}{188}{section*.197}; \kind{table}\,\code{table.\allowbreak{}fit\_\allowbreak{}curve()}~\manualpage{word:table.fit-curve}{189}{section*.198}; \kind{table}\,\code{table.\allowbreak{}fit\_\allowbreak{}formula()}~\manualpage{word:table.fit-formula}{189}{section*.199}; \kind{chart}\,\code{chart.\allowbreak{}fit\_\allowbreak{}formula()}~\manualpage{word:fit-formula}{123}{section*.41}; \kind{sdata}\,\code{sdata.\allowbreak{}find\_\allowbreak{}smallest()}~\manualpage{word:sdata.find-smallest}{190}{section*.201}; \kind{sdata}\,\code{sdata.\allowbreak{}find\_\allowbreak{}biggest()}~\manualpage{word:sdata.find-biggest}{191}{section*.202}; \kind{sdata}\,\code{sdata.\allowbreak{}solve()}~\manualpage{word:sdata.solve}{191}{section*.203}.}
\end{todo}
\begin{tryit}
\begin{enumerate}
  \item Growth table: fit the formula \code{\textquotedbl{}forest = a + b * (year - 2000)\textquotedbl{}} for each country. Which country loses forest fastest?
  \item A box has a square bottom of side x and a height of 10 - x, so its volume is x * x * (10 - x). For which x between 0 and 10 is the volume biggest?
\end{enumerate}
{\small Worked answers: \hyperref[answers:beginner:25]{page~\pageref*{answers:beginner:25}}, or \code{sdata.\allowbreak{}answers(25)}.}
\end{tryit}

\daystart{26}{Distributions and values between points}{Beginner track}{day:beginner:26}
\goal{Work out chances from a distribution, and estimate values between measured points. About an hour.}
\begin{warmup}
\begin{enumerate}
  \item Which word fits a straight line on a table? \fromday{25}\writeline
  \item Which word puts rows that are alike into groups? \fromday{23}\writeline
  \item You did many tests at once. Which word corrects the p-values? \fromday{19}\writeline
  \item How do you look at the first 5 rows of a table? \fromday{12}\writeline
\end{enumerate}
{\small Answers: \hyperref[answers:beginner:26]{page~\pageref*{answers:beginner:26}}, or \code{sdata.\allowbreak{}answers(26)}.}
\end{warmup}
\begin{todo}
\begin{itemize}
  \item Work through \code{sdata.\allowbreak{}lesson(20)}, examples 1 and 2: type each example yourself (do not copy and paste), run it and read what each line does. Then do the lesson's tasks and check them with \code{sdata.\allowbreak{}lesson(20,\allowbreak{} solutions=True)}.
  \item Work through \code{sdata.\allowbreak{}lesson(20)}, example 3: type each example yourself (do not copy and paste), run it and read what each line does. Then do the lesson's tasks and check them with \code{sdata.\allowbreak{}lesson(20,\allowbreak{} solutions=True)}.
\end{itemize}
\field{In the manual}{lesson~20, \manualpage{lesson:20}{104}{chapter.II.20}; ``The shape of the values: distributions'', \manualpage{ideas:distribution}{244}{section*.345}.}
\field{Words to look up}{\kind{sdata}\,\code{sdata.\allowbreak{}distribution()}~\manualpage{word:sdata.distribution}{193}{section*.208}; \kind{dist}\,\code{dist.\allowbreak{}chance()}~\manualpage{word:dist.chance}{195}{section*.212}; \kind{dist}\,\code{dist.\allowbreak{}percentile()}~\manualpage{word:dist.percentile}{196}{section*.214}; \kind{table}\,\code{table.\allowbreak{}fit\_\allowbreak{}distribution()}~\manualpage{word:table.fit-distribution}{194}{section*.209}; \kind{table}\,\code{table.\allowbreak{}values\_\allowbreak{}at()}~\manualpage{word:table.values-at}{192}{section*.205}; \kind{table}\,\code{table.\allowbreak{}map\_\allowbreak{}of()}~\manualpage{word:table.map-of}{193}{section*.206}.}
\end{todo}
\begin{tryit}
\begin{enumerate}
  \item Heights are normal with average 150 and standard deviation 8. What share is taller than 165? Which height do 95 \% of them stay below?
  \item Weather table: estimate the \code{\textquotedbl{}rain\textquotedbl{}} of each city at \code{month\_\allowbreak{}number} 2.5 and 7.5, joined by straight lines.
\end{enumerate}
{\small Worked answers: \hyperref[answers:beginner:26]{page~\pageref*{answers:beginner:26}}, or \code{sdata.\allowbreak{}answers(26)}.}
\end{tryit}

\daystart{27}{Small project: look, clean and sum up}{Beginner track}{day:beginner:27}
\goal{Choose a question and a file, then look at it, clean it and sum it up. About an hour.}
\begin{warmup}
\begin{enumerate}
  \item How do you make a normal distribution with average 50 and standard deviation 10? \fromday{26}\writeline
  \item Which word cuts a column of numbers into named groups? \fromday{15}\writeline
  \item How do you give the panels of a figure the letters A and B? \fromday{20}\writeline
  \item Which word tests whether survival differs between groups? \fromday{24}\writeline
\end{enumerate}
{\small Answers: \hyperref[answers:beginner:27]{page~\pageref*{answers:beginner:27}}, or \code{sdata.\allowbreak{}answers(27)}.}
\end{warmup}
\begin{todo}
\begin{itemize}
  \item Your small project, day 1. Choose data of your own (a file from school, work or a public data set), or use the \code{\textquotedbl{}school\textquotedbl{}} practice table.
  \item Write your question in one sentence. Open the file, look at it, clean what needs cleaning and make two summaries.
\end{itemize}
\field{In the manual}{lesson~9, \manualpage{lesson:9}{48}{chapter.II.9}; lesson~11, \manualpage{lesson:11}{57}{chapter.II.11}; lesson~12, \manualpage{lesson:12}{64}{chapter.II.12}.}
\field{Words to look up}{\kind{sdata}\,\code{sdata.\allowbreak{}table()}~\manualpage{word:sdata.table}{151}{section*.112}; \kind{table}\,\code{table.\allowbreak{}describe\_\allowbreak{}columns()}~\manualpage{word:table.describe-columns}{153}{section*.117}; \kind{table}\,\code{table.\allowbreak{}count\_\allowbreak{}empty\_\allowbreak{}cells()}~\manualpage{word:table.count-empty-cells}{163}{section*.142}; \kind{table}\,\code{table.\allowbreak{}summary\_\allowbreak{}of()}~\manualpage{word:table.summary-of}{179}{section*.179}.}
\end{todo}
\begin{tryit}
\begin{enumerate}
  \item Write your question. Open your file, describe its columns, deal with the empty cells and make two summaries for each group.
\end{enumerate}
{\small Worked answers: \hyperref[answers:beginner:27]{page~\pageref*{answers:beginner:27}}, or \code{sdata.\allowbreak{}answers(27)}.}
\end{tryit}

\daystart{28}{Small project: test, figure and results}{Beginner track}{day:beginner:28}
\goal{Answer your question with a test, a figure and two sentences of results. About an hour.}
\begin{warmup}
\begin{enumerate}
  \item Which word adds the values a model predicts as a new column? \fromday{22}\writeline
  \item Which option gives the rank-based test, for small groups that are not bell-shaped? \fromday{18}\writeline
  \item A word for one axis ends in ...? \fromday{11}\writeline
  \item How do you fit a curve with x and x squared? \fromday{25}\writeline
\end{enumerate}
{\small Answers: \hyperref[answers:beginner:28]{page~\pageref*{answers:beginner:28}}, or \code{sdata.\allowbreak{}answers(28)}.}
\end{warmup}
\begin{todo}
\begin{itemize}
  \item Your small project, day 2. Test the difference, make one clear figure, and write two sentences with the numbers as they were printed.
\end{itemize}
\field{In the manual}{lesson~14, \manualpage{lesson:14}{76}{chapter.II.14}; lesson~15, \manualpage{lesson:15}{81}{chapter.II.15}; ``Writing it up'', \manualpage{ideas:report}{250}{section*.368}.}
\field{Words to look up}{\kind{table}\,\code{table.\allowbreak{}compare\_\allowbreak{}averages()}~\manualpage{word:table.compare-averages}{183}{section*.186}; \kind{chart}\,\code{chart.\allowbreak{}boxes()}~\manualpage{word:boxes}{116}{section*.25}; \kind{chart}\,\code{chart.\allowbreak{}save()}~\manualpage{word:save}{112}{section*.18}.}
\end{todo}
\begin{tryit}
\begin{enumerate}
  \item Test your question, draw one figure and save it, and write two sentences of results.
\end{enumerate}
{\small Worked answers: \hyperref[answers:beginner:28]{page~\pageref*{answers:beginner:28}}, or \code{sdata.\allowbreak{}answers(28)}.}
\end{tryit}

\daystart{29}{Final test from memory}{Beginner track}{day:beginner:29}
\goal{Answer questions on the whole course and write a whole analysis from memory. About an hour.}
\begin{warmup}
\begin{enumerate}
  \item Which word tests whether groups differ in their make-up across many columns (PERMANOVA)? \fromday{23}\writeline
  \item Which word gives the chance of a value above 80? \fromday{26}\writeline
  \item Which word says how many values you need in each group? \fromday{24}\writeline
  \item Which word hands a table on to pandas? \fromday{20}\writeline
  \item Which word fits a formula of your own, such as growth? \fromday{25}\writeline
  \item Which file endings give a picture that stays sharp at any size? \fromday{11}\writeline
  \item What does every simplydata program start with? \fromday{2}\writeline
  \item Where in an error message do you look first? \fromday{3}\writeline
  \item How do you name the y-axis? \fromday{4}\writeline
  \item Which word lets you choose the colours of the groups? \fromday{5}\writeline
  \item Which word keeps only the rows where subject is Maths? \fromday{13}\writeline
  \item Which word draws one dot for each row, x against y? \fromday{6}\writeline
  \item How do you give the number 12 the name age? \fromday{1}\writeline
  \item Which word draws a straight line through the dots? \fromday{8}\writeline
\end{enumerate}
{\small Answers: \hyperref[answers:beginner:29]{page~\pageref*{answers:beginner:29}}, or \code{sdata.\allowbreak{}answers(29)}.}
\end{warmup}
\begin{todo}
\begin{itemize}
  \item Today's warm-up is the final test: answer every question on paper, from memory. Then check your answers with \code{sdata.\allowbreak{}answers()}.
  \item Then write the whole program of the task in an empty file, without looking anything up. Only when it is finished, check it with \code{sdata.\allowbreak{}help}.
\end{itemize}
\field{Words to look up}{\kind{sdata}\,\code{sdata.\allowbreak{}words()}~\manualpage{word:sdata.words}{238}{section*.332}; \kind{sdata}\,\code{sdata.\allowbreak{}help()}~\manualpage{word:sdata.help}{238}{section*.333}.}
\end{todo}
\begin{tryit}
\begin{enumerate}
  \item From memory, in an empty file: open the \code{\textquotedbl{}marks\textquotedbl{}} table, remove the rows with empty cells, give the average marks for each class, test the difference, and save a two-panel figure (boxes, and dots of hours against marks with a fitted line) as a PDF.
\end{enumerate}
{\small Worked answers: \hyperref[answers:beginner:29]{page~\pageref*{answers:beginner:29}}, or \code{sdata.\allowbreak{}answers(29)}.}
\end{tryit}

\daystart{30}{Look back, and what comes next}{Beginner track}{day:beginner:30}
\goal{Do your hardest day again, and take a first look at pandas underneath simplydata. About an hour.}
\begin{warmup}
\begin{enumerate}
  \item Which word fits a distribution to a column? \fromday{26}\writeline
  \item Which word finds where a formula is smallest? \fromday{25}\writeline
  \item Which word adds the principal components PC1 and PC2? \fromday{23}\writeline
  \item Which word gives the average of a column for each group? \fromday{16}\writeline
\end{enumerate}
{\small Answers: \hyperref[answers:beginner:30]{page~\pageref*{answers:beginner:30}}, or \code{sdata.\allowbreak{}answers(30)}.}
\end{warmup}
\begin{todo}
\begin{itemize}
  \item Look back over the 30 days with \code{sdata.\allowbreak{}day()}: it lists every day.
  \item Choose the three days you found hardest and do their tasks again, without looking back.
  \item Take a first look at pandas, the library underneath simplydata tables: \code{table.\allowbreak{}as\_\allowbreak{}pandas()} hands your table over, and the manual's dictionary says which library does the work of each word.
\end{itemize}
\field{Words to look up}{\kind{sdata}\,\code{sdata.\allowbreak{}day()}~\manualpage{word:sdata.day}{240}{section*.338}; \kind{sdata}\,\code{sdata.\allowbreak{}words()}~\manualpage{word:sdata.words}{238}{section*.332}; \kind{table}\,\code{table.\allowbreak{}as\_\allowbreak{}pandas()}~\manualpage{word:table.as-pandas}{212}{section*.255}.}
\end{todo}
\begin{tryit}
\begin{enumerate}
  \item Do the task of your hardest day again from memory. Then hand the marks table to pandas and print its first rows with pandas' own word \code{head()}.
\end{enumerate}
{\small Worked answers: \hyperref[answers:beginner:30]{page~\pageref*{answers:beginner:30}}, or \code{sdata.\allowbreak{}answers(30)}.}
\end{tryit}

\trackstart{Answers: Python track}{The warm-up answers and a worked answer to each task}{answers:python}
There is often more than one good answer: if yours runs and answers the question, it is right. Every piece of code here runs as it is (the tests run it).
\answerday{1}{Your first chart and your first table}{answers:python:1}
\par\noindent{\sffamily\small\bfseries\color{tryframe}Task 1.} Make a bar chart of the \code{\textquotedbl{}weather\textquotedbl{}} table with \code{x=\textquotedbl{}city\textquotedbl{}} and \code{y=\textquotedbl{}rain\textquotedbl{}}, and save it as \code{\textquotedbl{}rain.\allowbreak{}png\textquotedbl{}}.\par
\begin{lstlisting}[style=py]
import simplydata as sdata
weather = sdata.example("weather")
chart = sdata.chart(weather)
chart.bars(x="city", y="rain")
chart.save("rain.png")
\end{lstlisting}
\begin{answerbox}One bar for each city, as tall as its average rain over the 12 months.\end{answerbox}
\par\noindent{\sffamily\small\bfseries\color{tryframe}Task 2.} Draw dots of \code{\textquotedbl{}hours\textquotedbl{}} (x) against \code{\textquotedbl{}marks\textquotedbl{}} (y) from the marks table, with a title, a name for each axis and a caption at the bottom.\par
\begin{lstlisting}[style=py]
import simplydata as sdata
marks = sdata.example("marks")
chart = sdata.chart(marks)
chart.dots(x="hours", y="marks")
chart.title_top("More study, better marks?")
chart.label_x("Hours of study")
chart.label_y("Mark (out of 100)")
chart.caption_bottom("Made-up data")
chart.save("hours_and_marks.png")
\end{lstlisting}
\begin{answerbox}First draw, then add the words; the order of the word steps does not matter.\end{answerbox}
\par\noindent{\sffamily\small\bfseries\color{tryframe}Task 3.} Start a table of the weather, keep the rows with \code{\textquotedbl{}temperature\textquotedbl{}} above 15, and give the average \code{\textquotedbl{}rain\textquotedbl{}} for each city.\par
\begin{lstlisting}[style=py]
import simplydata as sdata
table = sdata.table(sdata.example("weather"))
table.keep_rows("temperature", above=15)
table.average_of("rain", for_each="city")
\end{lstlisting}
\begin{answerbox}\code{keep\_\allowbreak{}rows} says how many warm months are left; \code{average\_\allowbreak{}of} gives one average for each city.\end{answerbox}
{\small\hyperref[day:python:1]{Back to day 1}}
\answerday{2}{Colours, word families and saving}{answers:python:2}
\begin{enumerate}
  \item Which word lists the practice tables? \fromday{1}\\ $\rightarrow$ \code{sdata.\allowbreak{}example()}; \code{sdata.\allowbreak{}example(\textquotedbl{}marks\textquotedbl{})} gets one of them.
  \item How do you start a chart of a table called marks? \fromday{1}\\ $\rightarrow$ chart = \code{sdata.\allowbreak{}chart(marks)}
  \item Which word draws one bar for each name, as tall as the marks? \fromday{1}\\ $\rightarrow$ \code{chart.\allowbreak{}bars(x=\textquotedbl{}name\textquotedbl{},\allowbreak{} y=\textquotedbl{}marks\textquotedbl{})}
\end{enumerate}
\par\noindent{\sffamily\small\bfseries\color{tryframe}Task 1.} Weather table: draw lines of \code{\textquotedbl{}temperature\textquotedbl{}} over \code{\textquotedbl{}month\textquotedbl{}}, one colour for each city. Put the legend on top and use the style \code{\textquotedbl{}paper\textquotedbl{}}.\par
\begin{lstlisting}[style=py]
import simplydata as sdata
weather = sdata.example("weather")
chart = sdata.chart(weather)
chart.lines(x="month", y="temperature", colour="city")
chart.legend_top()
chart.style("paper")
chart.save("temperature.png")
\end{lstlisting}
\begin{answerbox}\code{colour=\textquotedbl{}city\textquotedbl{}} draws one line for each city and adds a legend; the style changes the whole look.\end{answerbox}
\par\noindent{\sffamily\small\bfseries\color{tryframe}Task 2.} Save one chart as \code{.\allowbreak{}png}, \code{.\allowbreak{}pdf} and \code{.\allowbreak{}svg}. Then guess, before you look anything up: which word hides the legend, and which puts it on the left? Try your guesses on a chart with \code{colour=}.\par
\begin{lstlisting}[style=py]
import simplydata as sdata
marks = sdata.example("marks")
chart = sdata.chart(marks)
chart.bars(x="subject", y="marks", colour="class")
chart.save("subjects.png")
chart.save("subjects.pdf")
chart.save("subjects.svg")
chart.legend_left()
chart.hide_legend()
chart.save("no_legend.png")
\end{lstlisting}
\begin{answerbox}\code{legend\_\allowbreak{}left} follows \code{legend\_\allowbreak{}bottom} and \code{legend\_\allowbreak{}top}; \code{hide\_\allowbreak{}legend} follows \code{show\_\allowbreak{}} and \code{hide\_\allowbreak{}}. A family is learned once.\end{answerbox}
{\small\hyperref[day:python:2]{Back to day 2}}
\answerday{3}{Every kind of chart}{answers:python:3}
\begin{enumerate}
  \item How do you give each class its own colour in a bar chart? \fromday{2}\\ $\rightarrow$ \code{colour=\textquotedbl{}class\textquotedbl{}} inside the chart step: \code{chart.\allowbreak{}bars(x=\textquotedbl{}subject\textquotedbl{},\allowbreak{} y=\textquotedbl{}marks\textquotedbl{},\allowbreak{} colour=\textquotedbl{}class\textquotedbl{})}
  \item Which word moves the legend below the chart? \fromday{2}\\ $\rightarrow$ \code{chart.\allowbreak{}legend\_\allowbreak{}bottom()}
  \item How do you save a chart as a PDF? \fromday{1}\\ $\rightarrow$ \code{chart.\allowbreak{}save(\textquotedbl{}my\_\allowbreak{}chart.\allowbreak{}pdf\textquotedbl{})}: the end of the file name decides the type.
\end{enumerate}
\par\noindent{\sffamily\small\bfseries\color{tryframe}Task 1.} Draw violins of \code{\textquotedbl{}temperature\textquotedbl{}} for each city in the weather table, with all the dots on top.\par
\begin{lstlisting}[style=py]
import simplydata as sdata
weather = sdata.example("weather")
chart = sdata.chart(weather)
chart.violins(x="city", y="temperature")
chart.all_dots(x="city", y="temperature")
chart.save("violins.png")
\end{lstlisting}
\begin{answerbox}A violin shows the shape of the spread; the dots show every month.\end{answerbox}
{\small\hyperref[day:python:3]{Back to day 3}}
\answerday{4}{Statistics on charts, and axes}{answers:python:4}
\begin{enumerate}
  \item Which chart shows how the values of one column are spread, in bins? \fromday{3}\\ $\rightarrow$ \code{chart.\allowbreak{}histogram(x=\textquotedbl{}marks\textquotedbl{},\allowbreak{} bins=8)}
  \item Which word puts a title above the chart? \fromday{1}\\ $\rightarrow$ \code{chart.\allowbreak{}title\_\allowbreak{}top(\textquotedbl{}My title\textquotedbl{})}
  \item Which two words show the spread of each group, with every value as a dot on top? \fromday{3}\\ $\rightarrow$ \code{chart.\allowbreak{}boxes(x=\textquotedbl{}class\textquotedbl{},\allowbreak{} y=\textquotedbl{}marks\textquotedbl{})} and \code{chart.\allowbreak{}all\_\allowbreak{}dots(x=\textquotedbl{}class\textquotedbl{},\allowbreak{} y=\textquotedbl{}marks\textquotedbl{})}
\end{enumerate}
\par\noindent{\sffamily\small\bfseries\color{tryframe}Task 1.} Weather table: draw bars of the average \code{\textquotedbl{}rain\textquotedbl{}} for each city with 95 \% confidence intervals, and write the numbers on the bars.\par
\begin{lstlisting}[style=py]
import simplydata as sdata
weather = sdata.example("weather")
chart = sdata.chart(weather)
chart.bars(x="city", y="rain", error="confidence interval")
chart.numbers_on_bars()
chart.save("rain_ci.png")
\end{lstlisting}
\begin{answerbox}A line is printed saying which error bars were drawn, so the caption can say it too.\end{answerbox}
\par\noindent{\sffamily\small\bfseries\color{tryframe}Task 2.} Growth table: draw lines of \code{\textquotedbl{}population\textquotedbl{}} over \code{\textquotedbl{}year\textquotedbl{}}, one for each country, with a line at x = 2010 and the y-axis starting at 0.\par
\begin{lstlisting}[style=py]
import simplydata as sdata
growth = sdata.example("growth")
chart = sdata.chart(growth)
chart.lines(x="year", y="population", colour="country")
chart.line_at_x(2010, label="2010")
chart.range_y(0, 60)
chart.save("growth.png")
\end{lstlisting}
\begin{answerbox}\code{range\_\allowbreak{}y} takes the lowest and the highest value of the axis; \code{line\_\allowbreak{}at\_\allowbreak{}x} draws a line going up.\end{answerbox}
{\small\hyperref[day:python:4]{Back to day 4}}
\answerday{5}{Several charts, and a figure for a journal}{answers:python:5}
\begin{enumerate}
  \item Which option makes a bar show the median instead of the average? \fromday{4}\\ $\rightarrow$ \code{show=\textquotedbl{}median\textquotedbl{}}
  \item Which word lets you choose the colours of the groups? \fromday{2}\\ $\rightarrow$ \code{chart.\allowbreak{}colours(\textquotedbl{}orange\textquotedbl{},\allowbreak{} \textquotedbl{}teal\textquotedbl{},\allowbreak{} \textquotedbl{}purple\textquotedbl{})}
  \item Which option adds error bars that say how sure we are of the average? \fromday{4}\\ $\rightarrow$ \code{error=\textquotedbl{}standard error\textquotedbl{}}, or \code{error=\textquotedbl{}confidence interval\textquotedbl{}} for the 95 \% interval.
\end{enumerate}
\par\noindent{\sffamily\small\bfseries\color{tryframe}Task 1.} Make a histogram and a box chart of the marks, and put them one above the other on one page.\par
\begin{lstlisting}[style=py]
import simplydata as sdata
marks = sdata.example("marks")
a = sdata.chart(marks).histogram(x="marks")
b = sdata.chart(marks).boxes(x="class", y="marks")
page = sdata.one_above_other(a, b)
page.save("two_charts.png")
\end{lstlisting}
\begin{answerbox}Each chart is made first; \code{one\_\allowbreak{}above\_\allowbreak{}other} puts them on one page.\end{answerbox}
\par\noindent{\sffamily\small\bfseries\color{tryframe}Task 2.} A figure for a journal: dots of \code{\textquotedbl{}hours\textquotedbl{}} against \code{\textquotedbl{}marks\textquotedbl{}}, style \code{\textquotedbl{}paper\textquotedbl{}}, 8.5 cm wide and 7 cm high, font size 8, saved as a \code{.\allowbreak{}tif} with 600 dpi.\par
\begin{lstlisting}[style=py]
import simplydata as sdata
marks = sdata.example("marks")
chart = sdata.chart(marks)
chart.dots(x="hours", y="marks")
chart.style("paper")
chart.size(width=8.5, height=7, unit="cm")
chart.font_size(8)
chart.save("figure1.tif", resolution=600)
\end{lstlisting}
\begin{answerbox}Journals often ask for 8.5 cm (one column) or 17 cm (two columns) and 300 to 600 dpi.\end{answerbox}
{\small\hyperref[day:python:5]{Back to day 5}}
\answerday{6}{Open a table, choose rows and columns}{answers:python:6}
\begin{enumerate}
  \item Which word draws one small chart for each class? \fromday{5}\\ $\rightarrow$ \code{chart.\allowbreak{}one\_\allowbreak{}chart\_\allowbreak{}per(\textquotedbl{}class\textquotedbl{})}
  \item Which word draws one dot for each row, x against y? \fromday{3}\\ $\rightarrow$ \code{chart.\allowbreak{}dots(x=\textquotedbl{}hours\textquotedbl{},\allowbreak{} y=\textquotedbl{}marks\textquotedbl{})}
  \item Which word puts two charts next to each other on one page? \fromday{5}\\ $\rightarrow$ \code{sdata.\allowbreak{}side\_\allowbreak{}by\_\allowbreak{}side(a,\allowbreak{} b)}
\end{enumerate}
\par\noindent{\sffamily\small\bfseries\color{tryframe}Task 1.} Open the \code{\textquotedbl{}shop\textquotedbl{}} practice table. How many rows does it have, and how many rows does each shop have?\par
\begin{lstlisting}[style=py]
import simplydata as sdata
table = sdata.table(sdata.example("shop"))
table.count_rows()
table.count_each("shop")
table.describe_columns()
\end{lstlisting}
\begin{answerbox}\code{count\_\allowbreak{}rows} counts all the rows, \code{count\_\allowbreak{}each} the rows of each shop; \code{describe\_\allowbreak{}columns} also shows the empty cells.\end{answerbox}
\par\noindent{\sffamily\small\bfseries\color{tryframe}Task 2.} Weather table: keep the city \code{\textquotedbl{}Sunnyvale\textquotedbl{}}, keep only the columns month, temperature and rain, sort by rain with the biggest first, and save it as \code{\textquotedbl{}sunny.\allowbreak{}csv\textquotedbl{}}.\par
\begin{lstlisting}[style=py]
import simplydata as sdata
table = sdata.table(sdata.example("weather"))
table.keep_rows("city", equal_to="Sunnyvale")
table.keep_columns("month", "temperature", "rain")
table.sort_by("rain", biggest_first=True)
table.save("sunny.csv")
table.steps()
\end{lstlisting}
\begin{answerbox}Each step prints what it did; \code{steps()} lists them all, in order.\end{answerbox}
{\small\hyperref[day:python:6]{Back to day 6}}
\answerday{7}{Review of the first week}{answers:python:7}
\begin{enumerate}
  \item How do you look at the first 5 rows of a table? \fromday{6}\\ $\rightarrow$ \code{table.\allowbreak{}show\_\allowbreak{}first\_\allowbreak{}rows(5)}
  \item How do you name the y-axis? \fromday{1}\\ $\rightarrow$ \code{chart.\allowbreak{}label\_\allowbreak{}y(\textquotedbl{}Marks out of 100\textquotedbl{})}
  \item Which word draws a straight line through the dots? \fromday{4}\\ $\rightarrow$ \code{chart.\allowbreak{}fit\_\allowbreak{}line()}, after \code{chart.\allowbreak{}dots(.\allowbreak{}.\allowbreak{}.\allowbreak{})}
  \item How do you put the bars in your own order? \fromday{2}\\ $\rightarrow$ \code{chart.\allowbreak{}order\_\allowbreak{}x(\textquotedbl{}C\textquotedbl{},\allowbreak{} \textquotedbl{}B\textquotedbl{},\allowbreak{} \textquotedbl{}A\textquotedbl{})}
  \item How do you put four charts in two rows, with the letters A to D? \fromday{5}\\ $\rightarrow$ \code{sdata.\allowbreak{}grid(a,\allowbreak{} b,\allowbreak{} c,\allowbreak{} d,\allowbreak{} per\_\allowbreak{}row=2,\allowbreak{} letters=True)}
  \item Which word draws a line with a dot at each point? \fromday{3}\\ $\rightarrow$ \code{chart.\allowbreak{}lines(x=\textquotedbl{}month\textquotedbl{},\allowbreak{} y=\textquotedbl{}temperature\textquotedbl{},\allowbreak{} dots=True)}
\end{enumerate}
\par\noindent{\sffamily\small\bfseries\color{tryframe}Task 1.} Marks table: keep only class \code{\textquotedbl{}A\textquotedbl{}}, draw bars of the average \code{\textquotedbl{}marks\textquotedbl{}} for each subject with standard error bars, give the chart a title and save it as a PDF.\par
\begin{lstlisting}[style=py]
import simplydata as sdata
table = sdata.table(sdata.example("marks"))
table.keep_rows("class", equal_to="A")
chart = table.chart()
chart.bars(x="subject", y="marks", error="standard error")
chart.title_top("Class A")
chart.save("class_a.pdf")
\end{lstlisting}
\begin{answerbox}Day 1 (bars, save), day 4 (\code{error=}) and day 6 (\code{keep\_\allowbreak{}rows}) in one small program.\end{answerbox}
{\small\hyperref[day:python:7]{Back to day 7}}
\answerday{8}{Clean the data and add columns}{answers:python:8}
\begin{enumerate}
  \item Which word describes each column: its kind, empty cells and numbers? \fromday{6}\\ $\rightarrow$ \code{table.\allowbreak{}describe\_\allowbreak{}columns()}
  \item Which word adds a second y-axis on the right? \fromday{5}\\ $\rightarrow$ \code{chart.\allowbreak{}second\_\allowbreak{}y(\textquotedbl{}rain\textquotedbl{})}
  \item The names under the bars bump into each other. Which word helps? \fromday{1}\\ $\rightarrow$ \code{chart.\allowbreak{}rotate\_\allowbreak{}labels\_\allowbreak{}x(45)}
\end{enumerate}
\par\noindent{\sffamily\small\bfseries\color{tryframe}Task 1.} Shop table: count the empty cells, fill the empty \code{\textquotedbl{}sales\textquotedbl{}} with the line between neighbours for each shop (sorted by \code{\textquotedbl{}date\textquotedbl{}}), and add a column \code{\textquotedbl{}weekday\textquotedbl{}} with the name of the weekday.\par
\begin{lstlisting}[style=py]
import simplydata as sdata
table = sdata.table(sdata.example("shop"))
table.count_empty_cells()
table.fill_empty_cells("sales", with_the="line between neighbours", sorted_by="date", for_each="shop")
table.add_date_part_column("weekday", from_column="date", part="weekday name")
table.show_first_rows()
\end{lstlisting}
\begin{answerbox}\code{for\_\allowbreak{}each=} keeps the two shops apart, so a gap is never filled from the other shop's days.\end{answerbox}
{\small\hyperref[day:python:8]{Back to day 8}}
\answerday{9}{Summaries, numbers and grids}{answers:python:9}
\begin{enumerate}
  \item Which word counts the empty cells in each column? \fromday{8}\\ $\rightarrow$ \code{table.\allowbreak{}count\_\allowbreak{}empty\_\allowbreak{}cells()}
  \item Which word counts the rows of each fruit? \fromday{6}\\ $\rightarrow$ \code{table.\allowbreak{}count\_\allowbreak{}each(\textquotedbl{}fruit\textquotedbl{})}
  \item If \code{keep\_\allowbreak{}rows} keeps rows, what would you guess takes columns away? \fromday{2}\\ $\rightarrow$ \code{remove\_\allowbreak{}columns}: verb + thing. \code{keep\_\allowbreak{}} and \code{remove\_\allowbreak{}} work on rows, columns and numbers.
\end{enumerate}
\par\noindent{\sffamily\small\bfseries\color{tryframe}Task 1.} Fruit table: the total \code{\textquotedbl{}count\textquotedbl{}} of each fruit, drawn as bars with the numbers on the bars.\par
\begin{lstlisting}[style=py]
import simplydata as sdata
table = sdata.table(sdata.example("fruit"))
totals = table.total_of("count", for_each="fruit")
chart = totals.chart()
chart.bars(x="fruit", y="total of count")
chart.numbers_on_bars()
chart.save("fruit_totals.png")
\end{lstlisting}
\begin{answerbox}A summary gives back a new table, with a column called \code{\textquotedbl{}total of count\textquotedbl{}}, ready for a chart.\end{answerbox}
\par\noindent{\sffamily\small\bfseries\color{tryframe}Task 2.} Make 200 random heights with average 150 and standard deviation 8 (\code{seed=1}). What is their 90th percentile? Draw them as a histogram.\par
\begin{lstlisting}[style=py]
import simplydata as sdata
heights = sdata.random_numbers(how_many=200, average=150, standard_deviation=8, seed=1)
heights.percentile(90)
chart = heights.chart()
chart.histogram(x="value")
chart.save("heights.png")
\end{lstlisting}
\begin{answerbox}About 157: nine heights in ten are below it. \code{seed=1} gives the same numbers every time.\end{answerbox}
{\small\hyperref[day:python:9]{Back to day 9}}
\answerday{10}{Is a difference real? Tests and effect sizes}{answers:python:10}
\begin{enumerate}
  \item Which word gives the average of a column for each group? \fromday{9}\\ $\rightarrow$ \code{table.\allowbreak{}average\_\allowbreak{}of(\textquotedbl{}marks\textquotedbl{},\allowbreak{} for\_\allowbreak{}each=\textquotedbl{}class\textquotedbl{})}
  \item Which word counts the rows? \fromday{6}\\ $\rightarrow$ \code{table.\allowbreak{}count\_\allowbreak{}rows()}
  \item Which chart shows how the values of one column are spread, in bins? \fromday{3}\\ $\rightarrow$ \code{chart.\allowbreak{}histogram(x=\textquotedbl{}marks\textquotedbl{},\allowbreak{} bins=8)}
\end{enumerate}
\par\noindent{\sffamily\small\bfseries\color{tryframe}Task 1.} School table: do the average marks differ between pupils who had breakfast and pupils who did not? Write down the difference, the effect size and the p-value.\par
\begin{lstlisting}[style=py]
import simplydata as sdata
table = sdata.table(sdata.example("school"))
table.average_of("marks", for_each="breakfast")
table.compare_averages("marks", for_each="breakfast")
\end{lstlisting}
\begin{answerbox}Two groups, so it is Welch's t-test; the sentence gives the difference, Hedges' g and the p-value.\end{answerbox}
{\small\hyperref[day:python:10]{Back to day 10}}
\answerday{11}{Counts, outliers and many tests}{answers:python:11}
\begin{enumerate}
  \item Which word tests whether the average marks differ between the classes? \fromday{10}\\ $\rightarrow$ \code{table.\allowbreak{}compare\_\allowbreak{}averages(\textquotedbl{}marks\textquotedbl{},\allowbreak{} for\_\allowbreak{}each=\textquotedbl{}class\textquotedbl{})}
  \item How do you fill the empty marks with the median? \fromday{8}\\ $\rightarrow$ \code{table.\allowbreak{}fill\_\allowbreak{}empty\_\allowbreak{}cells(\textquotedbl{}marks\textquotedbl{},\allowbreak{} with\_\allowbreak{}the=\textquotedbl{}median\textquotedbl{})}
  \item Which word checks by eye whether values are bell-shaped? \fromday{4}\\ $\rightarrow$ \code{chart.\allowbreak{}qq\_\allowbreak{}dots(\textquotedbl{}marks\textquotedbl{})}
\end{enumerate}
\par\noindent{\sffamily\small\bfseries\color{tryframe}Task 1.} School table: is the share of pupils who \code{\textquotedbl{}passed\textquotedbl{}} different with and without breakfast?\par
\begin{lstlisting}[style=py]
import simplydata as sdata
table = sdata.table(sdata.example("school"))
table.count_rows(for_each="breakfast", across="passed")
table.compare_counts("passed", for_each="breakfast")
\end{lstlisting}
\begin{answerbox}passed is a category (yes or no), so the test compares counts: the chi-square test, with Cramer's V as the effect size.\end{answerbox}
{\small\hyperref[day:python:11]{Back to day 11}}
\answerday{12}{Choosing a test, and writing it up}{answers:python:12}
\begin{enumerate}
  \item Which word tests whether the counts of categories differ between groups? \fromday{11}\\ $\rightarrow$ \code{table.\allowbreak{}compare\_\allowbreak{}counts(\textquotedbl{}result\textquotedbl{},\allowbreak{} for\_\allowbreak{}each=\textquotedbl{}class\textquotedbl{})}
  \item Which word tells whether two columns go up and down together? \fromday{9}\\ $\rightarrow$ \code{table.\allowbreak{}correlation\_\allowbreak{}between(\textquotedbl{}hours\textquotedbl{},\allowbreak{} \textquotedbl{}marks\textquotedbl{})}
  \item Which word sets the size of a figure in cm? \fromday{5}\\ $\rightarrow$ \code{chart.\allowbreak{}size(width=17,\allowbreak{} height=8,\allowbreak{} unit=\textquotedbl{}cm\textquotedbl{})}
\end{enumerate}
\par\noindent{\sffamily\small\bfseries\color{tryframe}Task 1.} Marks table: is each child's Maths mark different from the same child's Science mark? Choose the right word first.\par
\begin{lstlisting}[style=py]
import simplydata as sdata
table = sdata.table(sdata.example("marks"))
table.keep_columns("name", "subject", "marks")
table.turn_rows_into_columns(names_from="subject", values_from="marks")
table.compare_before_after("Maths", "Science")
\end{lstlisting}
\begin{answerbox}The same children twice, so the test must be paired. First one row per child, with a Maths and a Science column.\end{answerbox}
\par\noindent{\sffamily\small\bfseries\color{tryframe}Task 2.} School table: is the average \code{\textquotedbl{}sleep\textquotedbl{}} different from 8 hours?\par
\begin{lstlisting}[style=py]
import simplydata as sdata
table = sdata.table(sdata.example("school"))
table.compare_average_to("sleep", 8)
\end{lstlisting}
\begin{answerbox}One average against a number you choose: the one-sample t-test. Empty sleep cells are skipped, with a note.\end{answerbox}
{\small\hyperref[day:python:12]{Back to day 12}}
\answerday{13}{A whole analysis}{answers:python:13}
\begin{enumerate}
  \item The same children were measured twice. Which word compares the two? \fromday{12}\\ $\rightarrow$ \code{table.\allowbreak{}compare\_\allowbreak{}before\_\allowbreak{}after(\textquotedbl{}before\textquotedbl{},\allowbreak{} \textquotedbl{}after\textquotedbl{})}: the paired test.
  \item Which word says which pairs of groups differ? \fromday{10}\\ $\rightarrow$ \code{table.\allowbreak{}compare\_\allowbreak{}each\_\allowbreak{}pair(\textquotedbl{}marks\textquotedbl{},\allowbreak{} for\_\allowbreak{}each=\textquotedbl{}class\textquotedbl{})}
  \item Which word keeps only the rows where subject is Maths? \fromday{6}\\ $\rightarrow$ \code{table.\allowbreak{}keep\_\allowbreak{}rows(\textquotedbl{}subject\textquotedbl{},\allowbreak{} equal\_\allowbreak{}to=\textquotedbl{}Maths\textquotedbl{})}
\end{enumerate}
\par\noindent{\sffamily\small\bfseries\color{tryframe}Task 1.} From memory: weather table, compare the average \code{\textquotedbl{}temperature\textquotedbl{}} of the cities, draw boxes with all the dots, and save the figure as a PDF and the table as \code{\textquotedbl{}weather.\allowbreak{}xlsx\textquotedbl{}}.\par
\begin{lstlisting}[style=py]
import simplydata as sdata
table = sdata.table(sdata.example("weather"))
table.compare_averages("temperature", for_each="city")
chart = table.chart()
chart.boxes(x="city", y="temperature")
chart.all_dots(x="city", y="temperature")
chart.save("temperature.pdf")
table.save("weather.xlsx")
\end{lstlisting}
\begin{answerbox}Question, test, figure, saved table: the whole analysis in eight lines.\end{answerbox}
{\small\hyperref[day:python:13]{Back to day 13}}
\answerday{14}{Review of the second week}{answers:python:14}
\begin{enumerate}
  \item Which word lists every change you made to a table, for your methods? \fromday{13}\\ $\rightarrow$ \code{table.\allowbreak{}steps()}
  \item Which word makes all the text on a chart bigger or smaller? \fromday{1}\\ $\rightarrow$ \code{chart.\allowbreak{}font\_\allowbreak{}size(14)}
  \item A word for one axis ends in ...? \fromday{2}\\ $\rightarrow$ \_x or \_y: \code{label\_\allowbreak{}y}, \code{range\_\allowbreak{}x}, \code{log\_\allowbreak{}scale\_\allowbreak{}y}.
  \item Which word turns a column of text such as \code{\textquotedbl{}144\textquotedbl{}} into numbers? \fromday{8}\\ $\rightarrow$ \code{table.\allowbreak{}make\_\allowbreak{}numbers(\textquotedbl{}height\textquotedbl{})}
  \item Which word sets the y-axis from 0 to 100? \fromday{4}\\ $\rightarrow$ \code{chart.\allowbreak{}range\_\allowbreak{}y(0,\allowbreak{} 100)}
  \item Which word shows the rows with outliers? \fromday{11}\\ $\rightarrow$ \code{table.\allowbreak{}show\_\allowbreak{}rows\_\allowbreak{}with\_\allowbreak{}outliers(\textquotedbl{}marks\textquotedbl{},\allowbreak{} for\_\allowbreak{}each=\textquotedbl{}class\textquotedbl{})}
  \item Which word joins two tables on a column they share? \fromday{9}\\ $\rightarrow$ \code{table.\allowbreak{}join\_\allowbreak{}with(ages,\allowbreak{} using=\textquotedbl{}name\textquotedbl{})}
  \item How do you save a picture sharp enough for a journal? \fromday{5}\\ $\rightarrow$ \code{chart.\allowbreak{}save(\textquotedbl{}figure.\allowbreak{}tif\textquotedbl{},\allowbreak{} resolution=600)}, or save it as \code{.\allowbreak{}pdf}
\end{enumerate}
\par\noindent{\sffamily\small\bfseries\color{tryframe}Task 1.} School table: count the empty cells, remove the rows with an empty \code{\textquotedbl{}sleep\textquotedbl{}}, give the average sleep for each value of \code{\textquotedbl{}passed\textquotedbl{}}, test the difference, and draw boxes with all the dots.\par
\begin{lstlisting}[style=py]
import simplydata as sdata
table = sdata.table(sdata.example("school"))
table.count_empty_cells()
table.remove_rows_with_empty_cells("sleep")
table.average_of("sleep", for_each="passed")
table.compare_averages("sleep", for_each="passed")
chart = table.chart()
chart.boxes(x="passed", y="sleep")
chart.all_dots(x="passed", y="sleep")
chart.save("sleep.png")
\end{lstlisting}
\begin{answerbox}Clean, sum up, test, draw: the order of every analysis.\end{answerbox}
{\small\hyperref[day:python:14]{Back to day 14}}
\answerday{15}{Models: explain one column from several}{answers:python:15}
\begin{enumerate}
  \item How do you give the panels of a figure the letters A and B? \fromday{13}\\ $\rightarrow$ \code{letters=True} in \code{sdata.\allowbreak{}grid(.\allowbreak{}.\allowbreak{}.\allowbreak{})} or \code{sdata.\allowbreak{}side\_\allowbreak{}by\_\allowbreak{}side(.\allowbreak{}.\allowbreak{}.\allowbreak{})}
  \item Which word compares one average with a number you choose? \fromday{12}\\ $\rightarrow$ \code{table.\allowbreak{}compare\_\allowbreak{}average\_\allowbreak{}to(\textquotedbl{}marks\textquotedbl{},\allowbreak{} 60)}
  \item Which word adds a new column made from other columns? \fromday{8}\\ $\rightarrow$ \code{table.\allowbreak{}add\_\allowbreak{}column(\textquotedbl{}total\textquotedbl{},\allowbreak{} \textquotedbl{}Maths + Science\textquotedbl{})}, or \code{table.\allowbreak{}add\_\allowbreak{}column(\textquotedbl{}per\_\allowbreak{}hour\textquotedbl{},\allowbreak{} by\_\allowbreak{}dividing=[\textquotedbl{}marks\textquotedbl{},\allowbreak{} \textquotedbl{}hours\textquotedbl{}])}
  \item How do you start a table from the marks practice table? \fromday{1}\\ $\rightarrow$ table = \code{sdata.\allowbreak{}table(sdata.\allowbreak{}example(\textquotedbl{}marks\textquotedbl{}))}
\end{enumerate}
\par\noindent{\sffamily\small\bfseries\color{tryframe}Task 1.} School table: explain \code{\textquotedbl{}days\_\allowbreak{}absent\textquotedbl{}} using \code{\textquotedbl{}sleep\textquotedbl{}} and \code{\textquotedbl{}breakfast\textquotedbl{}} with \code{outcome=\textquotedbl{}counts\textquotedbl{}}. Does breakfast matter?\par
\begin{lstlisting}[style=py]
import simplydata as sdata
table = sdata.table(sdata.example("school"))
table.explain("days_absent", using=["sleep", "breakfast"], outcome="counts")
\end{lstlisting}
\begin{answerbox}Counts of days call for Poisson regression; read the line for breakfast and its p-value.\end{answerbox}
{\small\hyperref[day:python:15]{Back to day 15}}
\answerday{16}{Patterns across many columns}{answers:python:16}
\begin{enumerate}
  \item Which word explains marks from hours, sleep and breakfast together? \fromday{15}\\ $\rightarrow$ \code{table.\allowbreak{}explain(\textquotedbl{}marks\textquotedbl{},\allowbreak{} using=[\textquotedbl{}hours\textquotedbl{},\allowbreak{} \textquotedbl{}sleep\textquotedbl{},\allowbreak{} \textquotedbl{}breakfast\textquotedbl{}])}
  \item Which word hands a table on to pandas? \fromday{13}\\ $\rightarrow$ \code{table.\allowbreak{}as\_\allowbreak{}pandas()}
  \item Which word gives several summaries at once? \fromday{9}\\ $\rightarrow$ \code{table.\allowbreak{}summary\_\allowbreak{}of(\textquotedbl{}sales\textquotedbl{},\allowbreak{} for\_\allowbreak{}each=\textquotedbl{}shop\textquotedbl{},\allowbreak{} show=[\textquotedbl{}average\textquotedbl{},\allowbreak{} \textquotedbl{}median\textquotedbl{},\allowbreak{} \textquotedbl{}count\textquotedbl{}])}
  \item Which file endings give a picture that stays sharp at any size? \fromday{2}\\ $\rightarrow$ \code{.\allowbreak{}pdf} and \code{.\allowbreak{}svg}
\end{enumerate}
\par\noindent{\sffamily\small\bfseries\color{tryframe}Task 1.} Meadow table: add 3 clusters and count the rows of each cluster across \code{\textquotedbl{}field\textquotedbl{}}.\par
\begin{lstlisting}[style=py]
import simplydata as sdata
table = sdata.table(sdata.example("meadow"))
table.add_clusters(how_many=3)
table.count_rows(for_each="cluster", across="field")
\end{lstlisting}
\begin{answerbox}If each cluster is mostly one field, the fields have different plants.\end{answerbox}
{\small\hyperref[day:python:16]{Back to day 16}}
\answerday{17}{Survival, and planning a study}{answers:python:17}
\begin{enumerate}
  \item Which word adds the principal components PC1 and PC2? \fromday{16}\\ $\rightarrow$ \code{table.\allowbreak{}add\_\allowbreak{}principal\_\allowbreak{}components(how\_\allowbreak{}many=2)}
  \item Which word lists every change you made to a table, for your methods? \fromday{13}\\ $\rightarrow$ \code{table.\allowbreak{}steps()}
  \item A p-value answers \code{\textquotedbl{}could it be chance?\textquotedbl{}}. What answers \code{\textquotedbl{}how big is it?\textquotedbl{}} \fromday{10}\\ $\rightarrow$ The effect size: the difference in averages, Cohen's d or Hedges' g, omega-squared.
  \item Which two words show the spread of each group, with every value as a dot on top? \fromday{3}\\ $\rightarrow$ \code{chart.\allowbreak{}boxes(x=\textquotedbl{}class\textquotedbl{},\allowbreak{} y=\textquotedbl{}marks\textquotedbl{})} and \code{chart.\allowbreak{}all\_\allowbreak{}dots(x=\textquotedbl{}class\textquotedbl{},\allowbreak{} y=\textquotedbl{}marks\textquotedbl{})}
\end{enumerate}
\par\noindent{\sffamily\small\bfseries\color{tryframe}Task 1.} How many values does each group need to find a large difference (\code{effect\_\allowbreak{}size=0.\allowbreak{}8}) with power 0.8? And with a small one (0.2)?\par
\begin{lstlisting}[style=py]
import simplydata as sdata
sdata.how_many_needed(effect_size=0.8, power=0.8)
sdata.how_many_needed(effect_size=0.2, power=0.8)
\end{lstlisting}
\begin{answerbox}About 26 for a large difference but about 394 for a small one: small effects need big studies.\end{answerbox}
{\small\hyperref[day:python:17]{Back to day 17}}
\answerday{18}{Fitting lines, curves and formulas}{answers:python:18}
\begin{enumerate}
  \item Which word draws survival curves? \fromday{17}\\ $\rightarrow$ \code{chart.\allowbreak{}survival\_\allowbreak{}curves(time=\textquotedbl{}days\textquotedbl{},\allowbreak{} event=\textquotedbl{}died\textquotedbl{},\allowbreak{} colour=\textquotedbl{}treatment\textquotedbl{})}
  \item How do you explain a yes-or-no column such as \code{\textquotedbl{}passed\textquotedbl{}}? \fromday{15}\\ $\rightarrow$ \code{outcome=\textquotedbl{}yes or no\textquotedbl{}} (logistic regression)
  \item Which word checks whether values are bell-shaped (normal)? \fromday{11}\\ $\rightarrow$ \code{table.\allowbreak{}compare\_\allowbreak{}shape\_\allowbreak{}to(\textquotedbl{}marks\textquotedbl{},\allowbreak{} \textquotedbl{}normal\textquotedbl{})}
  \item Which word draws a line across at y = 50, with a label? \fromday{4}\\ $\rightarrow$ \code{chart.\allowbreak{}line\_\allowbreak{}at\_\allowbreak{}y(50,\allowbreak{} label=\textquotedbl{}pass mark\textquotedbl{})}
\end{enumerate}
\par\noindent{\sffamily\small\bfseries\color{tryframe}Task 1.} Growth table: fit the formula \code{\textquotedbl{}forest = a + b * (year - 2000)\textquotedbl{}} for each country. Which country loses forest fastest?\par
\begin{lstlisting}[style=py]
import simplydata as sdata
table = sdata.table(sdata.example("growth"))
table.fit_formula("forest = a + b * (year - 2000)", for_each="country")
\end{lstlisting}
\begin{answerbox}b is the change in forest each year: the country with the most negative b loses forest fastest.\end{answerbox}
{\small\hyperref[day:python:18]{Back to day 18}}
\answerday{19}{Finding, solving and distributions}{answers:python:19}
\begin{enumerate}
  \item Which word fits a straight line on a table? \fromday{18}\\ $\rightarrow$ \code{table.\allowbreak{}fit\_\allowbreak{}line(x=\textquotedbl{}hours\textquotedbl{},\allowbreak{} y=\textquotedbl{}marks\textquotedbl{})}
  \item Which word puts rows that are alike into groups? \fromday{16}\\ $\rightarrow$ \code{table.\allowbreak{}add\_\allowbreak{}clusters(how\_\allowbreak{}many=2)}
  \item What should a results sentence give? \fromday{12}\\ $\rightarrow$ The summary of each group with n, the effect size (with its interval), the name of the test and the exact p-value.
  \item Which word writes a note with an arrow on the chart? \fromday{5}\\ $\rightarrow$ \code{chart.\allowbreak{}note(\textquotedbl{}X marks the spot\textquotedbl{},\allowbreak{} x=10,\allowbreak{} y=4,\allowbreak{} arrow=True)}
\end{enumerate}
\par\noindent{\sffamily\small\bfseries\color{tryframe}Task 1.} A box has a square bottom of side x and a height of 10 - x, so its volume is x * x * (10 - x). For which x between 0 and 10 is the volume biggest?\par
\begin{lstlisting}[style=py]
import simplydata as sdata
sdata.find_biggest("x * x * (10 - x)", between=[0, 10])
\end{lstlisting}
\begin{answerbox}x = 6.67 (that is, 20/3), with a volume of about 148.\end{answerbox}
\par\noindent{\sffamily\small\bfseries\color{tryframe}Task 2.} Heights are normal with average 150 and standard deviation 8. What share is taller than 165? Which height do 95 \% of them stay below?\par
\begin{lstlisting}[style=py]
import simplydata as sdata
heights = sdata.distribution("normal", average=150, standard_deviation=8)
heights.chance(above=165)
heights.percentile(95)
\end{lstlisting}
\begin{answerbox}About 3 \% are taller than 165, and 95 \% are below about 163.\end{answerbox}
{\small\hyperref[day:python:19]{Back to day 19}}
\answerday{20}{Values between points}{answers:python:20}
\begin{enumerate}
  \item Which word finds where a formula is smallest? \fromday{19}\\ $\rightarrow$ \code{sdata.\allowbreak{}find\_\allowbreak{}smallest(\textquotedbl{}x**2 - 4*x + 1\textquotedbl{},\allowbreak{} between=[0,\allowbreak{} 5])}
  \item Which word tests whether survival differs between groups? \fromday{17}\\ $\rightarrow$ \code{table.\allowbreak{}compare\_\allowbreak{}survival(time=\textquotedbl{}days\textquotedbl{},\allowbreak{} event=\textquotedbl{}died\textquotedbl{},\allowbreak{} for\_\allowbreak{}each=\textquotedbl{}treatment\textquotedbl{})}
  \item How do you give the panels of a figure the letters A and B? \fromday{13}\\ $\rightarrow$ \code{letters=True} in \code{sdata.\allowbreak{}grid(.\allowbreak{}.\allowbreak{}.\allowbreak{})} or \code{sdata.\allowbreak{}side\_\allowbreak{}by\_\allowbreak{}side(.\allowbreak{}.\allowbreak{}.\allowbreak{})}
  \item Which word sorts with the biggest marks first? \fromday{6}\\ $\rightarrow$ \code{table.\allowbreak{}sort\_\allowbreak{}by(\textquotedbl{}marks\textquotedbl{},\allowbreak{} biggest\_\allowbreak{}first=True)}
\end{enumerate}
\par\noindent{\sffamily\small\bfseries\color{tryframe}Task 1.} Weather table: estimate the \code{\textquotedbl{}rain\textquotedbl{}} of each city at \code{month\_\allowbreak{}number} 2.5 and 7.5, joined by straight lines.\par
\begin{lstlisting}[style=py]
import simplydata as sdata
table = sdata.table(sdata.example("weather"))
rain = table.values_at("rain", at={"month_number": [2.5, 7.5]}, joined_by="straight lines", for_each="city")
rain.show()
\end{lstlisting}
\begin{answerbox}Each estimate lies halfway between the two months around it.\end{answerbox}
{\small\hyperref[day:python:20]{Back to day 20}}
\answerday{21}{Review, and the question of your project}{answers:python:21}
\begin{enumerate}
  \item Which word estimates a column between the points you measured? \fromday{20}\\ $\rightarrow$ \code{table.\allowbreak{}values\_\allowbreak{}at(\textquotedbl{}population\textquotedbl{},\allowbreak{} at=\{\textquotedbl{}year\textquotedbl{}: [2001,\allowbreak{} 2011]\})}
  \item Which word cuts a column of numbers into named groups? \fromday{8}\\ $\rightarrow$ \code{table.\allowbreak{}add\_\allowbreak{}group\_\allowbreak{}column(\textquotedbl{}grade\textquotedbl{},\allowbreak{} from\_\allowbreak{}column=\textquotedbl{}marks\textquotedbl{},\allowbreak{} cut\_\allowbreak{}at=[40,\allowbreak{} 60,\allowbreak{} 80],\allowbreak{} names=[\textquotedbl{}D\textquotedbl{},\allowbreak{} \textquotedbl{}C\textquotedbl{},\allowbreak{} \textquotedbl{}B\textquotedbl{},\allowbreak{} \textquotedbl{}A\textquotedbl{}])}
  \item Which word keeps only the rows with marks above 50? \fromday{1}\\ $\rightarrow$ \code{table.\allowbreak{}keep\_\allowbreak{}rows(\textquotedbl{}marks\textquotedbl{},\allowbreak{} above=50)}
  \item How do you make simplydata numbers from a list? \fromday{9}\\ $\rightarrow$ nums = \code{sdata.\allowbreak{}numbers([12,\allowbreak{} 15,\allowbreak{} 9,\allowbreak{} 20,\allowbreak{} 14])}
  \item Which option gives the rank-based test, for small groups that are not bell-shaped? \fromday{10}\\ $\rightarrow$ \code{based\_\allowbreak{}on=\textquotedbl{}ranks\textquotedbl{}}
  \item Which word adds the values a model predicts as a new column? \fromday{15}\\ $\rightarrow$ \code{table.\allowbreak{}add\_\allowbreak{}predicted\_\allowbreak{}column(\textquotedbl{}marks\textquotedbl{},\allowbreak{} using=[\textquotedbl{}hours\textquotedbl{},\allowbreak{} \textquotedbl{}breakfast\textquotedbl{}])}
  \item You did many tests at once. Which word corrects the p-values? \fromday{11}\\ $\rightarrow$ \code{add\_\allowbreak{}adjusted\_\allowbreak{}p\_\allowbreak{}value\_\allowbreak{}column(method=\textquotedbl{}holm\textquotedbl{})}, on the table of results
  \item A few values are much bigger than the rest. Which word helps? \fromday{4}\\ $\rightarrow$ \code{chart.\allowbreak{}log\_\allowbreak{}scale\_\allowbreak{}y()}
  \item How do you fit a curve with x and x squared? \fromday{18}\\ $\rightarrow$ \code{table.\allowbreak{}fit\_\allowbreak{}curve(x=\textquotedbl{}hours\textquotedbl{},\allowbreak{} y=\textquotedbl{}marks\textquotedbl{},\allowbreak{} degree=2)}
\end{enumerate}
\par\noindent{\sffamily\small\bfseries\color{tryframe}Task 1.} Trial table: test whether survival differs between the treatments, then explain the risk of dying from \code{\textquotedbl{}treatment\textquotedbl{}} and \code{\textquotedbl{}age\textquotedbl{}} together.\par
\begin{lstlisting}[style=py]
import simplydata as sdata
table = sdata.table(sdata.example("trial"))
table.compare_survival(time="days", event="died", for_each="treatment")
table.explain_survival(time="days", event="died", using=["treatment", "age"])
\end{lstlisting}
\begin{answerbox}The log-rank test compares the curves; the Cox model gives a hazard ratio for each column.\end{answerbox}
\par\noindent{\sffamily\small\bfseries\color{tryframe}Task 2.} Write your question in one sentence. Open your raw file and write down how many rows and columns it has, and which columns have empty cells.\par
\begin{lstlisting}[style=py]
import simplydata as sdata
sdata.table(sdata.example("school")).save("school_raw.csv")
table = sdata.table("school_raw.csv")
table.count_rows()
table.describe_columns()
\end{lstlisting}
\begin{answerbox}The first line only makes a raw file to practise on; with your own data, start at the second line. The school table has 180 rows and 8 columns; sleep has 19 empty cells.\end{answerbox}
{\small\hyperref[day:python:21]{Back to day 21}}
\answerday{22}{Project: clean}{answers:python:22}
\begin{enumerate}
  \item Which word tests whether groups differ in their make-up across many columns (PERMANOVA)? \fromday{16}\\ $\rightarrow$ \code{table.\allowbreak{}compare\_\allowbreak{}profiles(for\_\allowbreak{}each=\textquotedbl{}grazed\textquotedbl{})}
  \item Which word finds where the two sides of an equation are equal? \fromday{19}\\ $\rightarrow$ \code{sdata.\allowbreak{}solve(\textquotedbl{}x**3 - 2*x = 5\textquotedbl{},\allowbreak{} between=[0,\allowbreak{} 3])}
  \item Which word says how many values you need in each group? \fromday{17}\\ $\rightarrow$ \code{sdata.\allowbreak{}how\_\allowbreak{}many\_\allowbreak{}needed(effect\_\allowbreak{}size=0.\allowbreak{}5,\allowbreak{} power=0.\allowbreak{}8)}
  \item You made a mistake in the last step. Which word takes it back? \fromday{6}\\ $\rightarrow$ \code{table.\allowbreak{}undo()}; \code{table.\allowbreak{}start\_\allowbreak{}again()} goes back to the start.
\end{enumerate}
\par\noindent{\sffamily\small\bfseries\color{tryframe}Task 1.} Clean your table and save it as a new file. Write one line for each decision you made about empty cells.\par
\begin{lstlisting}[style=py]
import simplydata as sdata
table = sdata.table(sdata.example("school"))
table.count_empty_cells()
table.add_is_empty_column("sleep_was_empty", from_column="sleep")
table.rename_column("days_absent", to="absent")
table.save("school_clean.csv")
table.steps()
\end{lstlisting}
\begin{answerbox}Here the empty sleep cells are kept (tests skip them, and models can use multiple imputation); a column remembers which were empty.\end{answerbox}
{\small\hyperref[day:python:22]{Back to day 22}}
\answerday{23}{Project: explore}{answers:python:23}
\begin{enumerate}
  \item Which word fills an even map from scattered points, ready for a contour map? \fromday{20}\\ $\rightarrow$ \code{table.\allowbreak{}map\_\allowbreak{}of(\textquotedbl{}height\textquotedbl{},\allowbreak{} across=[\textquotedbl{}east\textquotedbl{},\allowbreak{} \textquotedbl{}north\textquotedbl{}])}
  \item Which word fits a formula of your own, such as growth? \fromday{18}\\ $\rightarrow$ \code{table.\allowbreak{}fit\_\allowbreak{}formula(\textquotedbl{}population = a * exp(b * (year - 2000))\textquotedbl{})}
  \item Which word gives the middle value of numbers? \fromday{9}\\ $\rightarrow$ \code{nums.\allowbreak{}median()}
  \item Which word shows the numbers on the y-axis as percents? \fromday{4}\\ $\rightarrow$ \code{chart.\allowbreak{}tick\_\allowbreak{}format\_\allowbreak{}y(\textquotedbl{}percent\textquotedbl{})}
\end{enumerate}
\par\noindent{\sffamily\small\bfseries\color{tryframe}Task 1.} Make three summaries and two charts of your table that help to answer your question.\par
\begin{lstlisting}[style=py]
import simplydata as sdata
table = sdata.table(sdata.example("school"))
table.count_each("breakfast")
table.summary_of("marks", for_each="breakfast", show=["average", "median", "standard deviation", "count"])
table.correlation_between("sleep", "marks")
chart = table.chart()
chart.histogram(x="marks", colour="breakfast")
chart.save("marks_histogram.png")
sdata.pairs(table, colour="breakfast").save("pairs.png")
\end{lstlisting}
\begin{answerbox}Each summary and chart is a small answer; the test comes later.\end{answerbox}
{\small\hyperref[day:python:23]{Back to day 23}}
\answerday{24}{Project: figures}{answers:python:24}
\begin{enumerate}
  \item Which word gives the average mark of each class? \fromday{1}\\ $\rightarrow$ \code{table.\allowbreak{}average\_\allowbreak{}of(\textquotedbl{}marks\textquotedbl{},\allowbreak{} for\_\allowbreak{}each=\textquotedbl{}class\textquotedbl{})}
  \item How do you make a normal distribution with average 50 and standard deviation 10? \fromday{19}\\ $\rightarrow$ \code{sdata.\allowbreak{}distribution(\textquotedbl{}normal\textquotedbl{},\allowbreak{} average=50,\allowbreak{} standard\_\allowbreak{}deviation=10)}
  \item Which word keeps rows 10 to 20, by their place in the table? \fromday{6}\\ $\rightarrow$ \code{table.\allowbreak{}keep\_\allowbreak{}rows\_\allowbreak{}by\_\allowbreak{}position(10,\allowbreak{} to=20)}
  \item How do you make 100 random numbers with average 50 and standard deviation 10, the same every time? \fromday{9}\\ $\rightarrow$ \code{sdata.\allowbreak{}random\_\allowbreak{}numbers(how\_\allowbreak{}many=100,\allowbreak{} average=50,\allowbreak{} standard\_\allowbreak{}deviation=10,\allowbreak{} seed=1)}
\end{enumerate}
\par\noindent{\sffamily\small\bfseries\color{tryframe}Task 1.} Make your main figure with two panels and save it ready for a journal.\par
\begin{lstlisting}[style=py]
import simplydata as sdata
table = sdata.table(sdata.example("school"))
a = table.chart()
a.boxes(x="breakfast", y="marks")
a.all_dots(x="breakfast", y="marks")
a.label_y("Mark (out of 100)")
b = table.chart()
b.dots(x="sleep", y="marks", colour="breakfast")
b.fit_line(confidence_band=True)
b.label_x("Sleep (hours)")
figure = sdata.grid(a, b, per_row=2, letters=True)
figure.style("paper")
figure.size(width=17, height=8, unit="cm")
figure.save("figure1.pdf")
figure.save("figure1.tif", resolution=600)
\end{lstlisting}
\begin{answerbox}Panel A shows the groups, panel B the link with sleep; the band is the 95 \% confidence interval.\end{answerbox}
{\small\hyperref[day:python:24]{Back to day 24}}
\answerday{25}{Project: tests}{answers:python:25}
\begin{enumerate}
  \item Which word gives the chance of a value above 80? \fromday{19}\\ $\rightarrow$ \code{dist.\allowbreak{}chance(above=80)}
  \item How do you get the number in row 1, column 2 of a grid? \fromday{9}\\ $\rightarrow$ \code{grid.\allowbreak{}number\_\allowbreak{}at(row=1,\allowbreak{} column=2)}
  \item Which word lists every step you did to a table? \fromday{6}\\ $\rightarrow$ \code{table.\allowbreak{}steps()}
  \item How do you give each class its own colour in a bar chart? \fromday{2}\\ $\rightarrow$ \code{colour=\textquotedbl{}class\textquotedbl{}} inside the chart step: \code{chart.\allowbreak{}bars(x=\textquotedbl{}subject\textquotedbl{},\allowbreak{} y=\textquotedbl{}marks\textquotedbl{},\allowbreak{} colour=\textquotedbl{}class\textquotedbl{})}
\end{enumerate}
\par\noindent{\sffamily\small\bfseries\color{tryframe}Task 1.} Run the test that answers your question. Write down the test's name, the effect size and the p-value.\par
\begin{lstlisting}[style=py]
import simplydata as sdata
table = sdata.table(sdata.example("school"))
table.compare_shape_to("marks", "normal", for_each="breakfast")
table.compare_averages("marks", for_each="breakfast")
table.correlation_between("sleep", "marks")
\end{lstlisting}
\begin{answerbox}Two groups of numbers: Welch's t-test. Two columns of numbers: Pearson's correlation.\end{answerbox}
{\small\hyperref[day:python:25]{Back to day 25}}
\answerday{26}{Project: models and fits}{answers:python:26}
\begin{enumerate}
  \item Which word fits a distribution to a column? \fromday{19}\\ $\rightarrow$ \code{table.\allowbreak{}fit\_\allowbreak{}distribution(\textquotedbl{}days\_\allowbreak{}absent\textquotedbl{},\allowbreak{} \textquotedbl{}negative binomial\textquotedbl{})}
  \item Which word draws one dot for each row, x against y? \fromday{3}\\ $\rightarrow$ \code{chart.\allowbreak{}dots(x=\textquotedbl{}hours\textquotedbl{},\allowbreak{} y=\textquotedbl{}marks\textquotedbl{})}
  \item The same children were measured twice. Which word compares the two? \fromday{12}\\ $\rightarrow$ \code{table.\allowbreak{}compare\_\allowbreak{}before\_\allowbreak{}after(\textquotedbl{}before\textquotedbl{},\allowbreak{} \textquotedbl{}after\textquotedbl{})}: the paired test.
  \item Which word draws one small chart for each class? \fromday{5}\\ $\rightarrow$ \code{chart.\allowbreak{}one\_\allowbreak{}chart\_\allowbreak{}per(\textquotedbl{}class\textquotedbl{})}
\end{enumerate}
\par\noindent{\sffamily\small\bfseries\color{tryframe}Task 1.} Explain your main column from two or three others, and say in one sentence what changed compared with the single test.\par
\begin{lstlisting}[style=py]
import simplydata as sdata
table = sdata.table(sdata.example("school"))
table.explain("marks", using=["hours", "sleep", "breakfast"])
table.fit_line(x="sleep", y="marks")
\end{lstlisting}
\begin{answerbox}Each line of the model is the effect of one column with the others held the same.\end{answerbox}
{\small\hyperref[day:python:26]{Back to day 26}}
\answerday{27}{Project: results text}{answers:python:27}
\begin{enumerate}
  \item Which word hands a table on to pandas? \fromday{13}\\ $\rightarrow$ \code{table.\allowbreak{}as\_\allowbreak{}pandas()}
  \item Which word explains marks from hours, sleep and breakfast together? \fromday{15}\\ $\rightarrow$ \code{table.\allowbreak{}explain(\textquotedbl{}marks\textquotedbl{},\allowbreak{} using=[\textquotedbl{}hours\textquotedbl{},\allowbreak{} \textquotedbl{}sleep\textquotedbl{},\allowbreak{} \textquotedbl{}breakfast\textquotedbl{}])}
  \item Which word tests whether the counts of categories differ between groups? \fromday{11}\\ $\rightarrow$ \code{table.\allowbreak{}compare\_\allowbreak{}counts(\textquotedbl{}result\textquotedbl{},\allowbreak{} for\_\allowbreak{}each=\textquotedbl{}class\textquotedbl{})}
  \item Which word tests whether the average marks differ between the classes? \fromday{10}\\ $\rightarrow$ \code{table.\allowbreak{}compare\_\allowbreak{}averages(\textquotedbl{}marks\textquotedbl{},\allowbreak{} for\_\allowbreak{}each=\textquotedbl{}class\textquotedbl{})}
\end{enumerate}
\par\noindent{\sffamily\small\bfseries\color{tryframe}Task 1.} Write a results paragraph, a methods paragraph and a figure caption for your project.\par
\begin{answerbox}For example: \code{\textquotedbl{}Pupils who had breakfast scored on average 3.\allowbreak{}4 marks more than pupils who did not (49.\allowbreak{}7 against 46.\allowbreak{}3; n = 115 and 65),\allowbreak{} but a difference this big could be chance (Welch's t-test,\allowbreak{} p = 0.\allowbreak{}051; Hedges' g = 0.\allowbreak{}31).\allowbreak{} Sleep was not recorded for 19 pupils; they were left out of the tests that use sleep.\allowbreak{}\textquotedbl{}} Use your own numbers, exactly as printed, and report a result that could be chance as plainly as one that could not.\end{answerbox}
{\small\hyperref[day:python:27]{Back to day 27}}
\answerday{28}{Review, and finish your project}{answers:python:28}
\begin{enumerate}
  \item Which word counts the empty cells in each column? \fromday{8}\\ $\rightarrow$ \code{table.\allowbreak{}count\_\allowbreak{}empty\_\allowbreak{}cells()}
  \item Which word draws survival curves? \fromday{17}\\ $\rightarrow$ \code{chart.\allowbreak{}survival\_\allowbreak{}curves(time=\textquotedbl{}days\textquotedbl{},\allowbreak{} event=\textquotedbl{}died\textquotedbl{},\allowbreak{} colour=\textquotedbl{}treatment\textquotedbl{})}
  \item Which word adds the principal components PC1 and PC2? \fromday{16}\\ $\rightarrow$ \code{table.\allowbreak{}add\_\allowbreak{}principal\_\allowbreak{}components(how\_\allowbreak{}many=2)}
  \item Which word estimates a column between the points you measured? \fromday{20}\\ $\rightarrow$ \code{table.\allowbreak{}values\_\allowbreak{}at(\textquotedbl{}population\textquotedbl{},\allowbreak{} at=\{\textquotedbl{}year\textquotedbl{}: [2001,\allowbreak{} 2011]\})}
  \item Which word fits a straight line on a table? \fromday{18}\\ $\rightarrow$ \code{table.\allowbreak{}fit\_\allowbreak{}line(x=\textquotedbl{}hours\textquotedbl{},\allowbreak{} y=\textquotedbl{}marks\textquotedbl{})}
  \item Which option makes a bar show the median instead of the average? \fromday{4}\\ $\rightarrow$ \code{show=\textquotedbl{}median\textquotedbl{}}
  \item Which word lists the practice tables? \fromday{1}\\ $\rightarrow$ \code{sdata.\allowbreak{}example()}; \code{sdata.\allowbreak{}example(\textquotedbl{}marks\textquotedbl{})} gets one of them.
  \item Which word gives the average of a column for each group? \fromday{9}\\ $\rightarrow$ \code{table.\allowbreak{}average\_\allowbreak{}of(\textquotedbl{}marks\textquotedbl{},\allowbreak{} for\_\allowbreak{}each=\textquotedbl{}class\textquotedbl{})}
  \item How do you look at the first 5 rows of a table? \fromday{6}\\ $\rightarrow$ \code{table.\allowbreak{}show\_\allowbreak{}first\_\allowbreak{}rows(5)}
  \item Which word moves the legend below the chart? \fromday{2}\\ $\rightarrow$ \code{chart.\allowbreak{}legend\_\allowbreak{}bottom()}
\end{enumerate}
\par\noindent{\sffamily\small\bfseries\color{tryframe}Task 1.} Without looking back: which word fits a formula, which finds where it is biggest, and which gives the chance of a value above a number? Use one of them on any table or formula.\par
\begin{lstlisting}[style=py]
import simplydata as sdata
dist = sdata.distribution("normal", average=50, standard_deviation=10)
dist.chance(above=70)
sdata.find_biggest("x * (10 - x)", between=[0, 10])
\end{lstlisting}
\begin{answerbox}\code{fit\_\allowbreak{}formula}, \code{find\_\allowbreak{}biggest} and chance: three families (\code{fit\_\allowbreak{}}, \code{find\_\allowbreak{}}, and the distribution words).\end{answerbox}
\par\noindent{\sffamily\small\bfseries\color{tryframe}Task 2.} Run your whole project script from the top. Check that it makes the clean table, the figure and the numbers in your text.\par
\begin{lstlisting}[style=py]
import simplydata as sdata
table = sdata.table(sdata.example("school"))
table.add_is_empty_column("sleep_was_empty", from_column="sleep")
table.compare_averages("marks", for_each="breakfast")
chart = table.chart()
chart.boxes(x="breakfast", y="marks")
chart.save("figure1.pdf")
table.save("school_clean.csv")
table.steps()
\end{lstlisting}
\begin{answerbox}One script from the raw file to the figure: anyone can run it again and get the same answers.\end{answerbox}
{\small\hyperref[day:python:28]{Back to day 28}}
\answerday{29}{Final test from memory}{answers:python:29}
\begin{enumerate}
  \item Which word finds where a formula is smallest? \fromday{19}\\ $\rightarrow$ \code{sdata.\allowbreak{}find\_\allowbreak{}smallest(\textquotedbl{}x**2 - 4*x + 1\textquotedbl{},\allowbreak{} between=[0,\allowbreak{} 5])}
  \item Which word compares one average with a number you choose? \fromday{12}\\ $\rightarrow$ \code{table.\allowbreak{}compare\_\allowbreak{}average\_\allowbreak{}to(\textquotedbl{}marks\textquotedbl{},\allowbreak{} 60)}
  \item Which word puts two charts next to each other on one page? \fromday{5}\\ $\rightarrow$ \code{sdata.\allowbreak{}side\_\allowbreak{}by\_\allowbreak{}side(a,\allowbreak{} b)}
  \item Which word draws a line with a dot at each point? \fromday{3}\\ $\rightarrow$ \code{chart.\allowbreak{}lines(x=\textquotedbl{}month\textquotedbl{},\allowbreak{} y=\textquotedbl{}temperature\textquotedbl{},\allowbreak{} dots=True)}
  \item How do you explain a yes-or-no column such as \code{\textquotedbl{}passed\textquotedbl{}}? \fromday{15}\\ $\rightarrow$ \code{outcome=\textquotedbl{}yes or no\textquotedbl{}} (logistic regression)
  \item Which word lists every change you made to a table, for your methods? \fromday{13}\\ $\rightarrow$ \code{table.\allowbreak{}steps()}
  \item Which word shows the rows with outliers? \fromday{11}\\ $\rightarrow$ \code{table.\allowbreak{}show\_\allowbreak{}rows\_\allowbreak{}with\_\allowbreak{}outliers(\textquotedbl{}marks\textquotedbl{},\allowbreak{} for\_\allowbreak{}each=\textquotedbl{}class\textquotedbl{})}
  \item Which word says which pairs of groups differ? \fromday{10}\\ $\rightarrow$ \code{table.\allowbreak{}compare\_\allowbreak{}each\_\allowbreak{}pair(\textquotedbl{}marks\textquotedbl{},\allowbreak{} for\_\allowbreak{}each=\textquotedbl{}class\textquotedbl{})}
  \item Which word fills an even map from scattered points, ready for a contour map? \fromday{20}\\ $\rightarrow$ \code{table.\allowbreak{}map\_\allowbreak{}of(\textquotedbl{}height\textquotedbl{},\allowbreak{} across=[\textquotedbl{}east\textquotedbl{},\allowbreak{} \textquotedbl{}north\textquotedbl{}])}
  \item How do you fit a curve with x and x squared? \fromday{18}\\ $\rightarrow$ \code{table.\allowbreak{}fit\_\allowbreak{}curve(x=\textquotedbl{}hours\textquotedbl{},\allowbreak{} y=\textquotedbl{}marks\textquotedbl{},\allowbreak{} degree=2)}
  \item Which word tests whether survival differs between groups? \fromday{17}\\ $\rightarrow$ \code{table.\allowbreak{}compare\_\allowbreak{}survival(time=\textquotedbl{}days\textquotedbl{},\allowbreak{} event=\textquotedbl{}died\textquotedbl{},\allowbreak{} for\_\allowbreak{}each=\textquotedbl{}treatment\textquotedbl{})}
  \item Which word puts rows that are alike into groups? \fromday{16}\\ $\rightarrow$ \code{table.\allowbreak{}add\_\allowbreak{}clusters(how\_\allowbreak{}many=2)}
  \item Which word tells whether two columns go up and down together? \fromday{9}\\ $\rightarrow$ \code{table.\allowbreak{}correlation\_\allowbreak{}between(\textquotedbl{}hours\textquotedbl{},\allowbreak{} \textquotedbl{}marks\textquotedbl{})}
  \item How do you fill the empty marks with the median? \fromday{8}\\ $\rightarrow$ \code{table.\allowbreak{}fill\_\allowbreak{}empty\_\allowbreak{}cells(\textquotedbl{}marks\textquotedbl{},\allowbreak{} with\_\allowbreak{}the=\textquotedbl{}median\textquotedbl{})}
\end{enumerate}
\par\noindent{\sffamily\small\bfseries\color{tryframe}Task 1.} From memory, in an empty file: open the \code{\textquotedbl{}marks\textquotedbl{}} table, remove the rows with empty cells, give the average marks for each class, test the difference, and save a two-panel figure (boxes, and dots of hours against marks with a fitted line) as a PDF.\par
\begin{lstlisting}[style=py]
import simplydata as sdata
table = sdata.table(sdata.example("marks"))
table.remove_rows_with_empty_cells()
table.average_of("marks", for_each="class")
table.compare_averages("marks", for_each="class")
a = table.chart()
a.boxes(x="class", y="marks")
b = table.chart()
b.dots(x="hours", y="marks")
b.fit_line()
figure = sdata.grid(a, b, per_row=2, letters=True)
figure.save("final.pdf")
\end{lstlisting}
\begin{answerbox}If you wrote this from memory, you know the core of simplydata. Words you missed: run their \code{sdata.\allowbreak{}help} again tomorrow.\end{answerbox}
{\small\hyperref[day:python:29]{Back to day 29}}
\answerday{30}{The same work in pandas and matplotlib}{answers:python:30}
\begin{enumerate}
  \item Which word estimates a column between the points you measured? \fromday{20}\\ $\rightarrow$ \code{table.\allowbreak{}values\_\allowbreak{}at(\textquotedbl{}population\textquotedbl{},\allowbreak{} at=\{\textquotedbl{}year\textquotedbl{}: [2001,\allowbreak{} 2011]\})}
  \item Which word finds where the two sides of an equation are equal? \fromday{19}\\ $\rightarrow$ \code{sdata.\allowbreak{}solve(\textquotedbl{}x**3 - 2*x = 5\textquotedbl{},\allowbreak{} between=[0,\allowbreak{} 3])}
  \item Which word fits a formula of your own, such as growth? \fromday{18}\\ $\rightarrow$ \code{table.\allowbreak{}fit\_\allowbreak{}formula(\textquotedbl{}population = a * exp(b * (year - 2000))\textquotedbl{})}
  \item Which word tests whether groups differ in their make-up across many columns (PERMANOVA)? \fromday{16}\\ $\rightarrow$ \code{table.\allowbreak{}compare\_\allowbreak{}profiles(for\_\allowbreak{}each=\textquotedbl{}grazed\textquotedbl{})}
\end{enumerate}
\par\noindent{\sffamily\small\bfseries\color{tryframe}Task 1.} In simplydata: keep classes A and B of the marks table, give the average marks of each class, test the difference and draw it as bars.\par
\begin{lstlisting}[style=py]
import simplydata as sdata
table = sdata.table(sdata.example("marks"))
table.keep_rows("class", one_of=["A", "B"])
table.average_of("marks", for_each="class")
table.compare_averages("marks", for_each="class")
chart = table.chart()
chart.bars(x="class", y="marks")
chart.save("simplydata_bars.png")
\end{lstlisting}
\begin{answerbox}The simplydata version: one word per step, and each answer printed in a sentence.\end{answerbox}
\par\noindent{\sffamily\small\bfseries\color{tryframe}Task 2.} Write the same analysis with pandas, SciPy and matplotlib, then compare the two line by line.\par
\begin{lstlisting}[style=py]
import matplotlib.pyplot as plt
import simplydata as sdata
from scipy import stats
df = sdata.example("marks")
df = df[df["class"].isin(["A", "B"])]
print(df.groupby("class")["marks"].mean())
a = df.loc[df["class"] == "A", "marks"]
b = df.loc[df["class"] == "B", "marks"]
print(stats.ttest_ind(a, b, equal_var=False))
fig, ax = plt.subplots()
means = df.groupby("class")["marks"].mean()
ax.bar(means.index, means.values)
fig.savefig("pandas_bars.png")
\end{lstlisting}
\begin{answerbox}The same in pandas, SciPy and matplotlib: \code{equal\_\allowbreak{}var=False} makes it Welch's t-test, as in simplydata.\end{answerbox}
{\small\hyperref[day:python:30]{Back to day 30}}

\trackstart{Answers: Beginner track}{The warm-up answers and a worked answer to each task}{answers:beginner}
There is often more than one good answer: if yours runs and answers the question, it is right. Every piece of code here runs as it is (the tests run it).
\answerday{1}{Python first steps: values and names}{answers:beginner:1}
\par\noindent{\sffamily\small\bfseries\color{tryframe}Task 1.} Give the text \code{\textquotedbl{}Bergen\textquotedbl{}} the name city and the number 2250 the name rain (the rain in a year, in mm). Print the words \code{\textquotedbl{}Rain in\textquotedbl{}}, the city, and the rain of an average month (rain / 12).\par
\begin{lstlisting}[style=py]
city = "Bergen"
rain = 2250
print("Rain in", city, rain / 12)
\end{lstlisting}
\begin{answerbox}It prints: Rain in Bergen 187.5. Text in quotes, numbers without, and print with commas between the values.\end{answerbox}
{\small\hyperref[day:beginner:1]{Back to day 1}}
\answerday{2}{Python first steps: lists and words}{answers:beginner:2}
\begin{enumerate}
  \item How do you write the text Asha in Python? \fromday{1}\\ $\rightarrow$ \code{\textquotedbl{}Asha\textquotedbl{}}, in quotes. Without quotes, Python thinks Asha is a name.
  \item How do you give the number 12 the name age? \fromday{1}\\ $\rightarrow$ age = 12. The name goes on the left of =, the value on the right.
  \item What does \code{print(age)} do? \fromday{1}\\ $\rightarrow$ It shows the value of age on the screen.
\end{enumerate}
\par\noindent{\sffamily\small\bfseries\color{tryframe}Task 1.} Make a list of five temperatures, turn it into simplydata numbers, and find the biggest one and the median.\par
\begin{lstlisting}[style=py]
import simplydata as sdata
temperatures = [4, 9, 13, 17, 20]
nums = sdata.numbers(temperatures)
nums.biggest()
nums.median()
\end{lstlisting}
\begin{answerbox}The biggest is 20 and the median (the middle value when they are in order) is 13.\end{answerbox}
{\small\hyperref[day:beginner:2]{Back to day 2}}
\answerday{3}{Reading an error, and your first chart}{answers:beginner:3}
\begin{enumerate}
  \item How do you write a list of the numbers 3, 5 and 2? \fromday{2}\\ $\rightarrow$ [3, 5, 2]: square brackets, commas between.
  \item In \code{nums.\allowbreak{}keep\_\allowbreak{}numbers(above=50)}, what is \code{above=50}? \fromday{2}\\ $\rightarrow$ An option with a name: the name, =, and the value, inside the round brackets.
  \item How do you write the text Asha in Python? \fromday{1}\\ $\rightarrow$ \code{\textquotedbl{}Asha\textquotedbl{}}, in quotes. Without quotes, Python thinks Asha is a name.
\end{enumerate}
\par\noindent{\sffamily\small\bfseries\color{tryframe}Task 1.} Find and fix the two mistakes in these lines, then run them: nums = \code{sdata.\allowbreak{}numbers([4,\allowbreak{} 8,\allowbreak{} 6)} and then \code{nums.\allowbreak{}averag()}\par
\begin{lstlisting}[style=py]
import simplydata as sdata
nums = sdata.numbers([4, 8, 6])
nums.average()
\end{lstlisting}
\begin{answerbox}The list needs its closing ] before the ), and the word is average. The average is 6.\end{answerbox}
\par\noindent{\sffamily\small\bfseries\color{tryframe}Task 2.} Make a bar chart of the \code{\textquotedbl{}weather\textquotedbl{}} table with \code{x=\textquotedbl{}city\textquotedbl{}} and \code{y=\textquotedbl{}rain\textquotedbl{}}, and save it as \code{\textquotedbl{}rain.\allowbreak{}png\textquotedbl{}}.\par
\begin{lstlisting}[style=py]
import simplydata as sdata
weather = sdata.example("weather")
chart = sdata.chart(weather)
chart.bars(x="city", y="rain")
chart.save("rain.png")
\end{lstlisting}
\begin{answerbox}One bar for each city, as tall as its average rain over the 12 months.\end{answerbox}
{\small\hyperref[day:beginner:3]{Back to day 3}}
\answerday{4}{Titles and labels}{answers:beginner:4}
\begin{enumerate}
  \item Where in an error message do you look first? \fromday{3}\\ $\rightarrow$ At the last line: it names the kind of error and says what went wrong.
  \item How do you give the number 12 the name age? \fromday{1}\\ $\rightarrow$ age = 12. The name goes on the left of =, the value on the right.
  \item What does NameError: name 'nmae' is not defined usually mean? \fromday{3}\\ $\rightarrow$ A name is misspelt, or was never given a value (with =) before it was used.
\end{enumerate}
\par\noindent{\sffamily\small\bfseries\color{tryframe}Task 1.} Draw dots of \code{\textquotedbl{}hours\textquotedbl{}} (x) against \code{\textquotedbl{}marks\textquotedbl{}} (y) from the marks table, with a title, a name for each axis and a caption at the bottom.\par
\begin{lstlisting}[style=py]
import simplydata as sdata
marks = sdata.example("marks")
chart = sdata.chart(marks)
chart.dots(x="hours", y="marks")
chart.title_top("More study, better marks?")
chart.label_x("Hours of study")
chart.label_y("Mark (out of 100)")
chart.caption_bottom("Made-up data")
chart.save("hours_and_marks.png")
\end{lstlisting}
\begin{answerbox}First draw, then add the words; the order of the word steps does not matter.\end{answerbox}
{\small\hyperref[day:beginner:4]{Back to day 4}}
\answerday{5}{Colours, groups and legends}{answers:beginner:5}
\begin{enumerate}
  \item Which word puts a title above the chart? \fromday{4}\\ $\rightarrow$ \code{chart.\allowbreak{}title\_\allowbreak{}top(\textquotedbl{}My title\textquotedbl{})}
  \item What does every simplydata program start with? \fromday{2}\\ $\rightarrow$ import simplydata as sdata
  \item How do you name the y-axis? \fromday{4}\\ $\rightarrow$ \code{chart.\allowbreak{}label\_\allowbreak{}y(\textquotedbl{}Marks out of 100\textquotedbl{})}
\end{enumerate}
\par\noindent{\sffamily\small\bfseries\color{tryframe}Task 1.} Weather table: draw lines of \code{\textquotedbl{}temperature\textquotedbl{}} over \code{\textquotedbl{}month\textquotedbl{}}, one colour for each city. Put the legend on top and use the style \code{\textquotedbl{}paper\textquotedbl{}}.\par
\begin{lstlisting}[style=py]
import simplydata as sdata
weather = sdata.example("weather")
chart = sdata.chart(weather)
chart.lines(x="month", y="temperature", colour="city")
chart.legend_top()
chart.style("paper")
chart.save("temperature.png")
\end{lstlisting}
\begin{answerbox}\code{colour=\textquotedbl{}city\textquotedbl{}} draws one line for each city and adds a legend; the style changes the whole look.\end{answerbox}
{\small\hyperref[day:beginner:5]{Back to day 5}}
\answerday{6}{Different chart types}{answers:beginner:6}
\begin{enumerate}
  \item How do you give each class its own colour in a bar chart? \fromday{5}\\ $\rightarrow$ \code{colour=\textquotedbl{}class\textquotedbl{}} inside the chart step: \code{chart.\allowbreak{}bars(x=\textquotedbl{}subject\textquotedbl{},\allowbreak{} y=\textquotedbl{}marks\textquotedbl{},\allowbreak{} colour=\textquotedbl{}class\textquotedbl{})}
  \item Python says SyntaxError. What do you check first? \fromday{3}\\ $\rightarrow$ Quote marks and brackets: every one that opens must be closed.
  \item Which word moves the legend below the chart? \fromday{5}\\ $\rightarrow$ \code{chart.\allowbreak{}legend\_\allowbreak{}bottom()}
\end{enumerate}
\par\noindent{\sffamily\small\bfseries\color{tryframe}Task 1.} Draw violins of \code{\textquotedbl{}temperature\textquotedbl{}} for each city in the weather table, with all the dots on top.\par
\begin{lstlisting}[style=py]
import simplydata as sdata
weather = sdata.example("weather")
chart = sdata.chart(weather)
chart.violins(x="city", y="temperature")
chart.all_dots(x="city", y="temperature")
chart.save("violins.png")
\end{lstlisting}
\begin{answerbox}A violin shows the shape of the spread; the dots show every month.\end{answerbox}
{\small\hyperref[day:beginner:6]{Back to day 6}}
\answerday{7}{Review of the first week}{answers:beginner:7}
\begin{enumerate}
  \item Which chart shows how the values of one column are spread, in bins? \fromday{6}\\ $\rightarrow$ \code{chart.\allowbreak{}histogram(x=\textquotedbl{}marks\textquotedbl{},\allowbreak{} bins=8)}
  \item The names under the bars bump into each other. Which word helps? \fromday{4}\\ $\rightarrow$ \code{chart.\allowbreak{}rotate\_\allowbreak{}labels\_\allowbreak{}x(45)}
  \item Which word gives the average of numbers called nums? \fromday{2}\\ $\rightarrow$ \code{nums.\allowbreak{}average()}
  \item Which word lets you choose the colours of the groups? \fromday{5}\\ $\rightarrow$ \code{chart.\allowbreak{}colours(\textquotedbl{}orange\textquotedbl{},\allowbreak{} \textquotedbl{}teal\textquotedbl{},\allowbreak{} \textquotedbl{}purple\textquotedbl{})}
  \item Which word lists the practice tables? \fromday{3}\\ $\rightarrow$ \code{sdata.\allowbreak{}example()}; \code{sdata.\allowbreak{}example(\textquotedbl{}marks\textquotedbl{})} gets one of them.
  \item What does \code{print(age)} do? \fromday{1}\\ $\rightarrow$ It shows the value of age on the screen.
\end{enumerate}
\par\noindent{\sffamily\small\bfseries\color{tryframe}Task 1.} Make numbers of five marks and give their average. Then draw bars of the marks table with \code{x=\textquotedbl{}subject\textquotedbl{}}, \code{y=\textquotedbl{}marks\textquotedbl{}}, one colour for each class, and a title. Save it.\par
\begin{lstlisting}[style=py]
import simplydata as sdata
nums = sdata.numbers([48, 71, 55, 88, 62])
nums.average()
marks = sdata.example("marks")
chart = sdata.chart(marks)
chart.bars(x="subject", y="marks", colour="class")
chart.title_top("Marks in each subject")
chart.save("subjects.png")
\end{lstlisting}
\begin{answerbox}Python lists (day 2), bars and save (day 3), titles (day 4) and colours (day 5).\end{answerbox}
{\small\hyperref[day:beginner:7]{Back to day 7}}
\answerday{8}{Averages, error bars and fit lines}{answers:beginner:8}
\begin{enumerate}
  \item Which two words show the spread of each group, with every value as a dot on top? \fromday{6}\\ $\rightarrow$ \code{chart.\allowbreak{}boxes(x=\textquotedbl{}class\textquotedbl{},\allowbreak{} y=\textquotedbl{}marks\textquotedbl{})} and \code{chart.\allowbreak{}all\_\allowbreak{}dots(x=\textquotedbl{}class\textquotedbl{},\allowbreak{} y=\textquotedbl{}marks\textquotedbl{})}
  \item How do you put the bars in your own order? \fromday{5}\\ $\rightarrow$ \code{chart.\allowbreak{}order\_\allowbreak{}x(\textquotedbl{}C\textquotedbl{},\allowbreak{} \textquotedbl{}B\textquotedbl{},\allowbreak{} \textquotedbl{}A\textquotedbl{})}
  \item How do you write the text Asha in Python? \fromday{1}\\ $\rightarrow$ \code{\textquotedbl{}Asha\textquotedbl{}}, in quotes. Without quotes, Python thinks Asha is a name.
\end{enumerate}
\par\noindent{\sffamily\small\bfseries\color{tryframe}Task 1.} Weather table: draw bars of the average \code{\textquotedbl{}rain\textquotedbl{}} for each city with 95 \% confidence intervals, and write the numbers on the bars.\par
\begin{lstlisting}[style=py]
import simplydata as sdata
weather = sdata.example("weather")
chart = sdata.chart(weather)
chart.bars(x="city", y="rain", error="confidence interval")
chart.numbers_on_bars()
chart.save("rain_ci.png")
\end{lstlisting}
\begin{answerbox}A line is printed saying which error bars were drawn, so the caption can say it too.\end{answerbox}
{\small\hyperref[day:beginner:8]{Back to day 8}}
\answerday{9}{Axes and reference lines}{answers:beginner:9}
\begin{enumerate}
  \item Which option makes a bar show the median instead of the average? \fromday{8}\\ $\rightarrow$ \code{show=\textquotedbl{}median\textquotedbl{}}
  \item Which word draws one dot for each row, x against y? \fromday{6}\\ $\rightarrow$ \code{chart.\allowbreak{}dots(x=\textquotedbl{}hours\textquotedbl{},\allowbreak{} y=\textquotedbl{}marks\textquotedbl{})}
  \item How do you write a list of the numbers 3, 5 and 2? \fromday{2}\\ $\rightarrow$ [3, 5, 2]: square brackets, commas between.
\end{enumerate}
\par\noindent{\sffamily\small\bfseries\color{tryframe}Task 1.} Growth table: draw lines of \code{\textquotedbl{}population\textquotedbl{}} over \code{\textquotedbl{}year\textquotedbl{}}, one for each country, with a line at x = 2010 and the y-axis starting at 0.\par
\begin{lstlisting}[style=py]
import simplydata as sdata
growth = sdata.example("growth")
chart = sdata.chart(growth)
chart.lines(x="year", y="population", colour="country")
chart.line_at_x(2010, label="2010")
chart.range_y(0, 60)
chart.save("growth.png")
\end{lstlisting}
\begin{answerbox}\code{range\_\allowbreak{}y} takes the lowest and the highest value of the axis; \code{line\_\allowbreak{}at\_\allowbreak{}x} draws a line going up.\end{answerbox}
{\small\hyperref[day:beginner:9]{Back to day 9}}
\answerday{10}{Several charts together}{answers:beginner:10}
\begin{enumerate}
  \item Which word sets the y-axis from 0 to 100? \fromday{9}\\ $\rightarrow$ \code{chart.\allowbreak{}range\_\allowbreak{}y(0,\allowbreak{} 100)}
  \item Which word draws a line with a dot at each point? \fromday{6}\\ $\rightarrow$ \code{chart.\allowbreak{}lines(x=\textquotedbl{}month\textquotedbl{},\allowbreak{} y=\textquotedbl{}temperature\textquotedbl{},\allowbreak{} dots=True)}
  \item How do you start a chart of a table called marks? \fromday{3}\\ $\rightarrow$ chart = \code{sdata.\allowbreak{}chart(marks)}
\end{enumerate}
\par\noindent{\sffamily\small\bfseries\color{tryframe}Task 1.} Make a histogram and a box chart of the marks, and put them one above the other on one page.\par
\begin{lstlisting}[style=py]
import simplydata as sdata
marks = sdata.example("marks")
a = sdata.chart(marks).histogram(x="marks")
b = sdata.chart(marks).boxes(x="class", y="marks")
page = sdata.one_above_other(a, b)
page.save("two_charts.png")
\end{lstlisting}
\begin{answerbox}Each chart is made first; \code{one\_\allowbreak{}above\_\allowbreak{}other} puts them on one page.\end{answerbox}
{\small\hyperref[day:beginner:10]{Back to day 10}}
\answerday{11}{A figure for a journal, and word families}{answers:beginner:11}
\begin{enumerate}
  \item Which word draws one small chart for each class? \fromday{10}\\ $\rightarrow$ \code{chart.\allowbreak{}one\_\allowbreak{}chart\_\allowbreak{}per(\textquotedbl{}class\textquotedbl{})}
  \item Which option adds error bars that say how sure we are of the average? \fromday{8}\\ $\rightarrow$ \code{error=\textquotedbl{}standard error\textquotedbl{}}, or \code{error=\textquotedbl{}confidence interval\textquotedbl{}} for the 95 \% interval.
  \item Which word makes all the text on a chart bigger or smaller? \fromday{4}\\ $\rightarrow$ \code{chart.\allowbreak{}font\_\allowbreak{}size(14)}
\end{enumerate}
\par\noindent{\sffamily\small\bfseries\color{tryframe}Task 1.} A figure for a journal: dots of \code{\textquotedbl{}hours\textquotedbl{}} against \code{\textquotedbl{}marks\textquotedbl{}}, style \code{\textquotedbl{}paper\textquotedbl{}}, 8.5 cm wide and 7 cm high, font size 8, saved as a \code{.\allowbreak{}tif} with 600 dpi.\par
\begin{lstlisting}[style=py]
import simplydata as sdata
marks = sdata.example("marks")
chart = sdata.chart(marks)
chart.dots(x="hours", y="marks")
chart.style("paper")
chart.size(width=8.5, height=7, unit="cm")
chart.font_size(8)
chart.save("figure1.tif", resolution=600)
\end{lstlisting}
\begin{answerbox}Journals often ask for 8.5 cm (one column) or 17 cm (two columns) and 300 to 600 dpi.\end{answerbox}
\par\noindent{\sffamily\small\bfseries\color{tryframe}Task 2.} Save one chart as \code{.\allowbreak{}png}, \code{.\allowbreak{}pdf} and \code{.\allowbreak{}svg}. Then guess, before you look anything up: which word hides the legend, and which puts it on the left? Try your guesses on a chart with \code{colour=}.\par
\begin{lstlisting}[style=py]
import simplydata as sdata
marks = sdata.example("marks")
chart = sdata.chart(marks)
chart.bars(x="subject", y="marks", colour="class")
chart.save("subjects.png")
chart.save("subjects.pdf")
chart.save("subjects.svg")
chart.legend_left()
chart.hide_legend()
chart.save("no_legend.png")
\end{lstlisting}
\begin{answerbox}\code{legend\_\allowbreak{}left} follows \code{legend\_\allowbreak{}bottom} and \code{legend\_\allowbreak{}top}; \code{hide\_\allowbreak{}legend} follows \code{show\_\allowbreak{}} and \code{hide\_\allowbreak{}}. A family is learned once.\end{answerbox}
{\small\hyperref[day:beginner:11]{Back to day 11}}
\answerday{12}{Open a table and look at it}{answers:beginner:12}
\begin{enumerate}
  \item Which word adds a second y-axis on the right? \fromday{11}\\ $\rightarrow$ \code{chart.\allowbreak{}second\_\allowbreak{}y(\textquotedbl{}rain\textquotedbl{})}
  \item Which word draws a line across at y = 50, with a label? \fromday{9}\\ $\rightarrow$ \code{chart.\allowbreak{}line\_\allowbreak{}at\_\allowbreak{}y(50,\allowbreak{} label=\textquotedbl{}pass mark\textquotedbl{})}
  \item How do you give each class its own colour in a bar chart? \fromday{5}\\ $\rightarrow$ \code{colour=\textquotedbl{}class\textquotedbl{}} inside the chart step: \code{chart.\allowbreak{}bars(x=\textquotedbl{}subject\textquotedbl{},\allowbreak{} y=\textquotedbl{}marks\textquotedbl{},\allowbreak{} colour=\textquotedbl{}class\textquotedbl{})}
\end{enumerate}
\par\noindent{\sffamily\small\bfseries\color{tryframe}Task 1.} Open the \code{\textquotedbl{}shop\textquotedbl{}} practice table. How many rows does it have, and how many rows does each shop have?\par
\begin{lstlisting}[style=py]
import simplydata as sdata
table = sdata.table(sdata.example("shop"))
table.count_rows()
table.count_each("shop")
table.describe_columns()
\end{lstlisting}
\begin{answerbox}\code{count\_\allowbreak{}rows} counts all the rows, \code{count\_\allowbreak{}each} the rows of each shop; \code{describe\_\allowbreak{}columns} also shows the empty cells.\end{answerbox}
{\small\hyperref[day:beginner:12]{Back to day 12}}
\answerday{13}{Choose rows and columns, and undo}{answers:beginner:13}
\begin{enumerate}
  \item How do you look at the first 5 rows of a table? \fromday{12}\\ $\rightarrow$ \code{table.\allowbreak{}show\_\allowbreak{}first\_\allowbreak{}rows(5)}
  \item Which word puts two charts next to each other on one page? \fromday{10}\\ $\rightarrow$ \code{sdata.\allowbreak{}side\_\allowbreak{}by\_\allowbreak{}side(a,\allowbreak{} b)}
  \item Which chart shows how the values of one column are spread, in bins? \fromday{6}\\ $\rightarrow$ \code{chart.\allowbreak{}histogram(x=\textquotedbl{}marks\textquotedbl{},\allowbreak{} bins=8)}
\end{enumerate}
\par\noindent{\sffamily\small\bfseries\color{tryframe}Task 1.} Weather table: keep the city \code{\textquotedbl{}Sunnyvale\textquotedbl{}}, keep only the columns month, temperature and rain, sort by rain with the biggest first, and save it as \code{\textquotedbl{}sunny.\allowbreak{}csv\textquotedbl{}}.\par
\begin{lstlisting}[style=py]
import simplydata as sdata
table = sdata.table(sdata.example("weather"))
table.keep_rows("city", equal_to="Sunnyvale")
table.keep_columns("month", "temperature", "rain")
table.sort_by("rain", biggest_first=True)
table.save("sunny.csv")
table.steps()
\end{lstlisting}
\begin{answerbox}Each step prints what it did; \code{steps()} lists them all, in order.\end{answerbox}
{\small\hyperref[day:beginner:13]{Back to day 13}}
\answerday{14}{Review of the second week}{answers:beginner:14}
\begin{enumerate}
  \item Which word keeps only the rows where subject is Maths? \fromday{13}\\ $\rightarrow$ \code{table.\allowbreak{}keep\_\allowbreak{}rows(\textquotedbl{}subject\textquotedbl{},\allowbreak{} equal\_\allowbreak{}to=\textquotedbl{}Maths\textquotedbl{})}
  \item Which word draws one bar for each name, as tall as the marks? \fromday{3}\\ $\rightarrow$ \code{chart.\allowbreak{}bars(x=\textquotedbl{}name\textquotedbl{},\allowbreak{} y=\textquotedbl{}marks\textquotedbl{})}
  \item Which word draws a straight line through the dots? \fromday{8}\\ $\rightarrow$ \code{chart.\allowbreak{}fit\_\allowbreak{}line()}, after \code{chart.\allowbreak{}dots(.\allowbreak{}.\allowbreak{}.\allowbreak{})}
  \item Which word sets the size of a figure in cm? \fromday{11}\\ $\rightarrow$ \code{chart.\allowbreak{}size(width=17,\allowbreak{} height=8,\allowbreak{} unit=\textquotedbl{}cm\textquotedbl{})}
  \item A few values are much bigger than the rest. Which word helps? \fromday{9}\\ $\rightarrow$ \code{chart.\allowbreak{}log\_\allowbreak{}scale\_\allowbreak{}y()}
  \item Which word describes each column: its kind, empty cells and numbers? \fromday{12}\\ $\rightarrow$ \code{table.\allowbreak{}describe\_\allowbreak{}columns()}
  \item How do you put four charts in two rows, with the letters A to D? \fromday{10}\\ $\rightarrow$ \code{sdata.\allowbreak{}grid(a,\allowbreak{} b,\allowbreak{} c,\allowbreak{} d,\allowbreak{} per\_\allowbreak{}row=2,\allowbreak{} letters=True)}
  \item How do you give the number 12 the name age? \fromday{1}\\ $\rightarrow$ age = 12. The name goes on the left of =, the value on the right.
\end{enumerate}
\par\noindent{\sffamily\small\bfseries\color{tryframe}Task 1.} Weather table: keep the city \code{\textquotedbl{}Frostby\textquotedbl{}}, sort by \code{\textquotedbl{}temperature\textquotedbl{}}, and draw a line of the temperature over the month, with a line at y = 0 and the label \code{\textquotedbl{}freezing\textquotedbl{}}.\par
\begin{lstlisting}[style=py]
import simplydata as sdata
table = sdata.table(sdata.example("weather"))
table.keep_rows("city", equal_to="Frostby")
table.sort_by("temperature")
table.show_first_rows(3)
chart = table.chart()
chart.lines(x="month_number", y="temperature")
chart.line_at_y(0, label="freezing")
chart.save("frostby.png")
\end{lstlisting}
\begin{answerbox}Tables (days 12 and 13) feed charts (days 3 to 11).\end{answerbox}
{\small\hyperref[day:beginner:14]{Back to day 14}}
\answerday{15}{Clean the data and add columns}{answers:beginner:15}
\begin{enumerate}
  \item Which word sorts with the biggest marks first? \fromday{13}\\ $\rightarrow$ \code{table.\allowbreak{}sort\_\allowbreak{}by(\textquotedbl{}marks\textquotedbl{},\allowbreak{} biggest\_\allowbreak{}first=True)}
  \item Which word counts the rows of each fruit? \fromday{12}\\ $\rightarrow$ \code{table.\allowbreak{}count\_\allowbreak{}each(\textquotedbl{}fruit\textquotedbl{})}
  \item Which word checks by eye whether values are bell-shaped? \fromday{8}\\ $\rightarrow$ \code{chart.\allowbreak{}qq\_\allowbreak{}dots(\textquotedbl{}marks\textquotedbl{})}
  \item What does \code{print(age)} do? \fromday{1}\\ $\rightarrow$ It shows the value of age on the screen.
\end{enumerate}
\par\noindent{\sffamily\small\bfseries\color{tryframe}Task 1.} Shop table: count the empty cells, fill the empty \code{\textquotedbl{}sales\textquotedbl{}} with the line between neighbours for each shop (sorted by \code{\textquotedbl{}date\textquotedbl{}}), and add a column \code{\textquotedbl{}weekday\textquotedbl{}} with the name of the weekday.\par
\begin{lstlisting}[style=py]
import simplydata as sdata
table = sdata.table(sdata.example("shop"))
table.count_empty_cells()
table.fill_empty_cells("sales", with_the="line between neighbours", sorted_by="date", for_each="shop")
table.add_date_part_column("weekday", from_column="date", part="weekday name")
table.show_first_rows()
\end{lstlisting}
\begin{answerbox}\code{for\_\allowbreak{}each=} keeps the two shops apart, so a gap is never filled from the other shop's days.\end{answerbox}
{\small\hyperref[day:beginner:15]{Back to day 15}}
\answerday{16}{Summaries, and from a table to a chart}{answers:beginner:16}
\begin{enumerate}
  \item Which word counts the empty cells in each column? \fromday{15}\\ $\rightarrow$ \code{table.\allowbreak{}count\_\allowbreak{}empty\_\allowbreak{}cells()}
  \item You made a mistake in the last step. Which word takes it back? \fromday{13}\\ $\rightarrow$ \code{table.\allowbreak{}undo()}; \code{table.\allowbreak{}start\_\allowbreak{}again()} goes back to the start.
  \item Which word shows the numbers on the y-axis as percents? \fromday{9}\\ $\rightarrow$ \code{chart.\allowbreak{}tick\_\allowbreak{}format\_\allowbreak{}y(\textquotedbl{}percent\textquotedbl{})}
  \item In \code{nums.\allowbreak{}keep\_\allowbreak{}numbers(above=50)}, what is \code{above=50}? \fromday{2}\\ $\rightarrow$ An option with a name: the name, =, and the value, inside the round brackets.
\end{enumerate}
\par\noindent{\sffamily\small\bfseries\color{tryframe}Task 1.} Fruit table: the total \code{\textquotedbl{}count\textquotedbl{}} of each fruit, drawn as bars with the numbers on the bars.\par
\begin{lstlisting}[style=py]
import simplydata as sdata
table = sdata.table(sdata.example("fruit"))
totals = table.total_of("count", for_each="fruit")
chart = totals.chart()
chart.bars(x="fruit", y="total of count")
chart.numbers_on_bars()
chart.save("fruit_totals.png")
\end{lstlisting}
\begin{answerbox}A summary gives back a new table, with a column called \code{\textquotedbl{}total of count\textquotedbl{}}, ready for a chart.\end{answerbox}
{\small\hyperref[day:beginner:16]{Back to day 16}}
\answerday{17}{Lists and grids of numbers}{answers:beginner:17}
\begin{enumerate}
  \item Which word gives the average of a column for each group? \fromday{16}\\ $\rightarrow$ \code{table.\allowbreak{}average\_\allowbreak{}of(\textquotedbl{}marks\textquotedbl{},\allowbreak{} for\_\allowbreak{}each=\textquotedbl{}class\textquotedbl{})}
  \item Which word keeps rows 10 to 20, by their place in the table? \fromday{13}\\ $\rightarrow$ \code{table.\allowbreak{}keep\_\allowbreak{}rows\_\allowbreak{}by\_\allowbreak{}position(10,\allowbreak{} to=20)}
  \item Which word draws one small chart for each class? \fromday{10}\\ $\rightarrow$ \code{chart.\allowbreak{}one\_\allowbreak{}chart\_\allowbreak{}per(\textquotedbl{}class\textquotedbl{})}
  \item How do you save a chart as a PDF? \fromday{3}\\ $\rightarrow$ \code{chart.\allowbreak{}save(\textquotedbl{}my\_\allowbreak{}chart.\allowbreak{}pdf\textquotedbl{})}: the end of the file name decides the type.
\end{enumerate}
\par\noindent{\sffamily\small\bfseries\color{tryframe}Task 1.} Make 200 random heights with average 150 and standard deviation 8 (\code{seed=1}). What is their 90th percentile? Draw them as a histogram.\par
\begin{lstlisting}[style=py]
import simplydata as sdata
heights = sdata.random_numbers(how_many=200, average=150, standard_deviation=8, seed=1)
heights.percentile(90)
chart = heights.chart()
chart.histogram(x="value")
chart.save("heights.png")
\end{lstlisting}
\begin{answerbox}About 157: nine heights in ten are below it. \code{seed=1} gives the same numbers every time.\end{answerbox}
{\small\hyperref[day:beginner:17]{Back to day 17}}
\answerday{18}{Is a difference real? Tests and effect sizes}{answers:beginner:18}
\begin{enumerate}
  \item How do you make simplydata numbers from a list? \fromday{17}\\ $\rightarrow$ nums = \code{sdata.\allowbreak{}numbers([12,\allowbreak{} 15,\allowbreak{} 9,\allowbreak{} 20,\allowbreak{} 14])}
  \item How do you fill the empty marks with the median? \fromday{15}\\ $\rightarrow$ \code{table.\allowbreak{}fill\_\allowbreak{}empty\_\allowbreak{}cells(\textquotedbl{}marks\textquotedbl{},\allowbreak{} with\_\allowbreak{}the=\textquotedbl{}median\textquotedbl{})}
  \item How do you save a picture sharp enough for a journal? \fromday{11}\\ $\rightarrow$ \code{chart.\allowbreak{}save(\textquotedbl{}figure.\allowbreak{}tif\textquotedbl{},\allowbreak{} resolution=600)}, or save it as \code{.\allowbreak{}pdf}
  \item Which word puts a title above the chart? \fromday{4}\\ $\rightarrow$ \code{chart.\allowbreak{}title\_\allowbreak{}top(\textquotedbl{}My title\textquotedbl{})}
\end{enumerate}
\par\noindent{\sffamily\small\bfseries\color{tryframe}Task 1.} School table: do the average marks differ between pupils who had breakfast and pupils who did not? Write down the difference, the effect size and the p-value.\par
\begin{lstlisting}[style=py]
import simplydata as sdata
table = sdata.table(sdata.example("school"))
table.average_of("marks", for_each="breakfast")
table.compare_averages("marks", for_each="breakfast")
\end{lstlisting}
\begin{answerbox}Two groups, so it is Welch's t-test; the sentence gives the difference, Hedges' g and the p-value.\end{answerbox}
{\small\hyperref[day:beginner:18]{Back to day 18}}
\answerday{19}{Counts, outliers and many tests}{answers:beginner:19}
\begin{enumerate}
  \item Which word tests whether the average marks differ between the classes? \fromday{18}\\ $\rightarrow$ \code{table.\allowbreak{}compare\_\allowbreak{}averages(\textquotedbl{}marks\textquotedbl{},\allowbreak{} for\_\allowbreak{}each=\textquotedbl{}class\textquotedbl{})}
  \item Which word tells whether two columns go up and down together? \fromday{16}\\ $\rightarrow$ \code{table.\allowbreak{}correlation\_\allowbreak{}between(\textquotedbl{}hours\textquotedbl{},\allowbreak{} \textquotedbl{}marks\textquotedbl{})}
  \item Which word counts the rows? \fromday{12}\\ $\rightarrow$ \code{table.\allowbreak{}count\_\allowbreak{}rows()}
  \item Which word moves the legend below the chart? \fromday{5}\\ $\rightarrow$ \code{chart.\allowbreak{}legend\_\allowbreak{}bottom()}
\end{enumerate}
\par\noindent{\sffamily\small\bfseries\color{tryframe}Task 1.} School table: is the share of pupils who \code{\textquotedbl{}passed\textquotedbl{}} different with and without breakfast?\par
\begin{lstlisting}[style=py]
import simplydata as sdata
table = sdata.table(sdata.example("school"))
table.count_rows(for_each="breakfast", across="passed")
table.compare_counts("passed", for_each="breakfast")
\end{lstlisting}
\begin{answerbox}passed is a category (yes or no), so the test compares counts: the chi-square test, with Cramer's V as the effect size.\end{answerbox}
{\small\hyperref[day:beginner:19]{Back to day 19}}
\answerday{20}{Choosing a test, and a whole analysis}{answers:beginner:20}
\begin{enumerate}
  \item Which word tests whether the counts of categories differ between groups? \fromday{19}\\ $\rightarrow$ \code{table.\allowbreak{}compare\_\allowbreak{}counts(\textquotedbl{}result\textquotedbl{},\allowbreak{} for\_\allowbreak{}each=\textquotedbl{}class\textquotedbl{})}
  \item Which word gives the middle value of numbers? \fromday{17}\\ $\rightarrow$ \code{nums.\allowbreak{}median()}
  \item Which word lists every step you did to a table? \fromday{13}\\ $\rightarrow$ \code{table.\allowbreak{}steps()}
  \item Which two words show the spread of each group, with every value as a dot on top? \fromday{6}\\ $\rightarrow$ \code{chart.\allowbreak{}boxes(x=\textquotedbl{}class\textquotedbl{},\allowbreak{} y=\textquotedbl{}marks\textquotedbl{})} and \code{chart.\allowbreak{}all\_\allowbreak{}dots(x=\textquotedbl{}class\textquotedbl{},\allowbreak{} y=\textquotedbl{}marks\textquotedbl{})}
\end{enumerate}
\par\noindent{\sffamily\small\bfseries\color{tryframe}Task 1.} Marks table: is each child's Maths mark different from the same child's Science mark? Choose the right word first.\par
\begin{lstlisting}[style=py]
import simplydata as sdata
table = sdata.table(sdata.example("marks"))
table.keep_columns("name", "subject", "marks")
table.turn_rows_into_columns(names_from="subject", values_from="marks")
table.compare_before_after("Maths", "Science")
\end{lstlisting}
\begin{answerbox}The same children twice, so the test must be paired. First one row per child, with a Maths and a Science column.\end{answerbox}
\par\noindent{\sffamily\small\bfseries\color{tryframe}Task 2.} School table: is the average \code{\textquotedbl{}sleep\textquotedbl{}} different from 8 hours?\par
\begin{lstlisting}[style=py]
import simplydata as sdata
table = sdata.table(sdata.example("school"))
table.compare_average_to("sleep", 8)
\end{lstlisting}
\begin{answerbox}One average against a number you choose: the one-sample t-test. Empty sleep cells are skipped, with a note.\end{answerbox}
\par\noindent{\sffamily\small\bfseries\color{tryframe}Task 3.} From memory: weather table, compare the average \code{\textquotedbl{}temperature\textquotedbl{}} of the cities, draw boxes with all the dots, and save the figure as a PDF and the table as \code{\textquotedbl{}weather.\allowbreak{}xlsx\textquotedbl{}}.\par
\begin{lstlisting}[style=py]
import simplydata as sdata
table = sdata.table(sdata.example("weather"))
table.compare_averages("temperature", for_each="city")
chart = table.chart()
chart.boxes(x="city", y="temperature")
chart.all_dots(x="city", y="temperature")
chart.save("temperature.pdf")
table.save("weather.xlsx")
\end{lstlisting}
\begin{answerbox}Question, test, figure, saved table: the whole analysis in eight lines.\end{answerbox}
{\small\hyperref[day:beginner:20]{Back to day 20}}
\answerday{21}{Review of the third week}{answers:beginner:21}
\begin{enumerate}
  \item The same children were measured twice. Which word compares the two? \fromday{20}\\ $\rightarrow$ \code{table.\allowbreak{}compare\_\allowbreak{}before\_\allowbreak{}after(\textquotedbl{}before\textquotedbl{},\allowbreak{} \textquotedbl{}after\textquotedbl{})}: the paired test.
  \item Which word turns a column of text such as \code{\textquotedbl{}144\textquotedbl{}} into numbers? \fromday{15}\\ $\rightarrow$ \code{table.\allowbreak{}make\_\allowbreak{}numbers(\textquotedbl{}height\textquotedbl{})}
  \item Which word writes a note with an arrow on the chart? \fromday{11}\\ $\rightarrow$ \code{chart.\allowbreak{}note(\textquotedbl{}X marks the spot\textquotedbl{},\allowbreak{} x=10,\allowbreak{} y=4,\allowbreak{} arrow=True)}
  \item Which word says which pairs of groups differ? \fromday{18}\\ $\rightarrow$ \code{table.\allowbreak{}compare\_\allowbreak{}each\_\allowbreak{}pair(\textquotedbl{}marks\textquotedbl{},\allowbreak{} for\_\allowbreak{}each=\textquotedbl{}class\textquotedbl{})}
  \item Which word joins two tables on a column they share? \fromday{16}\\ $\rightarrow$ \code{table.\allowbreak{}join\_\allowbreak{}with(ages,\allowbreak{} using=\textquotedbl{}name\textquotedbl{})}
  \item Which word shows the rows with outliers? \fromday{19}\\ $\rightarrow$ \code{table.\allowbreak{}show\_\allowbreak{}rows\_\allowbreak{}with\_\allowbreak{}outliers(\textquotedbl{}marks\textquotedbl{},\allowbreak{} for\_\allowbreak{}each=\textquotedbl{}class\textquotedbl{})}
  \item How do you make 100 random numbers with average 50 and standard deviation 10, the same every time? \fromday{17}\\ $\rightarrow$ \code{sdata.\allowbreak{}random\_\allowbreak{}numbers(how\_\allowbreak{}many=100,\allowbreak{} average=50,\allowbreak{} standard\_\allowbreak{}deviation=10,\allowbreak{} seed=1)}
  \item Which option makes a bar show the median instead of the average? \fromday{8}\\ $\rightarrow$ \code{show=\textquotedbl{}median\textquotedbl{}}
  \item How do you write the text Asha in Python? \fromday{1}\\ $\rightarrow$ \code{\textquotedbl{}Asha\textquotedbl{}}, in quotes. Without quotes, Python thinks Asha is a name.
\end{enumerate}
\par\noindent{\sffamily\small\bfseries\color{tryframe}Task 1.} School table: the average \code{\textquotedbl{}marks\textquotedbl{}} for each value of \code{\textquotedbl{}breakfast\textquotedbl{}}, a test of the difference, and boxes of the marks for each group.\par
\begin{lstlisting}[style=py]
import simplydata as sdata
table = sdata.table(sdata.example("school"))
table.average_of("marks", for_each="breakfast")
table.compare_averages("marks", for_each="breakfast")
chart = table.chart()
chart.boxes(x="breakfast", y="marks")
chart.save("breakfast.png")
\end{lstlisting}
\begin{answerbox}A summary, a test and a chart of the same question.\end{answerbox}
{\small\hyperref[day:beginner:21]{Back to day 21}}
\answerday{22}{Models: explain one column from several}{answers:beginner:22}
\begin{enumerate}
  \item Which word compares one average with a number you choose? \fromday{20}\\ $\rightarrow$ \code{table.\allowbreak{}compare\_\allowbreak{}average\_\allowbreak{}to(\textquotedbl{}marks\textquotedbl{},\allowbreak{} 60)}
  \item Which word checks whether values are bell-shaped (normal)? \fromday{19}\\ $\rightarrow$ \code{table.\allowbreak{}compare\_\allowbreak{}shape\_\allowbreak{}to(\textquotedbl{}marks\textquotedbl{},\allowbreak{} \textquotedbl{}normal\textquotedbl{})}
  \item Which word adds a new column made from other columns? \fromday{15}\\ $\rightarrow$ \code{table.\allowbreak{}add\_\allowbreak{}column(\textquotedbl{}total\textquotedbl{},\allowbreak{} \textquotedbl{}Maths + Science\textquotedbl{})}, or \code{table.\allowbreak{}add\_\allowbreak{}column(\textquotedbl{}per\_\allowbreak{}hour\textquotedbl{},\allowbreak{} by\_\allowbreak{}dividing=[\textquotedbl{}marks\textquotedbl{},\allowbreak{} \textquotedbl{}hours\textquotedbl{}])}
  \item Which option adds error bars that say how sure we are of the average? \fromday{8}\\ $\rightarrow$ \code{error=\textquotedbl{}standard error\textquotedbl{}}, or \code{error=\textquotedbl{}confidence interval\textquotedbl{}} for the 95 \% interval.
\end{enumerate}
\par\noindent{\sffamily\small\bfseries\color{tryframe}Task 1.} School table: explain \code{\textquotedbl{}days\_\allowbreak{}absent\textquotedbl{}} using \code{\textquotedbl{}sleep\textquotedbl{}} and \code{\textquotedbl{}breakfast\textquotedbl{}} with \code{outcome=\textquotedbl{}counts\textquotedbl{}}. Does breakfast matter?\par
\begin{lstlisting}[style=py]
import simplydata as sdata
table = sdata.table(sdata.example("school"))
table.explain("days_absent", using=["sleep", "breakfast"], outcome="counts")
\end{lstlisting}
\begin{answerbox}Counts of days call for Poisson regression; read the line for breakfast and its p-value.\end{answerbox}
{\small\hyperref[day:beginner:22]{Back to day 22}}
\answerday{23}{Patterns across many columns}{answers:beginner:23}
\begin{enumerate}
  \item Which word explains marks from hours, sleep and breakfast together? \fromday{22}\\ $\rightarrow$ \code{table.\allowbreak{}explain(\textquotedbl{}marks\textquotedbl{},\allowbreak{} using=[\textquotedbl{}hours\textquotedbl{},\allowbreak{} \textquotedbl{}sleep\textquotedbl{},\allowbreak{} \textquotedbl{}breakfast\textquotedbl{}])}
  \item What should a results sentence give? \fromday{20}\\ $\rightarrow$ The summary of each group with n, the effect size (with its interval), the name of the test and the exact p-value.
  \item Which word gives several summaries at once? \fromday{16}\\ $\rightarrow$ \code{table.\allowbreak{}summary\_\allowbreak{}of(\textquotedbl{}sales\textquotedbl{},\allowbreak{} for\_\allowbreak{}each=\textquotedbl{}shop\textquotedbl{},\allowbreak{} show=[\textquotedbl{}average\textquotedbl{},\allowbreak{} \textquotedbl{}median\textquotedbl{},\allowbreak{} \textquotedbl{}count\textquotedbl{}])}
  \item Which word sets the y-axis from 0 to 100? \fromday{9}\\ $\rightarrow$ \code{chart.\allowbreak{}range\_\allowbreak{}y(0,\allowbreak{} 100)}
\end{enumerate}
\par\noindent{\sffamily\small\bfseries\color{tryframe}Task 1.} Meadow table: add 3 clusters and count the rows of each cluster across \code{\textquotedbl{}field\textquotedbl{}}.\par
\begin{lstlisting}[style=py]
import simplydata as sdata
table = sdata.table(sdata.example("meadow"))
table.add_clusters(how_many=3)
table.count_rows(for_each="cluster", across="field")
\end{lstlisting}
\begin{answerbox}If each cluster is mostly one field, the fields have different plants.\end{answerbox}
{\small\hyperref[day:beginner:23]{Back to day 23}}
\answerday{24}{Survival, and planning a study}{answers:beginner:24}
\begin{enumerate}
  \item Which word adds the principal components PC1 and PC2? \fromday{23}\\ $\rightarrow$ \code{table.\allowbreak{}add\_\allowbreak{}principal\_\allowbreak{}components(how\_\allowbreak{}many=2)}
  \item Which word lists every change you made to a table, for your methods? \fromday{20}\\ $\rightarrow$ \code{table.\allowbreak{}steps()}
  \item How do you get the number in row 1, column 2 of a grid? \fromday{17}\\ $\rightarrow$ \code{grid.\allowbreak{}number\_\allowbreak{}at(row=1,\allowbreak{} column=2)}
  \item Which word puts two charts next to each other on one page? \fromday{10}\\ $\rightarrow$ \code{sdata.\allowbreak{}side\_\allowbreak{}by\_\allowbreak{}side(a,\allowbreak{} b)}
\end{enumerate}
\par\noindent{\sffamily\small\bfseries\color{tryframe}Task 1.} How many values does each group need to find a large difference (\code{effect\_\allowbreak{}size=0.\allowbreak{}8}) with power 0.8? And with a small one (0.2)?\par
\begin{lstlisting}[style=py]
import simplydata as sdata
sdata.how_many_needed(effect_size=0.8, power=0.8)
sdata.how_many_needed(effect_size=0.2, power=0.8)
\end{lstlisting}
\begin{answerbox}About 26 for a large difference but about 394 for a small one: small effects need big studies.\end{answerbox}
{\small\hyperref[day:beginner:24]{Back to day 24}}
\answerday{25}{Fitting, finding and solving}{answers:beginner:25}
\begin{enumerate}
  \item Which word draws survival curves? \fromday{24}\\ $\rightarrow$ \code{chart.\allowbreak{}survival\_\allowbreak{}curves(time=\textquotedbl{}days\textquotedbl{},\allowbreak{} event=\textquotedbl{}died\textquotedbl{},\allowbreak{} colour=\textquotedbl{}treatment\textquotedbl{})}
  \item How do you explain a yes-or-no column such as \code{\textquotedbl{}passed\textquotedbl{}}? \fromday{22}\\ $\rightarrow$ \code{outcome=\textquotedbl{}yes or no\textquotedbl{}} (logistic regression)
  \item A p-value answers \code{\textquotedbl{}could it be chance?\textquotedbl{}}. What answers \code{\textquotedbl{}how big is it?\textquotedbl{}} \fromday{18}\\ $\rightarrow$ The effect size: the difference in averages, Cohen's d or Hedges' g, omega-squared.
  \item If \code{keep\_\allowbreak{}rows} keeps rows, what would you guess takes columns away? \fromday{11}\\ $\rightarrow$ \code{remove\_\allowbreak{}columns}: verb + thing. \code{keep\_\allowbreak{}} and \code{remove\_\allowbreak{}} work on rows, columns and numbers.
\end{enumerate}
\par\noindent{\sffamily\small\bfseries\color{tryframe}Task 1.} Growth table: fit the formula \code{\textquotedbl{}forest = a + b * (year - 2000)\textquotedbl{}} for each country. Which country loses forest fastest?\par
\begin{lstlisting}[style=py]
import simplydata as sdata
table = sdata.table(sdata.example("growth"))
table.fit_formula("forest = a + b * (year - 2000)", for_each="country")
\end{lstlisting}
\begin{answerbox}b is the change in forest each year: the country with the most negative b loses forest fastest.\end{answerbox}
\par\noindent{\sffamily\small\bfseries\color{tryframe}Task 2.} A box has a square bottom of side x and a height of 10 - x, so its volume is x * x * (10 - x). For which x between 0 and 10 is the volume biggest?\par
\begin{lstlisting}[style=py]
import simplydata as sdata
sdata.find_biggest("x * x * (10 - x)", between=[0, 10])
\end{lstlisting}
\begin{answerbox}x = 6.67 (that is, 20/3), with a volume of about 148.\end{answerbox}
{\small\hyperref[day:beginner:25]{Back to day 25}}
\answerday{26}{Distributions and values between points}{answers:beginner:26}
\begin{enumerate}
  \item Which word fits a straight line on a table? \fromday{25}\\ $\rightarrow$ \code{table.\allowbreak{}fit\_\allowbreak{}line(x=\textquotedbl{}hours\textquotedbl{},\allowbreak{} y=\textquotedbl{}marks\textquotedbl{})}
  \item Which word puts rows that are alike into groups? \fromday{23}\\ $\rightarrow$ \code{table.\allowbreak{}add\_\allowbreak{}clusters(how\_\allowbreak{}many=2)}
  \item You did many tests at once. Which word corrects the p-values? \fromday{19}\\ $\rightarrow$ \code{add\_\allowbreak{}adjusted\_\allowbreak{}p\_\allowbreak{}value\_\allowbreak{}column(method=\textquotedbl{}holm\textquotedbl{})}, on the table of results
  \item How do you look at the first 5 rows of a table? \fromday{12}\\ $\rightarrow$ \code{table.\allowbreak{}show\_\allowbreak{}first\_\allowbreak{}rows(5)}
\end{enumerate}
\par\noindent{\sffamily\small\bfseries\color{tryframe}Task 1.} Heights are normal with average 150 and standard deviation 8. What share is taller than 165? Which height do 95 \% of them stay below?\par
\begin{lstlisting}[style=py]
import simplydata as sdata
heights = sdata.distribution("normal", average=150, standard_deviation=8)
heights.chance(above=165)
heights.percentile(95)
\end{lstlisting}
\begin{answerbox}About 3 \% are taller than 165, and 95 \% are below about 163.\end{answerbox}
\par\noindent{\sffamily\small\bfseries\color{tryframe}Task 2.} Weather table: estimate the \code{\textquotedbl{}rain\textquotedbl{}} of each city at \code{month\_\allowbreak{}number} 2.5 and 7.5, joined by straight lines.\par
\begin{lstlisting}[style=py]
import simplydata as sdata
table = sdata.table(sdata.example("weather"))
rain = table.values_at("rain", at={"month_number": [2.5, 7.5]}, joined_by="straight lines", for_each="city")
rain.show()
\end{lstlisting}
\begin{answerbox}Each estimate lies halfway between the two months around it.\end{answerbox}
{\small\hyperref[day:beginner:26]{Back to day 26}}
\answerday{27}{Small project: look, clean and sum up}{answers:beginner:27}
\begin{enumerate}
  \item How do you make a normal distribution with average 50 and standard deviation 10? \fromday{26}\\ $\rightarrow$ \code{sdata.\allowbreak{}distribution(\textquotedbl{}normal\textquotedbl{},\allowbreak{} average=50,\allowbreak{} standard\_\allowbreak{}deviation=10)}
  \item Which word cuts a column of numbers into named groups? \fromday{15}\\ $\rightarrow$ \code{table.\allowbreak{}add\_\allowbreak{}group\_\allowbreak{}column(\textquotedbl{}grade\textquotedbl{},\allowbreak{} from\_\allowbreak{}column=\textquotedbl{}marks\textquotedbl{},\allowbreak{} cut\_\allowbreak{}at=[40,\allowbreak{} 60,\allowbreak{} 80],\allowbreak{} names=[\textquotedbl{}D\textquotedbl{},\allowbreak{} \textquotedbl{}C\textquotedbl{},\allowbreak{} \textquotedbl{}B\textquotedbl{},\allowbreak{} \textquotedbl{}A\textquotedbl{}])}
  \item How do you give the panels of a figure the letters A and B? \fromday{20}\\ $\rightarrow$ \code{letters=True} in \code{sdata.\allowbreak{}grid(.\allowbreak{}.\allowbreak{}.\allowbreak{})} or \code{sdata.\allowbreak{}side\_\allowbreak{}by\_\allowbreak{}side(.\allowbreak{}.\allowbreak{}.\allowbreak{})}
  \item Which word tests whether survival differs between groups? \fromday{24}\\ $\rightarrow$ \code{table.\allowbreak{}compare\_\allowbreak{}survival(time=\textquotedbl{}days\textquotedbl{},\allowbreak{} event=\textquotedbl{}died\textquotedbl{},\allowbreak{} for\_\allowbreak{}each=\textquotedbl{}treatment\textquotedbl{})}
\end{enumerate}
\par\noindent{\sffamily\small\bfseries\color{tryframe}Task 1.} Write your question. Open your file, describe its columns, deal with the empty cells and make two summaries for each group.\par
\begin{lstlisting}[style=py]
import simplydata as sdata
table = sdata.table(sdata.example("school"))
table.describe_columns()
table.count_empty_cells()
table.remove_rows_with_empty_cells("sleep")
table.average_of("marks", for_each="breakfast")
table.summary_of("sleep", for_each="breakfast", show=["average", "count"])
\end{lstlisting}
\begin{answerbox}Question: do pupils who eat breakfast get better marks? Look, clean, sum up.\end{answerbox}
{\small\hyperref[day:beginner:27]{Back to day 27}}
\answerday{28}{Small project: test, figure and results}{answers:beginner:28}
\begin{enumerate}
  \item Which word adds the values a model predicts as a new column? \fromday{22}\\ $\rightarrow$ \code{table.\allowbreak{}add\_\allowbreak{}predicted\_\allowbreak{}column(\textquotedbl{}marks\textquotedbl{},\allowbreak{} using=[\textquotedbl{}hours\textquotedbl{},\allowbreak{} \textquotedbl{}breakfast\textquotedbl{}])}
  \item Which option gives the rank-based test, for small groups that are not bell-shaped? \fromday{18}\\ $\rightarrow$ \code{based\_\allowbreak{}on=\textquotedbl{}ranks\textquotedbl{}}
  \item A word for one axis ends in ...? \fromday{11}\\ $\rightarrow$ \_x or \_y: \code{label\_\allowbreak{}y}, \code{range\_\allowbreak{}x}, \code{log\_\allowbreak{}scale\_\allowbreak{}y}.
  \item How do you fit a curve with x and x squared? \fromday{25}\\ $\rightarrow$ \code{table.\allowbreak{}fit\_\allowbreak{}curve(x=\textquotedbl{}hours\textquotedbl{},\allowbreak{} y=\textquotedbl{}marks\textquotedbl{},\allowbreak{} degree=2)}
\end{enumerate}
\par\noindent{\sffamily\small\bfseries\color{tryframe}Task 1.} Test your question, draw one figure and save it, and write two sentences of results.\par
\begin{lstlisting}[style=py]
import simplydata as sdata
table = sdata.table(sdata.example("school"))
table.compare_averages("marks", for_each="breakfast")
chart = table.chart()
chart.boxes(x="breakfast", y="marks")
chart.all_dots(x="breakfast", y="marks")
chart.label_y("Mark (out of 100)")
chart.save("breakfast.png")
\end{lstlisting}
\begin{answerbox}Two sentences: the averages of both groups, then the test's name, the difference and the p-value.\end{answerbox}
{\small\hyperref[day:beginner:28]{Back to day 28}}
\answerday{29}{Final test from memory}{answers:beginner:29}
\begin{enumerate}
  \item Which word tests whether groups differ in their make-up across many columns (PERMANOVA)? \fromday{23}\\ $\rightarrow$ \code{table.\allowbreak{}compare\_\allowbreak{}profiles(for\_\allowbreak{}each=\textquotedbl{}grazed\textquotedbl{})}
  \item Which word gives the chance of a value above 80? \fromday{26}\\ $\rightarrow$ \code{dist.\allowbreak{}chance(above=80)}
  \item Which word says how many values you need in each group? \fromday{24}\\ $\rightarrow$ \code{sdata.\allowbreak{}how\_\allowbreak{}many\_\allowbreak{}needed(effect\_\allowbreak{}size=0.\allowbreak{}5,\allowbreak{} power=0.\allowbreak{}8)}
  \item Which word hands a table on to pandas? \fromday{20}\\ $\rightarrow$ \code{table.\allowbreak{}as\_\allowbreak{}pandas()}
  \item Which word fits a formula of your own, such as growth? \fromday{25}\\ $\rightarrow$ \code{table.\allowbreak{}fit\_\allowbreak{}formula(\textquotedbl{}population = a * exp(b * (year - 2000))\textquotedbl{})}
  \item Which file endings give a picture that stays sharp at any size? \fromday{11}\\ $\rightarrow$ \code{.\allowbreak{}pdf} and \code{.\allowbreak{}svg}
  \item What does every simplydata program start with? \fromday{2}\\ $\rightarrow$ import simplydata as sdata
  \item Where in an error message do you look first? \fromday{3}\\ $\rightarrow$ At the last line: it names the kind of error and says what went wrong.
  \item How do you name the y-axis? \fromday{4}\\ $\rightarrow$ \code{chart.\allowbreak{}label\_\allowbreak{}y(\textquotedbl{}Marks out of 100\textquotedbl{})}
  \item Which word lets you choose the colours of the groups? \fromday{5}\\ $\rightarrow$ \code{chart.\allowbreak{}colours(\textquotedbl{}orange\textquotedbl{},\allowbreak{} \textquotedbl{}teal\textquotedbl{},\allowbreak{} \textquotedbl{}purple\textquotedbl{})}
  \item Which word keeps only the rows where subject is Maths? \fromday{13}\\ $\rightarrow$ \code{table.\allowbreak{}keep\_\allowbreak{}rows(\textquotedbl{}subject\textquotedbl{},\allowbreak{} equal\_\allowbreak{}to=\textquotedbl{}Maths\textquotedbl{})}
  \item Which word draws one dot for each row, x against y? \fromday{6}\\ $\rightarrow$ \code{chart.\allowbreak{}dots(x=\textquotedbl{}hours\textquotedbl{},\allowbreak{} y=\textquotedbl{}marks\textquotedbl{})}
  \item How do you give the number 12 the name age? \fromday{1}\\ $\rightarrow$ age = 12. The name goes on the left of =, the value on the right.
  \item Which word draws a straight line through the dots? \fromday{8}\\ $\rightarrow$ \code{chart.\allowbreak{}fit\_\allowbreak{}line()}, after \code{chart.\allowbreak{}dots(.\allowbreak{}.\allowbreak{}.\allowbreak{})}
\end{enumerate}
\par\noindent{\sffamily\small\bfseries\color{tryframe}Task 1.} From memory, in an empty file: open the \code{\textquotedbl{}marks\textquotedbl{}} table, remove the rows with empty cells, give the average marks for each class, test the difference, and save a two-panel figure (boxes, and dots of hours against marks with a fitted line) as a PDF.\par
\begin{lstlisting}[style=py]
import simplydata as sdata
table = sdata.table(sdata.example("marks"))
table.remove_rows_with_empty_cells()
table.average_of("marks", for_each="class")
table.compare_averages("marks", for_each="class")
a = table.chart()
a.boxes(x="class", y="marks")
b = table.chart()
b.dots(x="hours", y="marks")
b.fit_line()
figure = sdata.grid(a, b, per_row=2, letters=True)
figure.save("final.pdf")
\end{lstlisting}
\begin{answerbox}If you wrote this from memory, you know the core of simplydata. Words you missed: run their \code{sdata.\allowbreak{}help} again tomorrow.\end{answerbox}
{\small\hyperref[day:beginner:29]{Back to day 29}}
\answerday{30}{Look back, and what comes next}{answers:beginner:30}
\begin{enumerate}
  \item Which word fits a distribution to a column? \fromday{26}\\ $\rightarrow$ \code{table.\allowbreak{}fit\_\allowbreak{}distribution(\textquotedbl{}days\_\allowbreak{}absent\textquotedbl{},\allowbreak{} \textquotedbl{}negative binomial\textquotedbl{})}
  \item Which word finds where a formula is smallest? \fromday{25}\\ $\rightarrow$ \code{sdata.\allowbreak{}find\_\allowbreak{}smallest(\textquotedbl{}x**2 - 4*x + 1\textquotedbl{},\allowbreak{} between=[0,\allowbreak{} 5])}
  \item Which word adds the principal components PC1 and PC2? \fromday{23}\\ $\rightarrow$ \code{table.\allowbreak{}add\_\allowbreak{}principal\_\allowbreak{}components(how\_\allowbreak{}many=2)}
  \item Which word gives the average of a column for each group? \fromday{16}\\ $\rightarrow$ \code{table.\allowbreak{}average\_\allowbreak{}of(\textquotedbl{}marks\textquotedbl{},\allowbreak{} for\_\allowbreak{}each=\textquotedbl{}class\textquotedbl{})}
\end{enumerate}
\par\noindent{\sffamily\small\bfseries\color{tryframe}Task 1.} Do the task of your hardest day again from memory. Then hand the marks table to pandas and print its first rows with pandas' own word \code{head()}.\par
\begin{lstlisting}[style=py]
import simplydata as sdata
table = sdata.table(sdata.example("marks"))
df = table.as_pandas()
print(df.head())
\end{lstlisting}
\begin{answerbox}\code{head()} is pandas' word for \code{show\_\allowbreak{}first\_\allowbreak{}rows()}. simplydata knows many such names: \code{sdata.\allowbreak{}help(\textquotedbl{}head\textquotedbl{})} points to the simplydata word.\end{answerbox}
{\small\hyperref[day:beginner:30]{Back to day 30}}

\end{document}
